% La période de mandat — Histoire 1re Tech — Cours signé OnBuch+
% Programme officiel MINESEC. Compiler avec : tectonic N06-periode-mandat.tex
\newcommand{\DOCMATIERE}{Histoire}
\newcommand{\DOCNIVEAU}{1re Tech}
\newcommand{\DOCMODULE}{Histoire – classe de Première technique}
\newcommand{\DOCLECON}{N06}
\newcommand{\DOCTITRE}{La période de mandat}
% OnBuch+ — préambule « cours signés » ENRICHI (compilé avec tectonic / XeLaTeX)
% Système visuel « Pop » d'OnBuch+ : fond crème (jamais blanc), bordures encre
% épaisses, ombres « dures » décalées, titres Archivo Black en capitales.
% Version MATHÉMATIQUES : graphes (pgfplots/tikz), boîtes pédagogiques
% (définition, propriété, méthode, attention, le savais-tu, exercice résolu,
% exercice, TP, bilan), page de garde et signature OnBuch+ ancrée en bas.
%
% Macros attendues dans main.tex AVANT \input{preamble} :
%   \DOCMATIERE  \DOCNIVEAU  \DOCMODULE  \DOCLECON (numéro)  \DOCTITRE
\documentclass[11pt]{article}

\usepackage{fontspec}
\usepackage[french]{babel}
\usepackage[a4paper,margin=2cm,top=3.4cm,bottom=2.8cm,headheight=1.4cm,footskip=1.3cm]{geometry}
\usepackage{amsmath,amssymb,amsfonts}
\usepackage{mathtools} % \xmapsto, \coloneqq... souvent utilisés par les modèles
\usepackage{cancel} % simplifications barrées (bilans de forces)
% \abs{z} (module, valeur absolue) et \norm{u} (norme), utilisés par les modèles
\providecommand{\abs}{}\let\abs\relax\DeclarePairedDelimiter{\abs}{\lvert}{\rvert}
\providecommand{\norm}{}\let\norm\relax\DeclarePairedDelimiter{\norm}{\lVert}{\rVert}
\usepackage[table]{xcolor} % [table] : fournit aussi \rowcolors (alternance de lignes)
\usepackage{pagecolor}
\usepackage{tikz}
\usetikzlibrary{arrows.meta,positioning,calc,shapes.geometric,shapes.misc,shapes.arrows,decorations.pathmorphing,decorations.pathreplacing,decorations.markings,patterns,shadows,mindmap,trees,fit,backgrounds,matrix,intersections,angles,quotes}
\usepackage{pgfplots}
\usepackage[version=4]{mhchem} % équations nucléaires \ce{^{235}_{92}U}
\usepackage[european,siunitx]{circuitikz} % schémas électriques
\pgfplotsset{compat=1.18}
\usepgfplotslibrary{fillbetween,groupplots} % aires entre courbes (calcul intégral)
\usepackage{enumitem}
\usepackage{fancyhdr}
% [most] charge toutes les bibliothèques tcolorbox (listings, minted, raster,
% documentation...) et gonfle inutilement la mémoire TeX sur un document long
% (~300 boîtes) : on ne charge que ce qui est réellement utilisé.
\usepackage[skins,breakable]{tcolorbox}
\usepackage{graphicx}
\usepackage{colortbl}
\usepackage{booktabs}
\usepackage{tabularx}
\usepackage{multirow}
\usepackage{diagbox} % case d'en-tête diagonale des tableaux à double entrée
% Colonnes extensibles centrée / gauche / droite (C, L, R), souvent utilisées par les modèles
\newcolumntype{C}{>{\centering\arraybackslash}X}
\newcolumntype{L}{>{\raggedright\arraybackslash}X}
\newcolumntype{R}{>{\raggedleft\arraybackslash}X}
\usepackage{array}
\usepackage{multicol}
\usepackage{siunitx}
% parse-numbers=false : accepte aussi \SI{2,5 \times 10^{-3}}{...}
\sisetup{locale=FR,output-decimal-marker={,},group-minimum-digits=4,parse-numbers=false}
\DeclareSIUnit\ohm{\ensuremath{\symup{\Omega}}} % U+2126 absent de la police
\DeclareSIUnit\division{div} % réglages d'oscilloscope (ms/div, V/div)
\DeclareSIUnit\tr{tr} % tours (vitesse de rotation, ex. \SI{1500}{\tr\per\minute})
\DeclareSIUnit\molar{M} % concentration molaire (ex. \SI{3}{\molar}, \SI{1}{\micro\molar})
\DeclareSIUnit\an{an} % année (le modèle confond parfois avec \year, un compteur TeX) — ex. \SI{5}{\kilo\meter\per\an}
\DeclareSIUnit\FCFA{FCFA} % franc CFA (ex. \SI{500}{\FCFA\per\kilo\gram})
\DeclareSIUnit\degreeBrix{\textdegree Brix} % teneur en sucre d'un sirop/jus (réfractométrie)
\DeclareSIUnit\month{mois} % le modèle utilise parfois \month, un compteur TeX, comme unité de durée
\DeclareSIUnit\microliter{\micro\liter} % le modèle écrit parfois \microliter en un seul token
\DeclareSIUnit\million{millions} % dénombrement (ex. \SI{15}{\million\per\mL} spermatozoïdes)
\usepackage{qrcode}
% \degree : commande fréquemment utilisée par les modèles pour un angle ou une
% température en dehors de \SI{}{} (non fournie par les paquets chargés).
\providecommand{\degree}{^{\circ}}
% \qed (fin de démonstration), utilisé par les modèles sans amsthm
\providecommand{\qed}{\hfill$\square$}
% \barre : double barre des valeurs interdites dans un tableau de variations (tkz-tab)
\providecommand{\barre}{\ensuremath{\|}}
% environnement proof (amsthm non chargé) utilisé par les modèles
\ifdefined\proof\else\newenvironment{proof}[1][Démonstration]{\par\noindent\textit{#1.}\ }{\qed\par}\fi
\usepackage{lastpage}
\usepackage{hyperref}
\hypersetup{hidelinks}

% ============================================================
% Palette « Pop » OnBuch+ — mêmes hex que la classe P (plus_ui.dart)
% ============================================================
\definecolor{poporange}{HTML}{F59321}
\definecolor{popdark}{HTML}{E07A0C}
\definecolor{popcream}{HTML}{FFF6E8}
\definecolor{popink}{HTML}{1C1714}
\definecolor{poppurple}{HTML}{7A5AE0}
\definecolor{popgreen}{HTML}{2E9E6A}
\definecolor{poppink}{HTML}{E0457B}
\definecolor{popblue}{HTML}{2F7DE1}
\definecolor{popgold}{HTML}{C98A12}
\definecolor{popmuted}{HTML}{8A7F76}
\definecolor{popline}{HTML}{EADFD2}
\definecolor{popgray}{HTML}{8A7F76}
\colorlet{popgrey}{popgray}
% Alias tolérants (le modèle nomme parfois les couleurs différemment)
\colorlet{popwhite}{white}
\colorlet{popblack}{popink}
\colorlet{popred}{poppink}
\colorlet{popyellow}{popgold}
% Teintes claires (fonds de tableaux/encadrés) — le modèle nomme parfois
% les couleurs avec un suffixe L (ex. popblueL) sans les avoir définies.
\colorlet{poporangeL}{poporange!20}
\colorlet{popdarkL}{popdark!20}
\colorlet{poppurpleL}{poppurple!20}
\colorlet{popgreenL}{popgreen!20}
\colorlet{poppinkL}{poppink!20}
\colorlet{popblueL}{popblue!20}
\colorlet{popgoldL}{popgold!20}
\colorlet{popredL}{popred!20}
\colorlet{popgrayL}{popgray!20}
% Teintes claires pour fonds de tableaux / zones
\colorlet{popblueL}{popblue!10!white}
\colorlet{poporangeL}{poporange!14!white}
\colorlet{popgreenL}{popgreen!12!white}
\colorlet{poppurpleL}{poppurple!12!white}
\colorlet{poppinkL}{poppink!12!white}
\colorlet{popgoldL}{popgold!14!white}

\pagecolor{popcream}

% ============================================================
% Typographie
% ============================================================
\defaultfontfeatures{Ligatures=TeX}
\setmainfont{PlusJakartaSans-Medium.ttf}[Path=../../../fonts/,BoldFont=PlusJakartaSans-Bold.ttf]
% Maths en sans-serif (Fira Math) avec chiffres et lettres droites de Plus
% Jakarta, pour que les formules se fondent dans le texte.
\usepackage[math-style=ISO,bold-style=ISO]{unicode-math}
\setmathfont{FiraMath-Regular.otf}[Path=../../../fonts/]
\setmathfont{PlusJakartaSans-Medium.ttf}[Path=../../../fonts/,range={up/num,up/{Latin,latin},\mathup}]
\usepackage{icomma} % virgule décimale sans espace en mode math
% Caractères mathématiques Unicode absents des polices de texte (Plus Jakarta,
% Archivo Black) : rendus par leur équivalent mathématique, en texte comme en
% formule — sinon ils s'affichent en carré vide (ex. « Arithmétique dans ℤ »).
\usepackage{newunicodechar}
\newunicodechar{ℝ}{\ensuremath{\mathbb{R}}}
\newunicodechar{ℤ}{\ensuremath{\mathbb{Z}}}
\newunicodechar{ℂ}{\ensuremath{\mathbb{C}}}
\newunicodechar{ℕ}{\ensuremath{\mathbb{N}}}
\newunicodechar{ℚ}{\ensuremath{\mathbb{Q}}}
\newunicodechar{𝔻}{\ensuremath{\mathbb{D}}}
\newunicodechar{α}{\ensuremath{\alpha}}
\newunicodechar{β}{\ensuremath{\beta}}
\newunicodechar{γ}{\ensuremath{\gamma}}
\newunicodechar{δ}{\ensuremath{\delta}}
\newunicodechar{ε}{\ensuremath{\varepsilon}}
\newunicodechar{θ}{\ensuremath{\theta}}
\newunicodechar{λ}{\ensuremath{\lambda}}
\newunicodechar{μ}{\ensuremath{\mu}}
\newunicodechar{σ}{\ensuremath{\sigma}}
\newunicodechar{τ}{\ensuremath{\tau}}
\newunicodechar{φ}{\ensuremath{\varphi}}
\newunicodechar{ψ}{\ensuremath{\psi}}
\newunicodechar{ω}{\ensuremath{\omega}}
\newunicodechar{Σ}{\ensuremath{\Sigma}}
% majuscules produites par \MakeUppercase dans les titres (α-aminés -> Α)
\newunicodechar{Α}{\ensuremath{\alpha}}
\newunicodechar{Β}{\ensuremath{\beta}}
\newunicodechar{Γ}{\ensuremath{\Gamma}}
\newunicodechar{Δ}{\ensuremath{\Delta}}
\newunicodechar{Ε}{\ensuremath{\varepsilon}}
\newunicodechar{Θ}{\ensuremath{\Theta}}
\newunicodechar{Λ}{\ensuremath{\Lambda}}
\newunicodechar{Μ}{\ensuremath{\mu}}
\newunicodechar{Τ}{\ensuremath{\tau}}
\newunicodechar{Φ}{\ensuremath{\Phi}}
\newunicodechar{Ψ}{\ensuremath{\Psi}}
\newunicodechar{Ω}{\ensuremath{\Omega}}
\newunicodechar{↦}{\ensuremath{\mapsto}}
\newunicodechar{∈}{\ensuremath{\in}}
\newunicodechar{∉}{\ensuremath{\notin}}
\newunicodechar{∧}{\ensuremath{\wedge}}
\newunicodechar{⇔}{\ensuremath{\Leftrightarrow}}
\newunicodechar{⟺}{\ensuremath{\Longleftrightarrow}}
\newunicodechar{⇒}{\ensuremath{\Rightarrow}}
\newunicodechar{⟹}{\ensuremath{\Longrightarrow}}
\newunicodechar{∘}{\ensuremath{\circ}}
\newunicodechar{∪}{\ensuremath{\cup}}
\newunicodechar{∩}{\ensuremath{\cap}}
\newunicodechar{≡}{\ensuremath{\equiv}}
\newunicodechar{⊂}{\ensuremath{\subset}}
\newunicodechar{⊥}{\ensuremath{\perp}}
\newunicodechar{∃}{\ensuremath{\exists}}
\newunicodechar{∀}{\ensuremath{\forall}}
\newunicodechar{∛}{\ensuremath{\sqrt[3]{\,}}}
\newunicodechar{∜}{\ensuremath{\sqrt[4]{\,}}}
\newunicodechar{⃗}{\ensuremath{{}^{\to}}}
\newunicodechar{ⁿ}{\ifmmode^{n}\else\textsuperscript{\ensuremath{n}}\fi}
\newunicodechar{ˣ}{\ifmmode^{x}\else\textsuperscript{\ensuremath{x}}\fi}
\newunicodechar{ᵇ}{\ifmmode^{b}\else\textsuperscript{\ensuremath{b}}\fi}
\newunicodechar{ʳ}{\ifmmode^{r}\else\textsuperscript{\ensuremath{r}}\fi}
\newunicodechar{ᵘ}{\ifmmode^{u}\else\textsuperscript{\ensuremath{u}}\fi}
\newunicodechar{ʸ}{\ifmmode^{y}\else\textsuperscript{\ensuremath{y}}\fi}
\newunicodechar{ᵃ}{\ifmmode^{a}\else\textsuperscript{\ensuremath{a}}\fi}
\newunicodechar{ᵉ}{\ifmmode^{e}\else\textsuperscript{\ensuremath{e}}\fi}
\newunicodechar{ᵏ}{\ifmmode^{k}\else\textsuperscript{\ensuremath{k}}\fi}
\newunicodechar{ᵖ}{\ifmmode^{p}\else\textsuperscript{\ensuremath{p}}\fi}
\newunicodechar{ᴺ}{\ifmmode^{N}\else\textsuperscript{\ensuremath{N}}\fi}
\newunicodechar{ᶜ}{\ifmmode^{c}\else\textsuperscript{\ensuremath{c}}\fi}
\newunicodechar{ʰ}{\ifmmode^{h}\else\textsuperscript{\ensuremath{h}}\fi}
\newunicodechar{ᵠ}{\ifmmode^{q}\else\textsuperscript{\ensuremath{q}}\fi}
\newunicodechar{ₙ}{\ifmmode_{n}\else\textsubscript{\ensuremath{n}}\fi}
\newunicodechar{ₐ}{\ifmmode_{a}\else\textsubscript{\ensuremath{a}}\fi}
\newunicodechar{ᵢ}{\ifmmode_{i}\else\textsubscript{\ensuremath{i}}\fi}
\newunicodechar{ᵣ}{\ifmmode_{r}\else\textsubscript{\ensuremath{r}}\fi}
\newunicodechar{ₖ}{\ifmmode_{k}\else\textsubscript{\ensuremath{k}}\fi}
\newunicodechar{ᵦ}{\ifmmode_{\beta}\else\textsubscript{\ensuremath{\beta}}\fi}
\newunicodechar{ₓ}{\ifmmode_{x}\else\textsubscript{\ensuremath{x}}\fi}
\newfontfamily\popdisplay{ArchivoBlack-Regular.ttf}[Path=../../../fonts/]
\newfontfamily\poplogo{SpaceGrotesk-Bold.ttf}[Path=../../../fonts/]
\newfontfamily\popsemi{PlusJakartaSans-SemiBold.ttf}[Path=../../../fonts/]
\newfontfamily\popxbold{PlusJakartaSans-ExtraBold.ttf}[Path=../../../fonts/]
\newfontfamily\popxboldls{PlusJakartaSans-ExtraBold.ttf}[Path=../../../fonts/,LetterSpace=3.0]

\setlist[enumerate]{leftmargin=1.7em,itemsep=5pt,topsep=4pt}
\setlist[itemize]{leftmargin=1.7em,itemsep=4pt,topsep=4pt,label={\color{poporange}\textbullet}}
\setlength{\parindent}{0pt}
\setlength{\parskip}{5pt}
\renewcommand{\arraystretch}{1.25}

% pgfplots : style Pop par défaut
\pgfplotsset{
  every axis/.append style={axis line style={popink,line width=1pt},tick style={popink},
    grid style={popline,line width=0.5pt},label style={font=\small},tick label style={font=\footnotesize}},
  popcurve/.style={poporange,line width=1.8pt,no markers,smooth},
}
% Même style disponible dans un \draw TikZ simple (hors pgfplots)
\tikzset{popcurve/.style={poporange,line width=1.8pt}}
\tikzset{popgrid/.style={popline,very thin}}
\colorlet{popcurve}{poporange} % le nom du style est parfois utilisé comme couleur

% ============================================================
% Pill / squiggle
% ============================================================
\newcommand{\poppill}[3]{%
  \tikz[baseline=-0.55ex]\node[draw=popink,line width=1.2pt,rounded corners=2.6pt,%
    fill=#2,inner xsep=6.5pt,inner ysep=3pt]{%
    \popxboldls\fontsize{7}{7}\selectfont\color{#3}\MakeUppercase{#1}};%
}
\newcommand{\popsquiggle}[1]{%
  \tikz[baseline=-0.4ex]{
    \draw[#1,line width=2.6pt,line cap=round]
      plot [smooth] coordinates {(0,0.09) (0.34,0.01) (0.68,0.21) (1.02,0.01) (1.36,0.10)};
  }%
}
% Mot-clé surligné (marqueur orange sous le texte)
\newcommand{\cle}[1]{{\popxbold\color{popdark}#1}}

% ============================================================
% En-tête / pied
% ============================================================
\pagestyle{fancy}
\fancyhf{}
\renewcommand{\headrulewidth}{0pt}
\renewcommand{\footrulewidth}{0pt}
\lhead{\raisebox{-0.15ex}{\poplogo\fontsize{13}{13}\selectfont\color{popink}OnBuch\color{poporange}+}}
\rhead{\poppill{\DOCMATIERE\ \textbullet\ \DOCNIVEAU\ \textbullet\ Leçon \DOCLECON}{white}{popink}}
\lfoot{\popxbold\fontsize{7.6}{7.6}\selectfont\color{popmuted}\MakeUppercase{Cours signé OnBuch+}\ \textbullet\ {\popsemi\DOCTITRE}}
\rfoot{\tikz[baseline=-0.6ex]\node[rounded corners=8.5pt,fill=popink,minimum height=17pt,minimum width=17pt,inner xsep=5pt,inner sep=0]{\popxbold\fontsize{8}{8}\selectfont\color{popcream}\thepage\,/\,\pageref*{LastPage}};}

% ============================================================
% Titres
% ============================================================
\newcommand{\chaptitre}[2]{%
  \begin{tcolorbox}[enhanced,colback=poporange,colframe=popink,boxrule=2.6pt,arc=11pt,
      drop shadow=popink,left=15pt,right=15pt,top=12pt,bottom=11pt,before skip=3pt,after skip=18pt]
    \poppill{#2}{popink}{popcream}\par\vspace{9pt}
    {\raggedright\hyphenpenalty=10000\exhyphenpenalty=10000\popdisplay\fontsize{22}{24}\selectfont\color{popcream}\MakeUppercase{#1}\par}
  \end{tcolorbox}
}
% Section numérotée : pastille orange avec le numéro + titre Archivo
\newcounter{popsec}
\newcommand{\coursec}[1]{%
  \stepcounter{popsec}\setcounter{popsub}{0}%
  \par\vspace{16pt}\noindent
  \begin{minipage}{\linewidth}
  \tikz[baseline=-0.7ex]\node[draw=popink,line width=1.6pt,fill=poporange,rounded corners=4pt,minimum size=22pt,inner sep=0]{\popdisplay\fontsize{12}{12}\selectfont\color{popcream}\thepopsec};%
  \hspace{8pt}{\popdisplay\fontsize{15}{17}\selectfont\color{popink}\MakeUppercase{#1}}\par
  \vspace{4pt}\noindent{\color{poporange}\rule{36pt}{3.6pt}}
  \end{minipage}\par\nopagebreak\vspace{8pt}
}
\newcounter{popsub}
\newcommand{\courssub}[1]{%
  \stepcounter{popsub}%
  \par\vspace{9pt}\noindent{\popxbold\fontsize{11.5}{13}\selectfont\color{popink}{\color{poporange}\thepopsec.\thepopsub}\hspace{6pt}#1}\par\nopagebreak\vspace{4pt}
}

% ============================================================
% Boîtes pédagogiques — même recette : fond blanc, bordure épaisse colorée,
% ombre dure encre, badge plein. Titre optionnel [#1].
% ============================================================
\tcbset{popbase/.style={breakable,enhanced,colback=white,boxrule=2.3pt,arc=9pt,
  shadow={3pt}{-3pt}{0pt}{popink},left=14pt,right=14pt,top=11pt,bottom=11pt,before skip=11pt,after skip=15pt,
  fontupper=\popsemi\fontsize{10.6}{14.5}\selectfont\color{popink},
  fontlower=\popsemi\fontsize{10.6}{14.5}\selectfont\color{popink}}}
% Coche dessinée (le glyphe ✓ n'existe pas dans les polices du document)
\newcommand{\popcheck}[1]{\tikz[baseline=-0.5ex]\draw[#1,line width=1.6pt,line cap=round,line join=round](-0.13,0.0)--(-0.04,-0.09)--(0.13,0.1);}
\providecommand{\coche}{\popcheck{popgreen}}
\providecommand{\resultat}[1]{\par\smallskip\noindent\fcolorbox{popgreen}{popgreenL}{\parbox{\dimexpr\linewidth-2\fboxsep-2\fboxrule\relax}{\textbf{Résultat.} #1}}\par\smallskip}
\newcommand{\popboxhead}[3]{\poppill{#1}{#2}{white}\ifx\relax#3\relax\else\hspace{6pt}{\popxbold\fontsize{10.6}{12}\selectfont\color{popink}#3}\fi\par\vspace{7pt}}

\newtcolorbox{aretenir}[1][]{popbase,colframe=popgreen,before upper={\popboxhead{À retenir}{popgreen}{#1}}}
\newtcolorbox{exemplebox}[1][]{popbase,colframe=poporange,before upper={\popboxhead{Exemple}{poporange}{#1}}}
\newtcolorbox{definition}[1][]{popbase,colframe=popblue,before upper={\popboxhead{Définition}{popblue}{#1}}}
\newtcolorbox{propriete}[1][]{popbase,colframe=poppurple,before upper={\popboxhead{Propriété}{poppurple}{#1}}}
\newtcolorbox{methode}[1][]{popbase,colframe=popgold,before upper={\popboxhead{Méthode}{popgold}{#1}}}
\newtcolorbox{attention}[1][]{popbase,colframe=poppink,before upper={\popboxhead{Attention}{poppink}{#1}}}
\newtcolorbox{experience}[1][]{popbase,colframe=popblue,colback=popblueL,before upper={\popboxhead{Expérience}{popblue}{#1}}}
% « Le savais-tu ? » : carte violette pleine (contexte, culture, Cameroun)
\newtcolorbox{savaistu}[1][]{popbase,colback=poppurple,colframe=popink,
  fontupper=\popsemi\fontsize{10.4}{14.2}\selectfont\color{white},
  before upper={\poppill{Le savais-tu\,?}{white}{poppurple}\ifx\relax#1\relax\else\hspace{6pt}{\popxbold\color{white}#1}\fi\par\vspace{7pt}}}
% Exercice résolu : énoncé en haut, \tcblower puis la solution sur fond crème
\newcounter{popexo}\newcounter{popexr}
\newtcolorbox{exoresolu}[1][]{popbase,colframe=poporange,colbacklower=popcream,
  segmentation style={popink,line width=1.2pt,dashed},
  before upper={\stepcounter{popexr}\popboxhead{Exercice résolu \thepopexr}{poporange}{#1}},
  before lower={\poppill{Solution}{popink}{popcream}\par\vspace{6pt}}}
% Exercice (non résolu en place) — la correction est dans la partie Corrigés
\newtcolorbox{exercice}[1][]{popbase,colframe=popink,
  before upper={\stepcounter{popexo}\popboxhead{Exercice \thepopexo}{popink}{#1}}}
% Corrigé
\newtcolorbox{corrige}[1][]{popbase,colframe=popgreen,colback=popgreenL,
  before upper={\popboxhead{Corrigé}{popgreen}{#1}}}
% Objectifs / prérequis / bilan
\newtcolorbox{objectifs}{popbase,colframe=popink,colback=poporangeL,
  before upper={\popboxhead{Objectifs de la leçon}{popink}{}}}
\newtcolorbox{prerequis}{popbase,colframe=popink,colback=popblueL,
  before upper={\popboxhead{Prérequis}{popblue}{}}}
\newtcolorbox{bilan}{popbase,colframe=popink,colback=white,boxrule=2.8pt,
  before upper={\popboxhead{Fiche bilan}{poporange}{L'essentiel à connaître pour l'examen}}}
% Figure encadrée avec légende
\newtcolorbox{popfigure}[1][]{enhanced,breakable=false,colback=white,colframe=popline,boxrule=1.4pt,arc=8pt,
  left=10pt,right=10pt,top=10pt,bottom=8pt,before skip=10pt,after skip=12pt,halign=center,
  lower separated=false,halign lower=center,
  fontlower=\popsemi\fontsize{9}{11.5}\selectfont\color{popmuted},
  #1}
\newcommand{\legende}[1]{\par\vspace{6pt}{\popsemi\fontsize{9}{11.5}\selectfont\color{popmuted}#1}}

% ============================================================
% Signature OnBuch+
% ============================================================
\newcommand{\poplogoinline}[1]{{\poplogo\fontsize{#1}{#1}\selectfont\color{popink}OnBuch\color{poporange}+}}
\newcommand{\signatureonbuch}{%
  \begin{tcolorbox}[enhanced,colback=popink,colframe=popink,boxrule=2.2pt,arc=10pt,
      drop shadow=poporange,left=14pt,right=14pt,top=10pt,bottom=10pt,before skip=0pt,after skip=0pt]
    \tikz[baseline=-0.7ex]\node[circle,fill=popgreen,draw=popcream,line width=1.3pt,minimum size=22pt,inner sep=0]{\popcheck{white}};%
    \hspace{9pt}\begin{minipage}[c]{0.62\linewidth}
      {\popxbold\fontsize{12}{13}\selectfont\color{popcream}Signé\ \textbullet\ {\poplogo OnBuch\color{poporange}+}}\par\vspace{2pt}
      {\raggedright\popsemi\fontsize{8}{10.5}\selectfont\color{popline}\MakeUppercase{Équipe pédagogique OnBuch+}\\\MakeUppercase{Conforme au programme officiel MINESEC}\par}
    \end{minipage}\hfill
    {\poplogo\fontsize{22}{22}\selectfont\color{popcream}OnBuch\color{poporange}+}
  \end{tcolorbox}%
}

% ============================================================
% Page de garde
% #1 titre · #2 matière · #3 niveau · #4 module · #5 n° leçon · #6 durée ·
% #7 description / objectifs (texte ou itemize)
% ============================================================
\newcommand{\pagegarde}[7]{%
  \thispagestyle{empty}
  \newgeometry{a4paper,margin=1.7cm,top=1.6cm,bottom=1.4cm}%
  \begin{tikzpicture}[remember picture,overlay]
    \fill[poporange!18] ([xshift=-3.2cm,yshift=-3.6cm]current page.north east) circle (4.6cm);
    \fill[poppurple!14] ([xshift=2.4cm,yshift=7.6cm]current page.south west) circle (3.8cm);
  \end{tikzpicture}%
  {\poplogo\fontsize{30}{30}\selectfont\color{popink}OnBuch\color{poporange}+}\hfill
  \raisebox{0.6ex}{\poppill{Cours signé \textbullet\ Premium}{popink}{popcream}}\par
  \vspace{1pt}\popsquiggle{poppurple}\par
  \vspace{8pt}
  \begin{tcolorbox}[enhanced,colback=poporange,colframe=popink,boxrule=2.8pt,arc=13pt,
      drop shadow=popink,left=18pt,right=18pt,top=12pt,bottom=13pt,after skip=10pt]
    \poppill{#2 \textbullet\ #3}{popink}{popcream}\hspace{5pt}\poppill{#4}{white}{popink}\par\vspace{8pt}
    {\popdisplay\fontsize{14}{14}\selectfont\color{popink}LEÇON #5}\par\vspace{4pt}
    {\raggedright\hyphenpenalty=10000\exhyphenpenalty=10000\popdisplay\fontsize{25}{27}\selectfont\color{popcream}\MakeUppercase{#1}\par}
    \vspace{8pt}
    {\popxbold\fontsize{9.5}{9.5}\selectfont\color{popink}\MakeUppercase{Durée conseillée : #6}}
  \end{tcolorbox}
  \begin{tcolorbox}[enhanced,colback=white,colframe=popink,boxrule=2.2pt,arc=10pt,
      drop shadow=popink,left=16pt,right=16pt,top=10pt,bottom=10pt,after skip=8pt]
    \poppill{Dans cette leçon}{poppurple}{white}\par\vspace{5pt}
    \popsemi\fontsize{9.3}{12.6}\selectfont\color{popink}#7
  \end{tcolorbox}
  \noindent\begin{minipage}[t]{0.487\linewidth}
    \begin{tcolorbox}[enhanced,colback=popgreen,colframe=popink,boxrule=2.2pt,arc=10pt,
        drop shadow=popink,left=11pt,right=11pt,top=10pt,bottom=10pt,height=3.1cm,valign=center]
      \tcbox[colback=white,colframe=popink,boxrule=1.3pt,arc=5pt,boxsep=0pt,nobeforeafter,
        left=3pt,right=3pt,top=3pt,bottom=3pt,valign=center]{\qrcode[height=1.9cm]{https://play.google.com/store/apps/details?id=com.luvvix.onbuch&hl=fr}}%
      \hspace{8pt}\begin{minipage}[c]{0.5\linewidth}
        {\popdisplay\fontsize{11}{12.5}\selectfont\color{white}\MakeUppercase{Télécharge OnBuch}}\par\vspace{4pt}
        {\raggedright\popsemi\fontsize{8.4}{11}\selectfont\color{white}Gratuit \textbullet\ Google Play\\Tuteur IA, annales, concours, résultats\par}
      \end{minipage}
    \end{tcolorbox}
  \end{minipage}\hfill
  \begin{minipage}[t]{0.487\linewidth}
    \begin{tcolorbox}[enhanced,colback=poppurple,colframe=popink,boxrule=2.2pt,arc=10pt,
        drop shadow=popink,left=11pt,right=11pt,top=10pt,bottom=10pt,height=3.1cm,valign=center]
      \tikz[baseline=-0.7ex]\node[circle,fill=white,draw=popink,line width=1.3pt,minimum size=1.6cm,inner sep=0]{\popdisplay\fontsize{24}{24}\selectfont\color{poppurple}+};%
      \hspace{8pt}\begin{minipage}[c]{0.6\linewidth}
        {\popdisplay\fontsize{11}{12.5}\selectfont\color{white}\MakeUppercase{Passe à OnBuch+}}\par\vspace{4pt}
        {\raggedright\popsemi\fontsize{8.4}{11}\selectfont\color{white}Tous les cours signés, Tuteur IA illimité, fiches PDF — onglet OnBuch+ de l'app\par}
      \end{minipage}
    \end{tcolorbox}
  \end{minipage}
  \vspace{8pt}
  \popcontenu
  \vfill
  \signatureonbuch
  \restoregeometry
  \newpage
}

% Rangée « Ce cours contient » (page de garde)
\newcommand{\poptile}[3]{% #1 couleur #2 gros texte #3 légende
  \begin{tcolorbox}[enhanced,colback=#1,colframe=popink,boxrule=2pt,arc=9pt,shadow={2.5pt}{-2.5pt}{0pt}{popink},
      left=6pt,right=6pt,top=9pt,bottom=9pt,halign=center,valign=center,height=1.75cm,width=\linewidth,nobeforeafter]
    {\popdisplay\fontsize{12.5}{13}\selectfont\color{white}\MakeUppercase{#2}}\par\vspace{3pt}
    {\centering\popsemi\fontsize{7.8}{9.5}\selectfont\color{white}#3\par}
  \end{tcolorbox}}
\newcommand{\popcontenu}{%
  \par\noindent{\popxboldls\fontsize{7.5}{7.5}\selectfont\color{popmuted}CE COURS SIGNÉ CONTIENT}\par\vspace{7pt}
  \noindent\begin{minipage}[t]{0.235\linewidth}\poptile{popblue}{Cours}{complet, illustré, pas à pas}\end{minipage}\hfill
  \begin{minipage}[t]{0.235\linewidth}\poptile{poporange}{Méthodes}{et exercices résolus}\end{minipage}\hfill
  \begin{minipage}[t]{0.235\linewidth}\poptile{poppink}{Exercices}{type examen + corrigés}\end{minipage}\hfill
  \begin{minipage}[t]{0.235\linewidth}\poptile{popgreen}{Fiche bilan}{l'essentiel à retenir}\end{minipage}\par}

% Sommaire
\newtcolorbox{sommaire}{enhanced,colback=white,colframe=popink,boxrule=2.2pt,arc=10pt,
  drop shadow=popink,left=18pt,right=18pt,top=13pt,bottom=13pt,before skip=6pt,after skip=20pt,
  before upper={\poppill{Sommaire}{poppurple}{white}\par\vspace{9pt}\popsemi\fontsize{10.6}{14.5}\selectfont\color{popink}}}

% Signature de fin de cours — poussée tout en bas de la dernière page
\newcommand{\findecours}{%
  \par\vfill\nopagebreak
  \begin{center}\popsquiggle{poporange}\end{center}\vspace{4pt}
  \signatureonbuch
}
\AtBeginDocument{\def\llbracket{\mathopen{[\mkern-3.5mu[}}\def\rrbracket{\mathclose{]\mkern-3.5mu]}}} % intervalles d'entiers (glyphe ⟦ absent de la police)
% Glyphes absents de Fira Math : redéfinis à partir de glyphes disponibles
\AtBeginDocument{%
  \def\setminus{\mathbin{\backslash}}\let\smallsetminus\setminus
  \def\vdots{\mathord{\vbox{\baselineskip3.2pt\lineskiplimit0pt\kern4pt\hbox{.}\hbox{.}\hbox{.}}}}%
  \def\ddots{\mathinner{\mkern1mu\raise6.5pt\hbox{.}\mkern2mu\raise3.6pt\hbox{.}\mkern2mu\raise0.7pt\hbox{.}\mkern1mu}}%
  \def\triangle{\mathord{\scalebox{0.9}{$\symup{\Delta}$}}}%
  \def\nabla{\mathord{\raisebox{\depth}{\scalebox{1}[-1]{$\symup{\Delta}$}}}}%
  \def\checkmark{\popcheck{popdark}}%
  \def\star{\mathbin{\ast}}\def\bigstar{\mathord{\ast}}%
  \def\diamond{\mathbin{\scalebox{0.6}{$\Diamond$}}}%
  \def\bigcup{\mathop{\mathchoice{\raisebox{-0.3em}{\scalebox{1.7}{$\cup$}}}{\scalebox{1.25}{$\cup$}}{\cup}{\cup}}\displaylimits}%
  \def\bigcap{\mathop{\mathchoice{\raisebox{-0.3em}{\scalebox{1.7}{$\cap$}}}{\scalebox{1.25}{$\cap$}}{\cap}{\cap}}\displaylimits}%
  \def\lesseqqgtr{\mathrel{\lessgtr}}\def\bigoplus{\mathop{\oplus}\displaylimits}%
}
\providecommand{\ssi}{si et seulement si} % abréviation fréquente
\AtBeginDocument{\providecommand{\wideparen}{\overparen}} % arc de cercle

%% FIGFIX-BEGIN v4 — correctifs globaux des figures (docs/onbuch-generation/tools/figfix)
\usepackage{adjustbox}\usepackage{etoolbox}
% toute figure TikZ/pgfplots plus large que la ligne est réduite à la largeur disponible
\BeforeBeginEnvironment{tikzpicture}{\begin{adjustbox}{max width=\linewidth}}
\AfterEndEnvironment{tikzpicture}{\end{adjustbox}}
% légendes pgfplots lisibles par-dessus les courbes
\pgfplotsset{every axis/.append style={legend style={fill=white,fill opacity=0.92,text opacity=1,draw=black!15,inner sep=3pt}}}
% étiquettes d'arêtes (syntaxe edge["..."]) avec fond blanc
\tikzset{every edge quotes/.append style={fill=white,inner sep=1.2pt}}
%% FIGFIX-END
\begin{document}

\pagegarde{La période de mandat}{Histoire}{1re Tech}{Histoire – classe de Première technique}{N06}{2 séances (fiche de progression harmonisée)}{Cette leçon étudie la partition du Cameroun entre la France et la Grande-Bretagne après la Première Guerre mondiale, dans le cadre du système des mandats de la Société des Nations. Elle analyse ensuite les conditions de vie des Camerounais dans les deux territoires sous administration française et anglaise entre 1922 et 1946.
\vspace{4pt}
\begin{multicols}{2}\small
\begin{itemize}
\item À la fin de cette leçon, je sais expliquer comment la défaite allemande de 1918 a entraîné la partition du Cameroun
\item À la fin de cette leçon, je sais définir le système des mandats de la SDN et situer le Cameroun dans la catégorie des mandats B
\item À la fin de cette leçon, je sais décrire sur une carte le partage du Cameroun entre la France et la Grande-Bretagne
\item À la fin de cette leçon, je sais comparer les modes d'administration française (administration directe) et anglaise (indirect rule) au Cameroun
\item À la fin de cette leçon, je sais décrire la vie quotidienne des Camerounais : travail forcé, impôts, réquisitions, éducation, santé
\item À la fin de cette leçon, je sais interpréter des documents (cartes, photographies, données chiffrées) pour établir une notion historique
\end{itemize}\end{multicols}}

\begin{sommaire}
\begin{multicols}{2}
\begin{enumerate}
\item De la guerre à la partition du Cameroun (1914-1922)
\item Le partage territorial : le Cameroun français et le Cameroun britannique
\item Les modes d'administration : administration directe et indirect rule
\item La vie économique et sociale sous mandat français
\item La vie dans le Cameroun sous administration britannique
\item Bilan et transition : du mandat à la tutelle (1939-1946)
\item Activité d'intégration
\item Méthodes et exercices résolus
\item Exercices
\item Corrigés des exercices
\item Fiche bilan
\end{enumerate}
\end{multicols}
\end{sommaire}

\begin{objectifs}
\begin{itemize}
\item À la fin de cette leçon, je sais expliquer comment la défaite allemande de 1918 a entraîné la partition du Cameroun
\item À la fin de cette leçon, je sais définir le système des mandats de la SDN et situer le Cameroun dans la catégorie des mandats B
\item À la fin de cette leçon, je sais décrire sur une carte le partage du Cameroun entre la France et la Grande-Bretagne
\item À la fin de cette leçon, je sais comparer les modes d'administration française (administration directe) et anglaise (indirect rule) au Cameroun
\item À la fin de cette leçon, je sais décrire la vie quotidienne des Camerounais : travail forcé, impôts, réquisitions, éducation, santé
\item À la fin de cette leçon, je sais interpréter des documents (cartes, photographies, données chiffrées) pour établir une notion historique
\item À la fin de cette leçon, je sais construire une carte et un croquis montrant la partition du Cameroun
\item À la fin de cette leçon, je sais rédiger et justifier une conclusion argumentée à partir de documents
\item À la fin de cette leçon, je sais appliquer les notions de la leçon à des situations concrètes de la vie courante au Cameroun
\end{itemize}
\end{objectifs}

\begin{prerequis}
\begin{itemize}
\item La colonisation allemande du Kamerun (1884-1916) vue en classe de Troisième
\item Les causes et le déroulement de la Première Guerre mondiale (1914-1918)
\item Les campagnes militaires franco-britanniques au Cameroun (1914-1916)
\item La notion de colonisation et de traité international
\item La localisation des grandes régions du Cameroun actuel
\end{itemize}
\end{prerequis}

\newpage

\coursec{De la guerre à la partition du Cameroun (1914-1922)}

\textbf{Situation d'accroche.} En 1919, à Yaoundé, un commerçant camerounais qui payait ses impôts en marks allemands se retrouve du jour au lendemain sous l'autorité d'un administrateur français, tandis que son cousin de Bamenda, à quelques centaines de kilomètres à l'ouest, doit désormais répondre à un officier britannique. La monnaie change, la langue de l'administration change, les lois changent. Pourtant, \cle{aucun Camerounais n'a été consulté}. Comment un même peuple a-t-il pu être partagé entre deux puissances étrangères ? Quelles conséquences ce partage a-t-il eues sur la vie quotidienne des populations ? C'est à ces questions que répond l'étude de la \cle{période de mandat} (1922-1946), dont cette première section pose les fondations : la guerre, la conquête et le partage juridique du territoire.

\courssub{Rappel : le Kamerun allemand (1884-1914)}

Le 12 juillet 1884, les commerçants allemands des maisons Woermann et Jantzen \& Thormählen font signer aux chefs Duala de Douala (les rois Bell et Akwa) un traité de protectorat : le \cle{Kamerun} devient un protectorat allemand. En trente ans, les Allemands étendent leur domination bien au-delà de la côte : après l'accord franco-allemand de 1911 (qui règle la crise du Maroc, dite « coup d'Agadir »), le Kamerun atteint son extension maximale, s'étirant \emph{du lac Tchad au fleuve Congo}, grâce à des annexions au détriment du Congo français (le « Neukamerun »).

L'administration allemande installe sa capitale d'abord à Douala, puis à \cle{Yaoundé} à partir de 1901. Elle développe les plantations (cacao, caoutchouc, huile de palme), lance le chemin de fer (lignes du Nord et du Centre) et fonde des postes administratifs comme \cle{Garoua} au nord. Mais cette domination repose sur la contrainte : impôts, travail forcé, répression des résistances (révoltes des peuples du Sud et de l'Ouest).

\begin{attention}[Ce qu'il faut retenir du Kamerun allemand]
En 1914, le Kamerun est \textbf{un seul territoire} sous une seule administration. Le partage franco-britannique n'existe pas encore : il est une conséquence directe de la Première Guerre mondiale. Ne pas projeter en 1914 la frontière de 1922.
\end{attention}

\begin{popfigure}
\begin{center}
\begin{tikzpicture}[scale=1.0]
% Contour très simplifié du Kamerun 1911
\draw[thick,popink,fill=popgoldL] (0.2,0) -- (1.6,-0.5) -- (3.2,-0.3) -- (4.6,0.4) -- (5.4,1.6) -- (5.0,3.0) -- (4.2,4.2) -- (3.0,5.0) -- (2.0,5.4) -- (1.2,4.6) -- (0.6,3.4) -- (0.0,2.2) -- (0.2,0) -- cycle;
% Côte atlantique
\draw[very thick,popblue] (0.2,0) -- (0.0,2.2) -- (0.6,3.4);
\node[popblue,font=\small,rotate=75] at (-0.45,1.7) {Océan Atlantique};
% Villes
\fill[popink] (0.35,1.5) circle (1.6pt); \node[right,font=\small] at (0.45,1.5) {Douala};
\fill[popink] (1.9,1.1) circle (1.6pt); \node[right,font=\small] at (2.0,1.1) {Yaoundé};
\fill[popink] (2.6,4.3) circle (1.6pt); \node[right,font=\small] at (2.7,4.3) {Garoua};
% Indications de limites
\node[font=\small\itshape,popmuted] at (2.9,5.6) {vers le lac Tchad};
\node[font=\small\itshape,popmuted] at (5.6,0.2) {vers le Congo};
\node[font=\small\itshape,popmuted] at (-0.3,4.4) {vers le Nigeria};
% Nord
\draw[-{Stealth},thick,popdark] (5.6,4.6) -- (5.6,5.3); \node[font=\small] at (5.6,5.5) {N};
\node[font=\bfseries,popdark] at (2.6,2.6) {KAMERUN (1911)};
\end{tikzpicture}
\end{center}
\legende{Croquis schématique du Kamerun allemand à son extension maximale (1911) : un territoire unique s'étendant du golfe de Guinée vers le lac Tchad et le Congo. Villes principales : Douala (port), Yaoundé (capitale), Garoua (Nord). Contour simplifié, sans précision cartographique.}
\end{popfigure}

\courssub{Le Cameroun dans la Première Guerre mondiale (1914-1916)}

Quand la Première Guerre mondiale éclate en août 1914, l'Afrique devient un champ de bataille secondaire : les Alliés (France, Grande-Bretagne, Belgique) veulent s'emparer des colonies allemandes. Le Kamerun, encerclé par le Congo français, la Guinée équatoriale espagnole (neutre) et le Nigeria britannique, est attaqué sur plusieurs fronts.

\begin{definition}[Campagnes conjointes franco-britanniques]
De 1914 à 1916, les troupes françaises (parties du Congo et du Tchad) et britanniques (parties du Nigeria), aidées de tirailleurs africains, mènent des \cle{campagnes militaires coordonnées} contre les forces allemandes du Kamerun. Les opérations sont difficiles : climat, maladies, relief forestier et résistance allemande ralentissent la progression alliée.
\end{definition}

\begin{savaistu}[Les chefs de la conquête]
C'est le général britannique \cle{Charles Macpherson Dobell} (commandant les forces alliées sur la côte) et le général français \cle{Joseph Gaudérique Aymerich} (commandant les colonnes venues du Congo et du Tchad) qui dirigèrent les principales campagnes contre les Allemands au Kamerun. Les troupes coloniales africaines — tirailleurs côté français, soldats de la \emph{West African Frontier Force} côté britannique — fournirent l'essentiel des combattants : des milliers d'Africains perdirent la vie dans cette guerre qui n'était pas la leur.
\end{savaistu}

L'étape décisive est la \cle{chute de Yaoundé} en \cle{janvier 1916} : les forces franco-britanniques convergent sur la capitale allemande. Le gouverneur Karl Ebermaier, encerclé, choisit alors de fuir avec ses troupes et une partie de la population allemande vers la \cle{Guinée espagnole} (Río Muni), territoire neutre, où ils sont internés. En février 1916, la dernière garnison allemande de Mora (Extrême-Nord) capitule : la conquête du Kamerun est achevée.

\begin{exemplebox}[Une observation qui fait réfléchir]
Dès 1916, le territoire est occupé par \emph{deux} armées victorieuses. Question de démarche : que faire d'une colonie conquise à deux ? Chaque occupant administre la zone qu'il contrôle militairement. C'est la logique du \emph{fait accompli} : la répartition militaire de 1916 préfigure le partage juridique de 1922. Observer qui occupe quoi permet donc de prévoir qui obtiendra quoi.
\end{exemplebox}

\courssub{L'occupation militaire et le partage provisoire de 1916}

Dès 1916, la France et la Grande-Bretagne concluent un \cle{accord de partage provisoire} du Kamerun (accord dit « de Picot », du nom du diplomate français Georges Picot associé aux négociations franco-britanniques), qui organise la répartition du territoire conquis :
\begin{itemize}
\item la France administre la plus grande partie du territoire (l'Est, le Centre, le Sud et le Nord), soit environ les quatre cinquièmes ;
\item la Grande-Bretagne administre deux bandes occidentales le long de la frontière du Nigeria (le Nord-Ouest autour de Bamenda et la zone côtière autour de Victoria/Buea).
\end{itemize}

Ce partage est \emph{provisoire} : il attend la confirmation de la paix. Mais sur le terrain, chaque puissance agit déjà comme chez elle : la France rattache administrativement sa zone au Congo français, les Britanniques administrent la leur depuis le Nigeria.

\courssub{Le traité de Versailles (1919) : l'Allemagne renonce à ses colonies}

La guerre se termine par l'armistice du 11 novembre 1918. La paix est signée au château de Versailles : le \cle{traité de Versailles} du \cle{28 juin 1919} règle le sort de l'Allemagne vaincue. Son \cle{article 119} est décisif pour le Cameroun :

\begin{aretenir}[Article 119 du traité de Versailles]
« L'Allemagne renonce à tous ses droits et titres sur ses possessions d'outre-mer » au profit des principales puissances alliées et associées. Autrement dit : l'Allemagne perd définitivement ses colonies, dont le Kamerun. Mais le traité ne dit pas encore \emph{qui} administrera ces territoires ni \emph{comment} : c'est le rôle de la Société des Nations.
\end{aretenir}

\courssub{La Société des Nations et le système des mandats}

\begin{definition}[La Société des Nations (SDN)]
La \cle{Société des Nations} (SDN) est une organisation internationale créée en 1919 (Pacte annexé au traité de Versailles, siège à Genève) pour maintenir la paix par la négociation et prévenir une nouvelle guerre mondiale. Elle est l'\cle{ancêtre de l'ONU} (créée en 1945).
\end{definition}

L'\cle{article 22 du Pacte} de la SDN institue le \cle{système des mandats}. L'idée officielle : les peuples des colonies allemandes et de l'Empire ottoman « ne sont pas encore capables de se diriger eux-mêmes » ; on confie donc leur « tutelle » à des nations « avancées », qui doivent les guider vers le développement. C'est le principe de la \emph{mission sacrée de civilisation}.

\begin{definition}[Le mandat]
Un \cle{mandat} est un territoire placé par la SDN sous l'administration d'une puissance tutrice (dite « puissance mandataire »), qui doit \cle{rendre compte chaque année} à la \cle{Commission permanente des mandats}, organe de contrôle siégeant à Genève. Le mandat comporte des obligations : interdiction de fortifier le territoire, respect des populations, liberté religieuse.
\end{definition}

L'article 22 distingue \cle{trois catégories de mandats} selon le « degré de développement » supposé des populations :

\begin{popfigure}
\begin{center}
\begin{tikzpicture}[scale=1.0]
% Mandat A
\draw[thick,popgreen,fill=popgreenL,rounded corners] (0,0) rectangle (4.2,3.2);
\node[font=\bfseries,popdark] at (2.1,2.8) {Mandat A};
\node[font=\small,text width=3.8cm,align=center] at (2.1,1.5) {Indépendance « proche »\\[2pt] Ex. : Irak, Syrie, Liban (ex-Empire ottoman)};
% Mandat B
\draw[thick,poporange,fill=poporangeL,rounded corners] (4.8,0) rectangle (9.0,3.2);
\node[font=\bfseries,popdark] at (6.9,2.8) {Mandat B};
\node[font=\small,text width=3.8cm,align=center] at (6.9,1.5) {Administration directe du mandataire\\[2pt] Ex. : \textbf{Cameroun}, Togo, Tanganyika, Rwanda-Urundi};
% Mandat C
\draw[thick,poppurple,fill=poppurpleL,rounded corners] (9.6,0) rectangle (13.8,3.2);
\node[font=\bfseries,popdark] at (11.7,2.8) {Mandat C};
\node[font=\small,text width=3.8cm,align=center] at (11.7,1.5) {Administré comme partie intégrante du mandataire\\[2pt] Ex. : Sud-Ouest africain, îles du Pacifique};
\end{tikzpicture}
\end{center}
\legende{Les trois catégories de mandats instituées par l'article 22 du Pacte de la SDN. Le Cameroun est classé en catégorie B : administré par la puissance mandataire sous contrôle annuel de la SDN, sans promesse d'indépendance proche.}
\end{popfigure}

Le Cameroun est classé en \cle{mandat de catégorie B} : la puissance mandataire y administre directement le territoire, mais sous le contrôle de la SDN et avec l'obligation de rapports annuels.

\begin{attention}[Erreur fréquente : le mandat n'est pas une annexion]
Le territoire sous mandat \textbf{n'appartient pas juridiquement} à la puissance mandataire. La France et la Grande-Bretagne \emph{administrent} le Cameroun sous le contrôle de la SDN ; elles ne peuvent pas l'annexer, ni y lever des troupes pour la métropole, ni le fortifier. Dans une copie, ne jamais écrire « le Cameroun devient une colonie française en 1922 » : c'est un territoire \emph{sous mandat}, statut distinct de la colonie.
\end{attention}

\courssub{La confirmation des mandats : 20 juillet 1922}

Après la conférence de San Remo (1920) et des négociations franco-britanniques, le Conseil de la SDN \cle{confirme les mandats le 20 juillet 1922} :
\begin{itemize}
\item la \cle{France} obtient le mandat sur environ les \cle{4/5 du territoire} : le « Cameroun français » ;
\item la \cle{Grande-Bretagne} obtient environ \cle{1/5 du territoire} : les « British Cameroons », formés de \cle{deux bandes} (Nord et Sud) le long de la \cle{frontière nigériane}, administrées comme une annexe du Nigeria.
\end{itemize}

\begin{exemplebox}[Exemple chiffré : le partage en superficie]
La France administre environ \cle{432 000 km\textsuperscript{2}} (Cameroun français), la Grande-Bretagne environ \cle{88 000 km\textsuperscript{2}} (British Cameroons). Vérification : $432 + 88 = 520$ milliers de km\textsuperscript{2} au total ; la part française est $\dfrac{432}{520} \approx 0,83$, soit environ $83\,\%$ du territoire — c'est bien l'ordre de grandeur des « quatre cinquièmes » ($4/5 = 80\,\%$).
\end{exemplebox}

\begin{popfigure}
\begin{center}
\begin{tikzpicture}[scale=1.0]
\draw[thick,popdark] (0,0) -- (12,0);
\foreach \x/\y in {0/1914, 2/1916, 4.5/1918, 6.5/1919, 9/1920, 12/1922}{
  \draw[thick,popdark] (\x,-0.08) -- (\x,0.08);
  \node[below,font=\small] at (\x,-0.1) {\y};
}
% Événements
\draw[-{Stealth},popblue,thick] (0,0) -- (0,0.8);
\node[above,font=\small,text width=2.6cm,align=center] at (0,0.8) {Août 1914 : début des campagnes franco-britanniques};
\draw[-{Stealth},popblue,thick] (2,0) -- (2,1.6);
\node[above,font=\small,text width=2.8cm,align=center] at (2,1.6) {Janv. 1916 : chute de Yaoundé ; partage provisoire};
\draw[-{Stealth},popmuted,thick] (4.5,0) -- (4.5,0.8);
\node[above,font=\small,text width=2.4cm,align=center] at (4.5,0.8) {Nov. 1918 : armistice};
\draw[-{Stealth},poporange,thick] (6.5,0) -- (6.5,1.6);
\node[above,font=\small,text width=2.8cm,align=center] at (6.5,1.6) {28 juin 1919 : Versailles, art. 119 ; création de la SDN};
\draw[-{Stealth},popgreen,thick] (12,0) -- (12,1.6);
\node[above,font=\small,text width=3.2cm,align=center] at (12,1.6) {20 juillet 1922 : confirmation des mandats (France 4/5, G.-B. 1/5)};
\end{tikzpicture}
\end{center}
\legende{Frise chronologique : de la conquête militaire (1914-1916) au partage juridique du Cameroun sous mandats de la SDN (1922).}
\end{popfigure}

\begin{exoresolu}[Exercice : analyser un document]
\textbf{Document.} Extrait de l'article 22 du Pacte de la SDN (1919) : « Aux peuples qui ne sont pas encore capables de se diriger eux-mêmes […] la tutelle de ces peuples doit être confiée à des nations avancées […]. Cette tutelle doit être exercée par elles en qualité de mandataires et au nom de la Société. »\\
\textbf{Questions.} 1) Qui exerce la tutelle des peuples du Cameroun après 1919 ? 2) Au nom de qui cette tutelle est-elle exercée ? 3) Pourquoi peut-on dire que le mandat n'est pas une annexion ?
\tcblower
\textbf{Corrigé.} 1) La tutelle est exercée par des « nations avancées » : pour le Cameroun, la France (4/5 du territoire) et la Grande-Bretagne (1/5), confirmées comme puissances mandataires le 20 juillet 1922. 2) Elle est exercée « au nom de la Société » des Nations : les mandataires agissent comme délégués de la SDN et doivent lui rendre compte chaque année devant la Commission permanente des mandats. 3) Le texte précise « en qualité de mandataires » : la puissance n'est pas propriétaire du territoire, elle en est seulement l'administratrice sous contrôle international. Il n'y a donc pas transfert de souveraineté, contrairement à une annexion.
\end{exoresolu}

\begin{aretenir}[L'essentiel de la section]
\begin{itemize}
\item 1884-1914 : le Kamerun est un territoire allemand unique, du Tchad au Congo.
\item 1914-1916 : campagnes franco-britanniques ; chute de Yaoundé (janvier 1916) ; fuite des Allemands vers la Guinée espagnole ; partage provisoire (1916).
\item 1919 : traité de Versailles (art. 119) — l'Allemagne renonce à ses colonies ; création de la SDN et système des mandats (art. 22, catégories A, B, C).
\item 20 juillet 1922 : la SDN confirme les mandats — France : 4/5 (environ 432 000 km\textsuperscript{2}) ; Grande-Bretagne : 1/5 (environ 88 000 km\textsuperscript{2}), en deux bandes le long du Nigeria.
\end{itemize}
\end{aretenir}

\begin{exercice}[Pour s'entraîner]
1) Expliquez en cinq lignes pourquoi la partition du Cameroun est « une conséquence de la Première Guerre mondiale ». 2) Construisez un tableau à deux colonnes comparant le Cameroun français et les British Cameroons en 1922 (superficie, puissance mandataire, position géographique). 3) Rédigez un paragraphe argumenté répondant à la problématique : « Comment un même peuple a-t-il pu être partagé entre deux puissances étrangères sans être consulté ? »
\end{exercice}

\coursec{Le partage territorial : le Cameroun français et le Cameroun britannique}

La signature du traité de Versailles en juin 1919 entérine la défaite allemande, mais elle ne règle pas immédiatement le sort des colonies. Pour le Cameroun, la conférence de paix valide \emph{de facto} le partage militaire opéré sur le terrain entre troupes françaises et britanniques depuis 1916. C'est dans ce contexte que se dessine, à Londres au printemps 1919, une frontière qui va structurer l'histoire politique, administrative et humaine du territoire pour plus de quatre décennies. Comprendre ce partage, c'est saisir l'origine des disparités régionales qui marquent encore le Cameroun contemporain.

\courssub{La ligne Milner-Simon : naissance d'une frontière artificielle}

Le 10 juillet 1919, les représentants de la France (Simon) et de la Grande-Bretagne (Milner) signent à Londres un accord définissant la ligne de démarcation entre leurs zones d'occupation respectives. Cette \textbf{ligne Milner-Simon} ne suit aucune frontière naturelle majeure (fleuve, crête de montagne) ni aucune limite ethnique préexistante. Elle résulte d'un compromis géostratégique : la France obtient la partie la plus vaste et la plus riche (accès à la mer par Douala, bassin du Sanaga, terres de l'Adamaoua), tandis que la Grande-Bretagne sécurise la continuité territoriale avec sa colonie du Nigeria, au nord comme au sud du Bénoué.

\begin{popfigure}
\begin{tikzpicture}[scale=0.85]
% --- Coordonnées de base (contour simplifié du Kamerun 1914) ---
% Ouest (Nigeria) : x=0. Nord (Tchad/Lac) : y=10. Est (RCA/Congo) : x=11. Sud (Congo/Gabon/Mer) : y=0.
\coordinate (WN) at (0,10);   % Nord-Ouest (tripoint Nigeria-Tchad-Cameroun)
\coordinate (WS) at (0,0);    % Sud-Ouest (côte/mer)
\coordinate (EN) at (10,10);  % Nord-Est (Tchad)
\coordinate (EE) at (11,2);   % Est (RCA)
\coordinate (ES) at (11,0);   % Sud-Est (Congo)

% Frontière Nord (Nigeria-Tchad) : WN -- EN
% Frontière Est (Tchad-RCA-Congo) : EN -- EE -- ES
% Frontière Sud (Congo-Gabon-Mer) : ES -- WS
% Frontière Ouest (Nigeria) : WS -- WN

% Ligne Milner-Simon : P1 (sur frontière Nord) -> P2 (sur frontière Est)
\coordinate (P1) at (2,10);   % Sur frontière Nord, près du Logone/Lac Tchad
\coordinate (P2) at (10,2);   % Sur frontière Est (EE-ES), près du Sangha

% Ligne Bénoué (séparation Nord/Sud britannique) : horizontale y=6 de x=0 à intersection avec Milner-Simon
% Eq Milner-Simon : y = -x + 12. Pour y=6 -> x=6.
\coordinate (B) at (6,6);    % Intersection Bénoué / Milner-Simon
\coordinate (Bw) at (0,6);    % Bénoué sur frontière Ouest

% --- ZONES ---
% 1. CAMEROUN FRANÇAIS (Est de la ligne Milner-Simon)
\fill[popblue!15] (P1) -- (EN) -- (EE) -- (ES) -- (P2) -- cycle;
\node[popblue, font=\bfseries\large, align=center] at (7,7) {CAMEROUN\\ FRANÇAIS\\(Mandat français)};

% 2. NORTHERN CAMEROONS (UK) : Ouest de Milner-Simon ET Nord du Bénoué (y>6)
\fill[poporange!15] (P1) -- (WN) -- (Bw) -- (B) -- cycle;
\node[poporange!80!black, font=\bfseries, align=center] at (1,8.5) {NORTHERN\\ CAMEROONS\\(UK)};

% 3. SOUTHERN CAMEROONS (UK) : Ouest de Milner-Simon ET Sud du Bénoué (y<6)
\fill[popgreen!15] (B) -- (Bw) -- (WS) -- (ES) -- (P2) -- cycle;
\node[popgreen!70!black, font=\bfseries, align=center] at (3,3) {SOUTHERN\\ CAMEROONS\\(UK)};

% --- FRONTIÈRES ET TRAITS ---
% Contour pays
\draw[popdark, thick] (WN) -- (EN) -- (EE) -- (ES) -- (WS) -- cycle;
% Ligne Milner-Simon
\draw[popdark, very thick, dashed] (P1) -- (P2) node[midway, above right, sloped, font=\small\itshape] {Ligne Milner-Simon (1919)};
% Ligne Bénoué (limite Nord/Sud UK)
\draw[popmuted, thick, densely dotted] (Bw) -- (B) node[midway, above, font=\small\itshape] {Bénoué (limite N/S UK)};

% --- VILLES REPÈRES ---
% Françaises
\node[circle, fill=popdark, inner sep=1.5pt, label=above right:\small Yaoundé] at (6,4) {};
\node[circle, fill=popdark, inner sep=1.5pt, label=below right:\small Douala] at (8.5,1) {};
\node[circle, fill=popdark, inner sep=1.5pt, label=above:\small Ngaoundéré] at (6,6.5) {};
\node[circle, fill=popdark, inner sep=1.5pt, label=above:\small Garoua] at (7,8) {};
\node[circle, fill=popdark, inner sep=1.5pt, label=above left:\small Maroua] at (4,9) {};
\node[circle, fill=popdark, inner sep=1.5pt, label=right:\small Bertoua] at (9,5) {};
% Britanniques
\node[circle, fill=popdark, inner sep=1.5pt, label=left:\small Bamenda] at (0.5,7.5) {};
\node[circle, fill=popdark, inner sep=1.5pt, label=left:\small Buea] at (0.5,2) {};
\node[circle, fill=popdark, inner sep=1.5pt, label=above left:\small Yola (Nigeria)] at (-0.5,6) {};

% --- LÉGENDE ET ORIENTATION ---
\begin{scope}[shift={(11.5,9)}]
\fill[popblue!15] (0,0) rectangle (0.6,0.4); \node[right, font=\small] at (0.7,0.2) {Mandat français (84 \%)};
\fill[poporange!15] (0,-0.6) rectangle (0.6,-0.2); \node[right, font=\small] at (0.7,-0.4) {Northern Cameroons (UK)};
\fill[popgreen!15] (0,-1.2) rectangle (0.6,-0.8); \node[right, font=\small] at (0.7,-1.0) {Southern Cameroons (UK)};
\draw[popdark, very thick, dashed] (0,-1.8) -- (0.6,-1.8); \node[right, font=\small] at (0.7,-1.8) {Ligne Milner-Simon};
\draw[popmuted, thick, densely dotted] (0,-2.4) -- (0.6,-2.4); \node[right, font=\small] at (0.7,-2.4) {Bénoué (limite UK)};
\draw[->, thick, popdark] (0.3, -3.2) -- (0.3, -4.0) node[below, font=\small\bfseries] {N};
\end{scope}

% Échelle
\draw[thick, popdark] (0.5,-0.5) -- (2.5,-0.5) node[midway, below, font=\small] {200 km};
\end{tikzpicture}
\legende{Carte schématique du partage du Cameroun (1919). La ligne Milner-Simon (pointillés) sépare le mandat français (bleu, 84 \% du territoire) des deux enclaves britanniques : Northern Cameroons (orange, au nord du Bénoué) et Southern Cameroons (vert, au sud), toutes deux administrées depuis le Nigeria. Le trait pointillé fin marque le Bénoué, frontière interne britannique.}
\end{popfigure}

La carte ci-dessus illustre la rupture brutale de l'unité territoriale du \emph{Kamerun} allemand. La ligne Milner-Simon, tracée sur une carte à petite échelle sans connaissance fine du terrain, coupe arbitrairement des chefferies, des marchés et des parcours de transhumance. Elle attribue à la France \textbf{environ 415 000 km²} (soit 84 \% du Kamerun de 1914, ~495 000 km²) et à la Grande-Bretagne \textbf{environ 80 000 km²} (16 \%), répartis en deux blocs disjoints par le fleuve Bénoué.

\begin{methode}[Lire et construire une carte historique]
Pour analyser une carte historique comme celle du partage du Cameroun, suivez ces étapes :
\begin{enumerate}
\item \textbf{Titre et datation :} Identifier le sujet (ex. « Partage du Cameroun, 1919 ») et la date de la situation représentée.
\item \textbf{Orientation :} Repérer le Nord (flèche) pour situer les espaces relatifs (Nord/Sud, Est/Ouest).
\item \textbf{Échelle :} Utiliser l'échelle graphique pour estimer les distances réelles (ex. largeur du Southern Cameroons, enclavement du Northern).
\item \textbf{Légende :} Décoder les symboles (couleurs, hachures, traits, points) : ici, trois couleurs pour trois administrations, trait pointillé épais pour la frontière internationale de 1919, trait pointillé fin pour la limite administrative britannique.
\item \textbf{Analyse spatiale :} Comparer les surfaces, noter la continuité ou la discontinuité des zones (ici, deux zones UK non contiguës), localiser les capitales (Yaoundé, Buea, administration via Lagos/Kaduna).
\item \textbf{Sources et critique :} Identifier la source (ici, reconstruction schématique d'après accords Milner-Simon) et ses limites (contours simplifiés, frontière approximative, Neukamerun non représenté car restitué à l'AEF).
\end{enumerate}
\end{methode}

\begin{exemplebox}[Repères chiffrés du partage (1919)]
\begin{center}
\begin{tabularx}{\linewidth}{|l|X|X|}\hline
\textbf{Critère} & \textbf{Cameroun français} & \textbf{Cameroun britannique} \\ \hline
\textbf{Superficie} & ~ 415 000 km² (84 \%) & ~ 80 000 km² (16 \%) \\ \hline
\textbf{Population (est. 1920)} & ~ 1 500 000 hab. & ~ 400 000 hab. \\ \hline
\textbf{Capitales administratives} & Yaoundé (siège du Commissaire de la République) & Buea (Southern), Garoua (Northern) \emph{mais supervision depuis Lagos (Nigeria)} \\ \hline
\textbf{Organisation} & Circonscriptions (puis régions), administration directe & Provinces du Nigeria, \emph{Indirect Rule} \\ \hline
\textbf{Accès à la mer} & Douala, Kribi, Victoria (Limbe) & Victoria (Limbe) seulement pour le Sud ; enclavement pour le Nord \\ \hline
\end{tabularx}
\end{center}
\legende{Principales caractéristiques comparées des deux mandats au lendemain du partage. Le Kamerun de 1914 (~495 000 km²) est la base du partage ; le Neukamerun (céédé en 1911) est restitué à l'AEF (Afrique-Équatoriale française) et ne fait pas partie du mandat.}
\end{exemplebox}

\begin{popfigure}
\begin{tikzpicture}
\begin{axis}[
    xbar, width=\linewidth, height=6cm,
    xmin=0, xmax=100,
    ytick={1,2,3},
    yticklabels={Southern Cameroons (UK), Northern Cameroons (UK), Cameroun français (France)},
    xlabel={Part de la superficie du Kamerun 1914 (\%)},
    nodes near coords, nodes near coords align={horizontal},
    bar width=12pt,
    enlarge y limits=0.3,
    axis lines*=left,
    tick style={draw=none},
    xmajorgrids=true, grid style={popline},
]
\addplot[popgreen!70!black, fill=popgreen!20] coordinates {(6,1)};
\addplot[poporange!70!black, fill=poporange!20] coordinates {(10,2)};
\addplot[popblue, fill=popblue!20] coordinates {(84,3)};
\end{axis}
\end{tikzpicture}
\legende{Répartition proportionnelle de la superficie du Kamerun (1914) entre les puissances mandataires. La France reçoit l'immense majorité du territoire, des ressources minières, forestières et du réseau ferré existant.}
\end{popfigure}

\courssub{Le Cameroun sous mandat français : étendue, capitale et organisation}

Le territoire confié à la France par la Société des Nations (SDN) s'étend de l'océan Atlantique au lac Tchad, et du fleuve Campo au nord du 10e parallèle. C'est un bloc compact, cohérent géographiquement, qui hérite de l'essentiel des infrastructures allemandes (chemin de fer Nord-Sud amorcé, port de Douala, plantations de la côte).

\begin{definition}[Puissance mandataire]
Une \cle{puissance mandataire} est un État membre de la Société des Nations (SDN) chargé, en vertu de l'article 22 du Pacte de la SDN, d'administrer un territoire anciennement possédé par une puissance vaincue (ici l'Allemagne), au nom de la communauté internationale et au profit des populations indigènes. Elle doit rendre des comptes annuels à la Commission des mandats de la SDN. Le mandat n'est pas une annexion : la souveraineté reste en suspens.
\end{definition}

La France choisit \textbf{Yaoundé} comme capitale administrative (siège du Commissaire de la République) dès 1921, déplaçant le centre de gravité de la côte (Douala) vers l'intérieur, au cœur du pays Béti, pour mieux contrôler l'Adamaoua et le Nord. Le territoire est découpé en \textbf{circonscriptions} (Yaoundé, Douala, Ngaoundéré, Garoua, Maroua, Bertoua, etc.), chacune dirigée par un administrateur français (chef de circonscription) disposant de pouvoirs étendus : police, justice, impôt, travail forcé (\emph{prestations}), recrutement militaire. Cette \textbf{administration directe} s'appuie sur un corps de fonctionnaires français et sur des auxiliaires africains (chefs de poste, gardes-cercles, interprètes) nommés et révoqués par l'administration, sans nécessairement respecter la légitimité coutumière.

\courssub{Le Cameroun sous mandat britannique : deux zones, une administration nigériane}

Le sort du \textbf{British Cameroons} est radicalement différent. La Grande-Bretagne ne le conçoit pas comme une entité administrative distincte, mais comme une extension de sa colonie du Nigeria. Le territoire est morcelé en deux zones sans continuité terrestre, séparées par le fleuve Bénoué :

\begin{itemize}
\item Les \textbf{Northern Cameroons} (Nord) : vaste zone de savane, peuplée de Peuls, Mafa, Kapsiki, etc., rattachée à la \textbf{province du Nord du Nigeria} (capitale Kaduna, puis Zaria). Le chef-lieu local est Garoua.
\item Les \textbf{Southern Cameroons} (Sud) : zone forestière et de hauts plateaux (Grassfields), peuplée de Bamileke, Bamoun, Tikar, Bassa, etc., rattachée à la \textbf{province de l'Est du Nigeria} (capitale Enugu). Le chef-lieu local est Buea (ancienne capitale allemande), siège du \emph{Resident} britannique.
\end{itemize}

\begin{attention}[Ne pas confondre Northern et Southern Cameroons]
\emph{Erreur fréquente :} Traiter le « Cameroun britannique » comme un bloc unique.
\begin{itemize}
\item \textbf{Northern Cameroons} : Zone sahélienne, islamisée, économie pastorale et cotonnière, administrée depuis \textbf{Kaduna/Zaria} (Nord Nigeria). Majorité musulmane.
\item \textbf{Southern Cameroons} : Zone forestière/altitudes, christianisée tôt, économie de plantation (cacao, café, palmiers), administrée depuis \textbf{Enugu} (Est Nigeria). Majorité chrétienne/animiste.
\item \textbf{Conséquence majeure :} En 1961, le Nord vote le rattachement au Nigeria, le Sud vote la réunification avec le Cameroun français. Cette divergence trouve sa racine dans cette administration séparée.
\end{itemize}
\end{attention}

Le statut juridique est celui d'une \textbf{province du Nigeria britannique}. Le Gouverneur du Nigeria nomme des \emph{Residents} pour chaque zone. La législation nigériane (code pénal, code fiscal, organisation judiciaire) s'y applique \emph{mutatis mutandis}. Les Camerounais de ces zones deviennent des « sujets britanniques protégés » du Nigeria, sans citoyenneté propre, et leurs représentants siégeront plus tard à l'Assemblée législative de Lagos, puis à la Chambre des représentants fédérale à Lagos.

\courssub{Conséquences humaines : familles et ethnies divisées}

Le tracé de la ligne Milner-Simon opère une \textbf{coupure anthropologique majeure}. Il ignore les aires culturelles, linguistiques et matrimoniales.

\begin{exemplebox}[Le cas des Grassfields et du peuple Bamoun]
Le royaume Bamoun (autour de Foumban) se retrouve dans le Cameroun français, mais ses dépendances occidentales et nombre de chefferies Tikar voisines (Bankim, Banyo - zone de transition) basculent côté britannique (Southern Cameroons pour les premières, Northern pour Banyo/Bankim selon les ajustements ultérieurs).
\begin{itemize}
\item \textbf{Réseaux commerciaux :} Les marchés hebdomadaires (ex. marché de Bafoussam, de Foumban) drainent des producteurs des deux côtés de la nouvelle frontière. L'instauration de la douane et de laissez-passer perturbe l'écoulement du café, du cola, du sel, du bétail.
\item \textbf{Lien sociaux :} Les mariages inter-ethnies (Bamileke/Bamoun/Tikar) deviennent transfrontaliers. Les chefferies traditionnelles (Fon, Sultan) voient leur autorité territoriale tronquée.
\item \textbf{Exemple chiffré :} On estime qu'en 1930, plus de \textbf{200 000 personnes} (sur ~2 millions) vivent dans des groupes ethniques coupés par la frontière. Les Bamileke (français) et les Bamenda (britanniques) partagent la même langue (ghomala'/bamileke) et les mêmes ancêtres mythiques.
\end{itemize}
\end{exemplebox}

Au Nord, la frontière sépare les Peuls de l'Adamaoua (français) de ceux du Bénoué (britanniques), et coupe le lamidat de Rey-Bouba (français) de ses pâturages saisonniers côté britannique. La libre circulation des nomades (Bororo, Mbororo) et des pèlerins vers La Mecque (via le Nord britannique puis le Soudan français) devient soumise à des formalités policières.

\courssub{Conséquences économiques : rupture des circuits, double monnaie, double fiscalité}

L'unité économique du \emph{Kamerun} allemand (chemin de fer unique, port unique de Douala, monnaie unique : Mark puis Mark-Or) est brisée.

\begin{itemize}
\item \textbf{Rupture des circuits commerciaux :} Le réseau ferré (Douala-Nkongsamba, Douala-Yaoundé achevée en 1927 ; l'extension vers Ngaoundéré ne sera achevée qu'en 1974) reste entièrement en zone française. Le Nord britannique (coton, bétail) et le Sud britannique (cacao, bois, café) doivent exporter via le Nigeria (ports de Lagos, Port Harcourt) par des routes précaires ou le fleuve Bénoué (navigation saisonnière). Le coût de transport double, la compétitivité chute.
\item \textbf{Double monnaie :} Zone française : \textbf{Franc CFA} (créé en 1945, mais dès les années 1920 franc français / mark de convention). Zone britannique : \textbf{Livre sterling / Shilling nigérian}. Le change est imposé par les administrations, défavorable aux producteurs camerounais.
\item \textbf{Double fiscalité :} Impôt de capitation (français) vs \emph{Native Revenue Ordinance} (britannique, impôt sur les huttes/chefferies). Taux, assiette, modes de recouvrement (contrainte physique vs prélèvement par chefs coutumiers) diffèrent. Les commerçants transfrontaliers paient deux fois.
\end{itemize}

\begin{savaistu}[La voix des populations à Genève : le droit de pétition]
L'article 22 du Pacte de la SDN prévoit que les populations des territoires sous mandat peuvent adresser des \textbf{pétitions} au Secrétariat de la SDN à Genève. Dès les années 1920, des notables camerounais (ex. le prince \textbf{Alexandre Douala Manga Bell}, fils de Rudolf exécuté par les Allemands, ou des chefs de l'Ouest comme le Sultan Njoya) utilisent ce droit pour dénoncer : le travail forcé (\emph{prestations}), l'accaparement des terres (plantations), l'arbitraire des chefs de poste français, ou l'absorption dans l'administration nigériane côté britannique. Ces pétitions, bien que souvent classées sans suite par la Commission des mandats (dominée par les puissances mandataires), constituent une \textbf{source historique précieuse} : elles prouvent la conscience politique naissante et l'existence d'une opinion publique camerounaise transfrontalière.
\end{savaistu}

\courssub{Le contrôle de la Société des Nations : rapports et pétitions}

Le mandat n'est pas une souveraineté pleine. La France et la Grande-Bretagne ont l'obligation juridique de remettre chaque année un \textbf{rapport annuel} à la Commission permanente des mandats (CPM) à Genève. Ces rapports détaillent : budget, santé, éducation, travaux publics, justice indigène, recrutement militaire, production économique.

\begin{exoresolu}[Analyse d'un extrait de rapport annuel français (1935)]
\textbf{Énoncé :} Extrait du rapport 1935 : « La circonscription de Ngaoundéré a fourni 1 200 porteurs pour le ravitaillement du poste de Tibati. Le taux de mortalité parmi les porteurs a été de 4 \%. Les recettes de l'impôt de capitation ont progressé de 12 \%. »

\textbf{Question :} Que révèle cet extrait sur la réalité de l'administration directe française ?

\tcblower

\textbf{Solution détaillée :}
\begin{enumerate}
\item \textbf{Identification de la source :} Rapport officiel de la puissance mandataire (France) à la SDN, 1935. Source administrative, autocensure probable, langage bureaucratique.
\item \textbf{Analyse interne :}
\begin{itemize}
\item « 1 200 porteurs » : Recrutement forcé (\emph{prestations}) pour des besoins logistiques militaires/administratifs (absence de routes praticables, pas de camions en nombre suffisant).
\item « Taux de mortalité 4 \% » : Soit 48 morts officiels pour un seul convoi. Conditions inhumaines (charge excessive, distance, alimentation insuffisante, maladies). Chiffre probablement sous-estimé par l'administration.
\item « Recettes impôt +12 \% » : Pression fiscale accrue, indicateur de la priorité budgétaire (autofinancement de la colonie, « mise en valeur ») sur le bien-être des populations.
\end{itemize}
\item \textbf{Mise en perspective :} Ce rapport, lu à Genève, est souvent édulcoré. La SDN manque de moyens d'enquête sur place (pas d'inspecteurs permanents). Les pétitions des populations (v. boîte « Le savais-tu ? ») contredisent souvent le ton lissé des rapports officiels.
\item \textbf{Conclusion :} L'administration directe française repose sur la contrainte (travail forcé, impôt lourd) et l'extraction de ressources, masquée par un langage technocratique dans les rapports à la SDN.
\end{enumerate}
\end{exoresolu}

\begin{aretenir}[Bilan du partage territorial (1919-1939)]
\begin{itemize}
\item \textbf{Frontière artificielle :} Ligne Milner-Simon (1919), tracée sans concertation des populations, coupant ethnies (Bamileke/Bamenda, Bamoun/Tikar, Peuls), bassins de vie et circuits commerciaux.
\item \textbf{Deux modèles administratifs :} \textbf{Administration directe} française (fonctionnaires, circonscriptions, capitale Yaoundé, franc CFA) vs \textbf{Indirect Rule} britannique (intégration au Nigeria, chefferies instrumentalisées, livre sterling, capitales Lagos/Enugu/Kaduna).
\item \textbf{Deux Cameroun britanniques :} Northern (enclavé, pastoral, musulman, tourné vers le Nord Nigeria) et Southern (côtier/forestier, plantation, christianisé, tourné vers l'Est Nigeria).
\item \textbf{Héritage conflictuel :} Double monnaie, double fiscalité, double système éducatif (français/anglais), double système judiciaire. Ces divergences structurent la décolonisation (1960-1961) et les tensions actuelles (crise anglophone).
\item \textbf{Contrôle international formel :} Rapports annuels à la SDN + droit de pétition (peu effectif mais trace écrite majeure pour l'historien).
\end{itemize}
\end{aretenir}

La partition de 1919 n'est pas une simple ligne sur une carte : elle installe deux logiques d'État, deux langues officielles, deux systèmes juridiques et éducatifs sur un même espace historique. C'est cette dualité institutionnelle, bien plus que les différences ethniques, qui fera la complexité de la réunification de 1961 et la fragilité de l'unité nationale camerounaise. La section suivante examinera précisément comment ces deux administrations ont concrètement géré les populations : l'administration directe française face à l'\emph{Indirect Rule} britannique.

\coursec{Les modes d'administration : administration directe et indirect rule}

La partition du Cameroun en 1919, entérinée par les mandats de la Société des Nations (SDN) en 1922, ne se limite pas à un tracé de frontières : elle impose aux populations deux logiques politiques, juridiques et administratives radicalement différentes. Si la France et le Royaume-Uni partagent l'objectif affiché de « mise en valeur » du territoire au profit de la puissance mandataire, les instruments mis en œuvre pour y parvenir divergent profondément. La France applique une \emph{administration directe}, calquée sur le modèle algérien, où l'État colonial substitue ses propres agents aux autorités traditionnelles. Le Royaume-Uni privilégie l'\emph{indirect rule} (règlement indirect), théorisé par Frederick Lugard au Nigeria, qui prétend gouverner \emph{à travers} les institutions locales. Cette section analyse la structure, le fonctionnement et les conséquences de ces deux systèmes sur le quotidien des Camerounais.

\courssub{L'administration directe française : une pyramide centralisée}

\textbf{Intuition et principe.}\par
Dès 1916, l'administration militaire française met en place une structure hiérarchique rigide pour contrôler le territoire conquis. L'objectif est double : assurer la souveraineté française face aux revendications allemandes ou britanniques, et extraire les ressources nécessaires au financement de la colonie (politique de l'autofinancement ou « mise en valeur »). Contrairement à la logique britannique, l'administration française ne cherche pas à préserver les chefferies traditionnelles comme rouages de l'État ; elle les \emph{substitue} par des agents de l'État français, reléguant les chefs au rang d'auxiliaires subalternes.

\begin{definition}[Administration directe (modèle français)]
Système colonial dans lequel la puissance mandataire gouverne par l'intermédiaire de ses propres fonctionnaires (européens ou assimilés), placés à tous les échelons de la hiérarchie administrative. Les autorités traditionnelles sont soit écartées, soit transformées en \emph{chefs auxiliaires} nommés, révoqués et contrôlés par l'administration, sans autonomie politique ni judiciaire propre.
\end{definition}

\textbf{La hiérarchie administrative : le Commissaire de la République au sommet.}\par
Le mandat français est dirigé par un \textbf{Commissaire de la République} (titre porté par le gouverneur après 1922), résidant à Yaoundé. Il concentre les pouvoirs exécutif, législatif (par décrets) et judiciaire (grâces). Il rend des comptes au Ministère des Colonies à Paris et au Conseil de la SDN (rapports annuels).

Le territoire est découpé en \textbf{circonscriptions} (une vingtaine dans les années 1930), chacune dirigée par un \textbf{Commissaire de circonscription} (souvent un administrateur de carrière). Les circonscriptions sont subdivisées en \textbf{subdivisions} (chef-lieu : poste administratif), dirigées par des \textbf{Chefs de subdivision} (administrateurs adjoints ou officiers de réserve). C'est à ce niveau que s'exerce l'administration de proximité : perception de l'impôt, recrutement de la main-d'œuvre (prestations), justice sommaire.

\begin{popfigure}
\begin{tikzpicture}[
    level distance=1.8cm,
    level 1/.style={sibling distance=5.5cm},
    level 2/.style={sibling distance=2.8cm},
    level 3/.style={sibling distance=1.5cm},
    level 4/.style={sibling distance=1.0cm},
    every node/.style={draw=popblue, thick, rounded corners, fill=popblueL, minimum width=2.2cm, minimum height=0.9cm, font=\small, text centered, text width=2.2cm},
    edge from parent/.style={draw=popdark, thick, -{Stealth[length=3mm]}},
    edge from parent path={(\tikzparentnode.south) -- ++(0,-0.4) -| (\tikzchildnode.north)}
]
\node[align=center]{Ministère des Colonies\\ (Paris)}
    child { node[align=center]{Commissaire de la République\\ (Yaoundé)}
        child { node[align=center]{Commissaire de\\ Circonscription 1}
            child { node[align=center]{Chef de\\ Subdivision A}
                child { node {Chef de Canton} child { node {Chef de Village} } }
                child { node {Chef de Canton} child { node {Chef de Village} } }
            }
            child { node[align=center]{Chef de\\ Subdivision B} }
        }
        child { node[align=center]{Commissaire de\\ Circonscription 2} }
        child { node {$\cdots$ (env. 20 circonscriptions)} }
    };
\end{tikzpicture}
\legende{Hiéarchie de l'administration directe française au Cameroun (années 1930). Le pouvoir descend du sommet vers la base sans retour horizontal. Les chefs de canton et de village sont des agents subalternes de l'État.}
\end{popfigure}

\textbf{La politique d'association et l'autofinancement : « La colonie doit payer la colonie ».}\par
Le principe de la \emph{mise en valeur} impose que le budget du Cameroun français soit équilibré par ses propres ressources (impôts, droits de douane, concessions forestières et minières). L'administration investit peu dans l'éducation ou la santé (laissées aux missions), mais massivement dans les infrastructures d'extraction : chemin de fer (Douala–Yaoundé–Ngaoundéré), ports, routes d'évacuation. Cette pression fiscale explique la dureté de la perception de l'impôt de capitation (impôt personnel) et le recours massif aux \emph{prestations} (travaux forcés non rémunérés ou sous-payés) pour l'entretien des pistes et la construction des bâtiments administratifs.

\textbf{Les chefs auxiliaires : auxiliaires de l'administration, non représentants du peuple.}\par
Pour combler le manque de fonctionnaires européens (environ 400 administrateurs pour 2 millions d'habitants dans les années 1930), l'administration institue une chefferie « officielle » à trois niveaux :
\begin{itemize}
    \item \textbf{Chefs supérieurs} (régions ethniques vastes, ex. : Sultanat du Bamoun, Lamidat de Rey Bouba) : intermédiaires de prestige, mais surveillés de près.
    \item \textbf{Chefs de canton} (regroupement de villages) : pivots de la perception de l'impôt et du recrutement de la main-d'œuvre. Ils sont \emph{nommés} par arrêté du Commissaire de la République, souvent choisis hors des lignées traditionnelles pour leur loyauté.
    \item \textbf{Chefs de village} : base de la pyramide, chargés du recensement, de la corvée, de l'ordre public.
\end{itemize}
Ces chefs perçoivent une commission sur l'impôt perçu (5 à 10 \%) et des avantages en nature. Leur légitimité traditionnelle est souvent brisée : ils deviennent des \emph{agents de l'État} perçus comme des oppresseurs par leurs propres « administrés ».

\begin{exemplebox}[Chiffre clé : le maillage territorial dans les années 1930]
En 1935, le Cameroun français compte \textbf{19 circonscriptions} (Douala, Yaoundé, Ngaoundéré, Garoua, Maroua, Bertoua, Bafoussam, etc.), subdivisées en \textbf{environ 80 subdivisions} et plus de \textbf{300 cantons}. Ce maillage dense permet un contrôle territorial effectif mais coûteux, financé par l'impôt de capitation (6 à 15 francs/or par homme valide selon les régions) et le travail forcé.
\end{exemplebox}

\textbf{Le régime de l'indigénat : une justice d'exception.}\par
Le fondement juridique de la domination repose sur le \emph{Code de l'indigénat} (décrets de 1924, 1928, appliqués au Cameroun par arrêtés). Les Camerounais sont des \emph{sujets français}, non des \emph{citoyens}. Ils n'ont pas de droits politiques (pas de vote, pas d'éligibilité) et sont justiciables de tribunaux d'exception.

\begin{definition}[Indigénat (zone française)]
Régime juridique d'exception appliqué aux « sujets » des colonies. Il permet aux administrateurs (chefs de subdivision, commissaires de circonscription) de prononcer, \textbf{sans procès ni avocat}, des peines de prison (jusqu'à 15 jours, renouvelables), d'amende, de déportation ou de réquisition de biens/travaux, pour des infractions vaguement définies (vagabondage, insoumission, manque de respect à l'administration). La « justice » est administrative, expéditive et discrétionnaire.
\end{definition}

\courssub{L'indirect rule britannique : gouverner par les chefferies traditionnelles}

\textbf{Intuition et principe.}\par
Le \emph{British Cameroons} (Nord et Sud) est administré comme une extension du Nigeria. Le gouverneur du Nigeria (siège à Lagos, puis Kaduna) nomme un \textbf{Resident} (pour le Nord) et un \textbf{Senior Resident} (pour le Sud, basé à Buea) qui supervisent des \textbf{District Officers} (DO). Mais la spécificité majeure est le recours aux \emph{Native Authorities} (Autorités indigènes).

\begin{definition}[Indirect Rule (Règlement indirect)]
Système colonial britannique théorisé par Lord Lugard : l'administration coloniale gouverne \emph{à travers} les institutions politiques, judiciaires et fiscales préexistantes (chefferies, conseils de notables, coutumes). Les chefs traditionnels reconnus (\emph{Native Authorities}) conservent leurs prérogatives (perception de l'impôt, police, justice coutumière) sous la \emph{supervision} (contrôle indirect) des officiers de district britanniques.
\end{definition}

\textbf{Le rôle pivot du Nigeria : une annexe administrative.}\par
Le Cameroun britannique n'a pas de budget propre, pas de service de santé ou d'éducation distincts du Nigeria. Les recettes (impôts, douanes) partent au Trésor fédéral nigérian ; les dépenses sont décidées à Lagos/Kaduna. Le \emph{Northern Cameroons} est rattaché à la \emph{Northern Region} du Nigeria (capitales : Zungeru puis Kaduna) ; le \emph{Southern Cameroons} à la \emph{Eastern Region} (capitale : Enugu) puis devient une région quasi-fédérée en 1954. Cette intégration prive les Camerounais britanniques d'une identité administrative propre et les soumet aux politiques nigérianes (ex. : éducation en anglais, droit coutumier codifié par les Britanniques).

\textbf{Les Native Authorities : des chefs « traditionnels » renforcés.}\par
Au Nord (zone sahélienne, islamisée), les \emph{Lamibe} (chefs peuls) et le \emph{Sultan de Mandara} sont confirmés comme \emph{First Class Chiefs}. Ils perçoivent l'impôt (\emph{Jangali} sur le bétail, \emph{Kharaj} sur les terres), gèrent les tribunaux alkali (droit musulman) et les prisons natives. Au Sud (zone forestière, chefferies plus fragmentées), les Britanniques créent ou renforcent des \emph{Warrant Chiefs} (chefs à mandat) là où le pouvoir était diffus, ou s'appuient sur les \emph{Fons} (Bamileke, Grassfields) et les \emph{Obas} (zone côtière). Ces chefs disposent de \textbf{trésoreries locales} (Native Treasuries) financées par l'impôt local, qui paient les écoles, dispensaires, routes et salaires des courriers/criers. Ils ont donc un \emph{intérêt matériel} à collaborer.

\begin{attention}[Piège fréquent : confondre chefs auxiliaires (France) et Native Authorities (UK)]
\textbf{Erreur :} « Dans les deux zones, les chefs sont des marionnettes de l'administration. »
\textbf{Réalité :} En zone \textbf{française}, le chef de canton est un \emph{agent subalterne} révocable à tout moment, sans budget propre, exécutant les ordres du Chef de subdivision. En zone \textbf{britannique}, le \emph{Native Authority} (Fon, Lamido, Sultan) gère un \emph{budget propre} (Native Treasury), rend la justice selon la coutume (sous réserve de la \emph{repugnancy clause} : rien de « répugnant à la justice naturelle »), et négocie avec le District Officer. Son pouvoir est \textbf{réel, bien que borné} par la suprématie britannique.
\end{attention}

\courssub{Analyse comparative : centralisation française contre décentralisation britannique}

Le tableau ci-dessous synthétise les différences structurelles. Il ne s'agit pas d'opposer un « bon » et un « mauvais » système, mais de comprendre comment chaque logique façonne différemment le rapport à l'État, à la loi et à l'autorité chez les Camerounais.

\begin{popfigure}
\begin{tikzpicture}
% Styles
\tikzset{
    header/.style={fill=popblue, font=\bfseries\small, text=white, minimum height=1.0cm, align=center},
    headerA/.style={header, fill=popblue!80!poppurple},
    headerB/.style={header, fill=poporange!80!popdark},
    cell/.style={font=\small, minimum height=0.85cm, align=center, inner xsep=3pt},
    cellL/.style={cell, align=left, text width=3.8cm},
    cellC/.style={cell, text width=3.2cm},
    rowA/.style={fill=popblueL!50},
    rowB/.style={fill=poporangeL!50},
    grid/.style={draw=popline, thin}
}
% Dimensions
\def\colA{3.8cm}
\def\colB{3.2cm}
\def\colC{3.2cm}
\def\rowh{0.85cm}
% Header
\node[headerA, minimum width=\colA, grid] at (0,0) {Critère de comparaison};
\node[headerB, minimum width=\colB, grid, anchor=west] at (\colA,0) {Zone française (Admin. directe)};
\node[headerB, minimum width=\colC, grid, anchor=west] at ({\colA+\colB},0) {Zone britannique (Indirect Rule)};
% Rows
\foreach \i/\crit/\fr/\uk [count=\y from 1] in {
    {{1}/{Autorité centrale}/{Commissaire de la République (Yaoundé) — pouvoir discrétionnaire total}/{Gouverneur du Nigeria (Lagos/Kaduna) → Residents / District Officers}},
    {{2}/{Relais locaux}/{Chefs auxiliaires (nommés, révoqués, commissionnés) — agents de l'État}/{Native Authorities (chefs traditionnels reconnus) — budgets propres (Native Treasuries)}},
    {{3}/{Justice}/{Justice administrative (indigénat) + tribunaux français pour Européens}/{Tribunaux coutumiers (Alkali, Native Courts) supervisés par DO ; appel au High Court}},
    {{4}/{Fiscalité}/{Impôt de capitation + prestations (travail forcé) → Budget général colonie}/{Impôts locaux (Jangali, taux) → Native Treasuries (gérées par chefs) + budget fédéral Nigeria}},
    {{5}/{Éducation et Santé}/{Laissées aux missions (subventionnées), peu d'écoles d'État}/{Intégrées au système nigérian (missions + Native Authority schools), plus d'écoles primaires au Sud}},
    {{6}/{Code du travail}/{Code du travail indigène (1931) : encadre strict, bas salaires, travail forcé}/{Ordonnances nigérianes : travail forcé aboli plus tôt (1930s), salaires légèrement supérieurs au Sud}}
} {
    \pgfmathsetmacro{\yy}{-\y*0.85};
    \node[cellL, \ifodd\y rowA\else rowB\fi, grid, anchor=north west] at (0,\yy) {\crit};
    \node[cellC, \ifodd\y rowA\else rowB\fi, grid, anchor=north west] at (\colA,\yy) {\fr};
    \node[cellC, \ifodd\y rowA\else rowB\fi, grid, anchor=north west] at ({\colA+\colB},\yy) {\uk};
}
% Bottom line
\pgfmathsetmacro{\yy}{-7*0.85};
\draw[popline, thick] (0,\yy) -- ({\colA+\colB+\colC},\yy);
\end{tikzpicture}
\legende{Tableau comparatif des modes d'administration français et britannique au Cameroun sous mandat (années 1920–1930). Source : synthèse d'après rapports annuels SDN et archives coloniales.}
\end{popfigure}

\textbf{Synthèse de la comparaison.}\par
\begin{itemize}
    \item \textbf{Centralisation vs Décentralisation budgétaire :} La France recentre tout à Yaoundé (budget unique) ; le UK décentralise le budget vers les \emph{Native Treasuries}, créant une classe de chefs gestionnaires.
    \item \textbf{Droit unique vs Pluralisme juridique :} Zone française : code de l'indigénat + code civil pour les rares « citoyens » (originaires des « Quatre Communes » du Sénégal, quasi inexistants au Cameroun). Zone britannique : coexistence du droit anglais (pour les Européens, affaires commerciales), du droit coutumier codifié (Native Courts) et de la charia (Nord).
    \item \textbf{Relation chef/population :} En zone française, le chef est perçu comme l'\emph{émissaire de l'oppresseur} (recruteur de porteurs, percepteur d'impôt). En zone britannique, le chef est le \emph{garant de l'autonomie locale} face à l'administration lointaine de Lagos/Enugu, ce qui renforce parfois son autorité traditionnelle (surtout au Nord).
\end{itemize}

\courssub{Le statut des administrés : ni citoyens, ni sujets à part entière}

\textbf{Zone française : le sujet de l'indigénat.}\par
Le Camerounais est un \emph{sujet français}. Il doit allégeance à la France mais ne jouit pas des droits du citoyen (libertés d'association, de presse, de réunion, droit de vote). L'accès à la citoyenneté est théoriquement possible (naturalisation) mais les conditions sont draconiennes (renonciation au statut personnel coutumier, service militaire, diplôme, « mérite exceptionnel ») : \textbf{moins de 50 naturalisations} sur toute la période du mandat. La masse de la population vit sous le régime de l'indigénat : corvées (prestations), impôt de capitation, interdiction de circuler sans laissez-passer, répression des mouvements politiques (ex. : interdiction de l'USCC, Union des Syndicats Confédérés du Cameroun, en 1939).

\textbf{Zone britannique : le « British Protected Person » (BPP).}\par
Juridiquement, le Cameroun britannique est un \emph{territoire sous protection}, non une colonie de peuplement. Les habitants sont des \emph{British Protected Persons}. Ils ne sont pas des sujets britanniques (pas de passeport britannique, pas de droit de séjour au UK). Ils relèvent de la juridiction des \emph{Native Courts} pour le droit personnel (mariage, succession, terre) et du droit anglais pour les crimes graves et le commerce. Ce statut offre une marge de manœuvre : les élites du Sud (Bakweri, Duala migrants, Bamileke) utilisent les \emph{Native Authorities} et les tribunaux pour défendre leurs terres contre les domaines allemands séquestrés (gérés par le \emph{Custodian of Enemy Property}, futur \emph{Cameroon Development Corporation}).

\begin{exoresolu}[Analyse de document : Extrait du Rapport annuel du Commissaire de la République au Conseil de la SDN (1934)]
\textbf{Énoncé :} « La politique d'association, si elle a permis de maintenir l'ordre par l'intermédiaire des chefs traditionnels, se heurte à l'insuffisance du recrutement indigène dans les cadres moyens. L'impôt de capitation a été perçu à 98\% dans la circonscription de Yaoundé, grâce à l'activité des chefs de canton. Les prestations ont fourni 120 000 journées-hommes pour l'entretien de la route Yaoundé–Bertoua. »
\textbf{Question :} Identifiez les mécanismes de l'administration directe illustrés par ce texte et leurs limites.

\tcblower
\textbf{Solution détaillée :}
\begin{enumerate}
    \item \textbf{Mécanismes identifiés :}
    \begin{itemize}
        \item \emph{Politique d'association / Chefs auxiliaires :} « intermédiaire des chefs traditionnels », « activité des chefs de canton ». L'administration utilise la chefferie comme courroie de transmission (perception impôt, recrutement main-d'œuvre).
        \item \emph{Centralisation financière :} L'impôt (capitation) finance le budget local (route Yaoundé–Bertoua). Taux de perception élevé (98\%) = contrainte forte.
        \item \emph{Travail forcé (prestations) :} « 120 000 journées-hommes » = main-d'œuvre non salariée réquisitionnée administrativement pour infrastructure d'évacuation.
    \end{itemize}
    \item \textbf{Limites révélées :}
    \begin{itemize}
        \item \emph{Insuffisance des cadres moyens indigènes :} L'administration directe ne forme pas d'élite administrative locale (pas d'école d'administration, peu d'écoles primaires). Les chefs restent des exécutants, pas des gestionnaires formés.
        \item \emph{Contrainte vs Adhésion :} Le taux de 98\% trahit la contrainte (indigénat, prison pour défaut de paiement), non le consentement.
        \item \emph{Extractivisme :} La route relie Yaoundé (centre décision) à Bertoua (ressources forestières/minières), pas les villages entre eux.
    \end{itemize}
\end{enumerate}
\end{exoresolu}

\textbf{Transition vers la section 4.}\par
Ces deux systèmes administratifs, par leur fiscalité, leur justice et leur gestion de la main-d'œuvre, structurent en profondeur la vie économique et sociale. La section suivante examinera comment les Camerounais vivent concrètement cette domination : économie de plantation, travail forcé, migrations, résistance silencieuse et naissance d'un syndicalisme naissant en zone française ; économie de traite, chefferies renforcées et premières revendications politiques en zone britannique.

\coursec{La vie économique et sociale sous mandat français}

Nous venons de voir que la France administre son territoire sous mandat par l'\cle{administration directe} : des administrateurs français commandent, des chefs auxiliaires exécutent. Mais administrer pour quoi faire ? La réponse est simple : pour exploiter. La Société des Nations (SDN) a certes imposé à la France de veiller au « bien-être et au développement » des populations. En réalité, l'administration française organise d'abord la \cle{mise en valeur} du territoire, c'est-à-dire l'exploitation de ses richesses. Cette section étudie comment cette politique transforme profondément la vie économique et sociale des Camerounais entre 1922 et 1939.

\courssub{La mise en valeur : exploiter les ressources du territoire}

\begin{definition}[La mise en valeur]
La \cle{mise en valeur} est l'exploitation économique des ressources d'un territoire (forêts, terres agricoles, mines) organisée par la puissance coloniale au profit de la métropole et de l'administration coloniale. Les produits sont exportés vers l'Europe ; les profits reviennent essentiellement aux compagnies et aux colons européens.
\end{definition}

L'intuition est la suivante : le Cameroun possède des richesses (bois précieux, terres fertiles), mais celles-ci ne rapportent rien à la France tant qu'elles ne sont pas extraites, transportées et vendues. La mise en valeur repose donc sur deux piliers : la \textbf{production} (exploitations forestières et plantations) et les \textbf{infrastructures} (chemins de fer, routes, port) pour évacuer les produits.

\textbf{L'exploitation forestière.} La forêt équatoriale camerounaise fournit des bois précieux très recherchés en Europe, notamment l'\cle{ébène}, bois noir et dur utilisé pour l'ébénisterie de luxe. Des compagnies concessionnaires obtiennent de vastes territoires et emploient des villageois pour abattre et transporter les grumes jusqu'aux fleuves ou aux routes.

\textbf{Les plantations.} L'administration encourage la création de grandes plantations européennes : \cle{cacao} et \cle{café} dans le Centre et l'Ouest, \cle{bananes} dans le Littoral (la Compagnie fruitière s'installe près de Douala), \cle{hévéa} (arbre à caoutchouc) dans le Sud. Ces cultures d'exportation sont destinées au marché européen, non à la consommation locale. Les paysans camerounais sont aussi contraints de cultiver eux-mêmes le cacao ou le café pour payer leurs impôts.

\begin{popfigure}
\begin{center}
\begin{tikzpicture}[scale=1.0, every node/.style={font=\small}]
% Contour simplifié du Cameroun français
\draw[thick, popdark, fill=popcream] (0.5,0) -- (3.5,0) -- (4.2,1.2) -- (4.8,2.5) -- (4.4,4) -- (3.6,5.2) -- (2.6,5.6) -- (1.6,5.0) -- (1.2,3.8) -- (0.4,2.6) -- (0.2,1.2) -- cycle;
% Zones de plantations
\fill[popgreen!45] (0.7,0.3) -- (2.6,0.3) -- (2.9,1.4) -- (1.0,1.6) -- cycle;
\fill[popgold!50] (1.2,1.8) -- (3.0,1.6) -- (3.4,3.0) -- (1.5,3.2) -- cycle;
% Forêt / ébène
\fill[popgreen!20] (1.6,3.4) -- (3.3,3.2) -- (3.0,4.6) -- (1.8,4.4) -- cycle;
% Villes
\fill[popink] (0.9,0.9) circle (2pt) node[left=1pt, popink]{\textbf{Douala}};
\fill[popink] (2.6,2.3) circle (2pt) node[right=1pt, popink]{\textbf{Yaoundé}};
\fill[popink] (1.7,1.5) circle (1.5pt) node[below=1pt, popink]{Edéa};
% Transcamerounais
\draw[line width=2pt, poporange] (0.9,0.9) -- (1.7,1.5) -- (2.6,2.3);
\draw[line width=2pt, poporange, dashed] (2.6,2.3) -- (3.1,3.6);
% Port
\draw[popblue, thick] (0.9,0.9) -- (0.55,0.55);
\node[popblue] at (0.25,0.45) {Port};
% Nord
\draw[->, thick, popdark] (4.5,4.6) -- (4.5,5.4) node[above]{N};
% Légende interne
\node[anchor=west] at (0.6,5.15) {\textcolor{popgreen!70!black}{\rule{8pt}{8pt}} Plantations (bananes, cacao)};
\node[anchor=west] at (0.6,4.75) {\textcolor{popgold!80!black}{\rule{8pt}{8pt}} Cacao, café, hévéa};
\node[anchor=west] at (0.6,4.35) {\textcolor{popgreen!50!black}{\rule{8pt}{8pt}} Forêt (ébène)};
\end{tikzpicture}
\end{center}
\legende{Croquis des axes de mise en valeur du Cameroun français : le Transcamerounais (trait orange) relie le port de Douala à Yaoundé (achevé en 1927) ; les zones colorées indiquent les principales productions d'exportation. Croquis schématique, sans précision cartographique.}
\end{popfigure}

\courssub{Les grands travaux : chemin de fer, routes et port de Douala}

Pour évacuer les produits vers l'Europe, il faut des infrastructures. L'administration lance de \cle{grands travaux} :

\begin{itemize}
\item le \textbf{chemin de fer Douala--Yaoundé}, appelé \cle{Transcamerounais}, achevé en \textbf{1927} : il relie la capitale administrative au port et traverse des zones forestières difficiles ;
\item la construction de \textbf{routes} carrossables reliant les chefs-lieux de circonscription ;
\item l'aménagement du \textbf{port de Douala}, principal point d'exportation du cacao, du bois et des bananes.
\end{itemize}

\begin{exemplebox}[Le Transcamerounais : un chantier coûteux en vies humaines]
La construction du chemin de fer Douala--Yaoundé mobilise des \textbf{dizaines de milliers de travailleurs forcés}, recrutés par les chefs dans les villages, parfois très loin du chantier. Les conditions sont terribles : travail épuisant en forêt équatoriale, paludisme, accidents, nourriture insuffisante. Le chantier coûte de \textbf{nombreuses vies humaines} : on estime que plusieurs milliers de travailleurs y ont laissé leur vie. Cet exemple montre que les infrastructures coloniales, souvent présentées comme des « bienfaits », ont été payées par la sueur et le sang des Camerounais.
\end{exemplebox}

\courssub{Le travail forcé, la prestation et l'impôt de capitation}

D'où viennent ces travailleurs ? Pas d'un marché libre de l'emploi : l'administration les obtient par la contrainte.

\begin{definition}[Le travail forcé et la prestation]
Le \cle{travail forcé} est un travail imposé par l'administration coloniale aux populations, sans contrat libre. La \cle{prestation} est une forme de travail obligatoire : chaque adulte doit fournir un certain nombre de journées de travail par an (corvées de \textbf{portage}, de construction \textbf{routière} et \textbf{ferroviaire}). Il est mal payé ou non payé.
\end{definition}

\begin{attention}[Erreur fréquente : travail forcé n'est pas emploi salarié]
Ne confonds pas travail forcé et travail salarié. Le salarié \emph{choisit} son emploi, signe un contrat et reçoit un salaire. Le travailleur forcé est \emph{imposé} par l'administration, encadré par des gardes, mal ou non payé, et puni en cas de refus ou de fuite. Dans une copie, écrire « les Camerounais étaient employés à la construction du chemin de fer » est une erreur : il faut écrire qu'ils y étaient \textbf{contraints}.
\end{attention}

\begin{definition}[L'impôt de capitation]
L'\cle{impôt de capitation} est un impôt payé par chaque adulte (par « tête »), en \textbf{argent monnaie}. Comme les villageois pratiquent une économie de subsistance sans monnaie, cet impôt les oblige à gagner de l'argent : ils doivent vendre des récoltes, partir travailler dans les plantations ou dans les villes. C'est un outil qui force l'entrée des Camerounais dans l'\textbf{économie monétaire} et fournit de la main-d'œuvre aux colons.
\end{definition}

Le mécanisme est donc un cercle : l'impôt pousse les hommes vers les plantations et les chantiers ; les réquisitions de \textbf{vivres} et de \textbf{main-d'œuvre} s'ajoutent à la charge, surtout pendant la \textbf{Seconde Guerre mondiale} (1939-1945), où le Cameroun doit nourrir et soutenir l'effort de guerre de la France libre. Les villages, privés de leurs hommes, voient les femmes, les enfants et les vieillards assurer seuls les cultures vivrières, ce qui provoque des disettes.

\courssub{Les conditions sanitaires et les campagnes de médecine préventive}

La mise en valeur se heurte à un obstacle : les \textbf{maladies endémiques} qui déciment les populations et la main-d'œuvre. Le \cle{paludisme} (transmis par les moustiques) et la \cle{trypanosomiase} ou \textbf{maladie du sommeil} (transmise par la mouche tsé-tsé) tuent des milliers de personnes, surtout dans les régions forestières.

L'administration organise alors des campagnes de \cle{médecine préventive} : prospection médicale (tournées dans les villages pour dépister les malades), traitements, assainissement. C'est une politique réellement bénéfique, mais elle répond aussi à un calcul économique : une main-d'œuvre saine est plus productive.

\begin{savaistu}[Le docteur Jamot et la maladie du sommeil]
Le médecin français \textbf{Eugène Jamot} organisa dans les années 1920-1930 la lutte contre la maladie du sommeil dans la région d'\textbf{Ayos} (Centre-Cameroun). Ses équipes parcouraient les villages pour dépister et soigner systématiquement les malades. Grâce à cette méthode de prospection mobile, des milliers de vies furent sauvées et la maladie recula fortement. Le docteur Jamot reste une figure marquante de l'histoire sanitaire du Cameroun.
\end{savaistu}

\courssub{L'enseignement : former une petite élite}

L'administration et les missions chrétiennes (catholiques et protestantes) ouvrent des écoles. Le système est pyramidal et très sélectif :

\begin{itemize}
\item les \textbf{écoles de village}, souvent tenues par les missions, enseignent la lecture, l'écriture, le calcul et le catéchisme ;
\item les \textbf{écoles primaires régionales} accueillent les meilleurs élèves ;
\item l'\textbf{école primaire supérieure} forme le sommet de la pyramide : instituteurs, infirmiers, employés de bureau, interprètes.
\end{itemize}

L'objectif colonial n'est pas d'instruire le peuple entier, mais de former une \cle{petite élite} auxiliaire, utile à l'administration et aux entreprises. Pourtant, cette élite scolarisée jouera plus tard un rôle majeur : c'est elle qui rédigera les pétitions et animera les premiers mouvements de revendication.

\courssub{Résistances et pétitions : la voix des Camerounais à la SDN}

Face au travail forcé, à l'impôt et aux réquisitions, les Camerounais ne restent pas passifs. Les résistances prennent plusieurs formes : fuites des chantiers, refus de payer l'impôt, révoltes locales, mais aussi une arme nouvelle : la \cle{pétition}. Le statut de mandat oblige la France à rendre des comptes à la SDN : des Camerounais lettrés adressent donc des pétitions à la \textbf{Société des Nations} pour dénoncer les abus de l'administration.

Des \textbf{associations} se créent pour défendre les intérêts des populations. La plus célèbre est la \cle{Jeucafra} (\emph{Jeunesse camerounaise française}), fondée en \textbf{1938}, qui regroupe de jeunes évolués et revendique l'amélioration des conditions de vie et l'égalité des droits.

\begin{exoresolu}[Exercice résolu : expliquer un mécanisme économique colonial]
\textbf{Énoncé.} Explique en quelques lignes comment l'impôt de capitation fournit de la main-d'œuvre aux plantations européennes au Cameroun sous mandat français.
\tcblower
\textbf{Solution.} L'impôt de capitation doit être payé par chaque adulte \emph{en argent monnaie}. Or les paysans camerounais vivent d'une économie de subsistance (ils produisent leur nourriture) et ne possèdent presque pas de monnaie. Pour se procurer l'argent de l'impôt, les hommes sont donc obligés de vendre des récoltes (cacao, café) ou de partir travailler contre un salaire dans les plantations européennes, sur les chantiers ou dans les villes. L'impôt agit ainsi comme un mécanisme de contrainte : il force l'entrée dans l'économie monétaire et met une main-d'œuvre abondante et bon marché à la disposition des colons.
\end{exoresolu}

\begin{popfigure}
\begin{center}
\begin{tikzpicture}
\begin{axis}[
    ybar, bar width=14pt,
    width=0.85\linewidth, height=6cm,
    symbolic x coords={1922,1926,1930,1934,1938},
    xtick=data,
    ylabel={Exportations (milliers de tonnes, indicatif)},
    ymin=0, ymax=60,
    legend style={at={(0.02,0.97)}, anchor=north west, font=\small},
    nodes near coords, every node near coord/.append style={font=\tiny},
    axis lines=left,
]
\addplot[fill=popgold!80] coordinates {(1922,8) (1926,18) (1930,28) (1934,22) (1938,35)};
\addplot[fill=popgreen!70] coordinates {(1922,12) (1926,20) (1930,30) (1934,25) (1938,40)};
\addplot[fill=poporange!80] coordinates {(1922,2) (1926,8) (1930,15) (1934,12) (1938,25)};
\legend{Cacao, Bois, Bananes}
\end{axis}
\end{tikzpicture}
\end{center}
\legende{Évolution indicative des exportations du Cameroun français entre 1922 et 1938 (ordres de grandeur). On observe une croissance générale, freinée vers 1930-1934 par la crise économique mondiale, puis une reprise. Source : données reconstituées à visée pédagogique.}
\end{popfigure}

\begin{aretenir}[L'essentiel de la section]
\begin{itemize}
\item La \textbf{mise en valeur} organise l'exploitation du bois (ébène) et des plantations (cacao, café, bananes, hévéa) au profit de la métropole.
\item Les \textbf{grands travaux} (Transcamerounais achevé en 1927, routes, port de Douala) servent à exporter ces produits.
\item Cette économie repose sur la contrainte : \textbf{travail forcé}, \textbf{prestations}, \textbf{impôt de capitation}, \textbf{réquisitions}.
\item Des progrès réels existent dans la \textbf{santé} (docteur Jamot contre la maladie du sommeil) et l'\textbf{enseignement}, qui forme une petite élite.
\item Les Camerounais résistent : fuites, révoltes, \textbf{pétitions à la SDN}, associations comme la \textbf{Jeucafra} (1938).
\end{itemize}
\end{aretenir}

Cette vie économique et sociale marquée par la contrainte prépare les revendications de l'après-guerre. Reste à comparer cette situation avec celle de l'autre territoire camerounais : la prochaine section étudie la vie dans le Cameroun sous administration britannique, où l'\emph{indirect rule} produit des effets différents.

\coursec{La vie dans le Cameroun sous administration britannique}

Dans la section précédente, nous avons vu que le Cameroun sous mandat français connaît une administration centrale forte, dirigée depuis Yaoundé par un commissaire de la République, avec une politique économique active (grands travaux, plantations, mise en valeur) mais aussi des contraintes lourdes pour les populations (impôts, travail forcé, indigénat). À l'ouest, de l'autre côté de la ligne de partage, la situation est radicalement différente : la Grande-Bretagne administre ses deux étroites bandes de territoire camerounais non pas directement, mais à travers sa grande colonie voisine, le \cle{Nigeria}. Cette particularité — un territoire sans administration propre, gouverné « à distance » — explique l'essentiel de la vie quotidienne des Camerounais britanniques : une économie tournée vers le Nigeria, un système politique fondé sur les chefs traditionnels, un enseignement confié aux missions, et un profond sentiment d'abandon qui nourrira les premières revendications.

\courssub{Une administration à distance : le British Cameroons gouverné depuis le Nigeria}

\begin{definition}[British Cameroons]
Le \cle{British Cameroons} désigne la partie du Cameroun placée sous mandat britannique par la Société des Nations (SDN) en 1922. Il est constitué de deux bandes de territoire non contiguës le long de la frontière nigériane : le \cle{Northern Cameroons} au nord et le \cle{Southern Cameroons} au sud, séparés l'un de l'autre par une portion du Cameroun français.
\end{definition}

Contrairement à la France, qui crée à Yaoundé une administration complète et autonome, la Grande-Bretagne fait le choix de l'économie : elle ne veut pas dépenser pour un territoire étroit, peu peuplé et qu'elle juge secondaire. Le British Cameroons est donc \textbf{administré comme une simple partie du Nigeria} :

\begin{itemize}
\item le \cle{Southern Cameroons} est rattaché à la province de l'Est du Nigeria (région de l'Est après 1939), dont le chef-lieu est \textbf{Enugu} ;
\item le \cle{Northern Cameroons} est rattaché aux provinces du Nord du Nigeria (région du Nord après 1939), administrées depuis \textbf{Kaduna} ;
\item l'ensemble dépend du gouverneur du Nigeria, installé à \textbf{Lagos} (capitale fédérale).
\end{itemize}

Un \cle{Résident} britannique, installé à \textbf{Buea}, représente l'autorité coloniale dans le Southern Cameroons, mais il n'a pas de réelle autonomie : les décisions importantes (budget, grands travaux, législation) sont prises à Enugu ou à Lagos. Dans le Northern Cameroons, l'administration est assurée par des officiers de district rattachés à l'administration de la région du Nord nigériane.

\begin{popfigure}
\begin{center}
\begin{tikzpicture}[scale=1, every node/.style={font=\small}]
% Nigeria
\fill[popgreenL] (-0.2,-0.4) rectangle (4.2,5.6);
\draw[popdark, thick] (-0.2,-0.4) rectangle (4.2,5.6);
\node[popdark, font=\bfseries] at (2,5.2) {NIGERIA (britannique)};
% Cameroun français
\fill[popblueL] (5.4,-0.4) rectangle (9.2,5.6);
\draw[popdark, thick] (5.4,-0.4) rectangle (9.2,5.6);
\node[popdark, font=\bfseries] at (7.3,5.2) {CAMEROUN FRANÇAIS};
% Northern Cameroons
\fill[poporangeL] (4.2,3.1) rectangle (5.4,5.6);
\draw[poporange, thick] (4.2,3.1) rectangle (5.4,5.6);
\node[align=center, font=\footnotesize\bfseries, popdark] at (4.8,4.35) {Northern\\Cameroons};
% Southern Cameroons
\fill[poporangeL] (4.2,0.2) rectangle (5.4,3.1);
\draw[poporange, thick] (4.2,0.2) rectangle (5.4,3.1);
\node[align=center, font=\footnotesize\bfseries, popdark] at (4.8,1.65) {Southern\\Cameroons};
% Villes
\fill[popink] (1.2,0.4) circle (2pt); \node[left, font=\footnotesize] at (1.2,0.4) {Lagos};
\fill[popink] (3.4,1.6) circle (2pt); \node[left, font=\footnotesize] at (3.4,1.6) {Enugu};
\fill[popink] (3.6,3.4) circle (2pt); \node[left, font=\footnotesize] at (3.6,3.4) {Kaduna};
\fill[popink] (4.9,0.7) circle (2pt); \node[right, font=\footnotesize] at (4.9,0.7) {Buea};
\fill[popink] (4.85,2.5) circle (2pt); \node[right, font=\footnotesize] at (4.85,2.5) {Victoria};
% Liaisons administratives
\draw[-{Stealth[length=3mm]}, very thick, poppurple] (3.4,1.6) -- (4.85,1.4);
\draw[-{Stealth[length=3mm]}, very thick, poppurple] (1.2,0.4) -- (4.6,0.55);
\draw[-{Stealth[length=3mm]}, very thick, poppurple] (3.6,3.4) -- (4.7,4.3);
% Nord
\draw[-{Stealth[length=2.5mm]}, popdark] (9.0,5.0) -- (9.0,5.5);
\node[font=\footnotesize] at (9.0,4.85) {N};
\end{tikzpicture}
\end{center}
\legende{Carte schématique du British Cameroons (1922-1946). 1 : Northern Cameroons rattaché à la région du Nord du Nigeria (Kaduna). 2 : Southern Cameroons rattaché à la région de l'Est (Enugu). 3 : flèches violettes : liaisons administratives vers Enugu, Kaduna et Lagos. 4 : Buea, résidence du Résident britannique (Southern Cameroons). Croquis sans précision cartographique.}
\end{popfigure}

\begin{attention}[Erreur fréquente à éviter]
Ne jamais écrire que le British Cameroons possède une administration propre et autonome comme le Cameroun français. C'est faux : le territoire n'a \textbf{ni gouverneur, ni budget, ni assemblée législative propres}. Il est gouverné depuis le Nigeria (Enugu, Kaduna, Lagos), ce qui explique sa marginalisation. Au Probatoire, confondre les deux systèmes fait perdre des points.
\end{attention}

\courssub{L'indirect rule en pratique : Native Authorities, Native Courts, Native Treasuries}

Comment gouverner un territoire avec très peu de fonctionnaires britanniques ? La réponse de la Grande-Bretagne s'appelle l'\cle{indirect rule} (« administration indirecte »), théorisée par Lord Lugard au Nigeria : au lieu d'administrer directement les populations, on gouverne \textbf{à travers les chefs traditionnels}, que l'on reconnaît, contrôle et rémunère.

\begin{definition}[Native Authority]
Une \cle{Native Authority} (« autorité indigène ») est une autorité traditionnelle — chef, lamido, fon ou conseil de notables — reconnue officiellement par l'administration britannique et dotée de pouvoirs administratifs, fiscaux et judiciaires sur sa population. Elle perçoit l'impôt, fait appliquer les ordres coloniaux et rend la justice coutumière, sous le contrôle du Résident britannique.
\end{definition}

Le dispositif repose sur trois rouages complémentaires :

\begin{itemize}
\item les \cle{Native Authorities} : les chefs reconnus deviennent des auxiliaires de l'administration ; ils lèvent l'impôt et reçoivent en échange un pourcentage des recettes ;
\item les \cle{Native Courts} (tribunaux coutumiers) : ils jugent les litiges selon la coutume locale, sous la surveillance du Résident qui peut réviser les jugements ;
\item les \cle{Native Treasuries} (trésoreries indigènes) : elles gèrent les recettes fiscales locales et financent les petits services (marchés, chemins, salaires des chefs).
\end{itemize}

\begin{popfigure}
\begin{center}
\begin{tikzpicture}[
  box/.style={draw=popdark, thick, rounded corners=2pt, fill=popblueL, align=center, minimum width=3.6cm, minimum height=0.9cm, font=\small},
  boxo/.style={box, fill=poporangeL},
  boxg/.style={box, fill=popgreenL},
  arr/.style={-{Stealth[length=3mm]}, very thick, popdark}]
\node[boxo, align=center] (res) at (0,4.5){\textbf{Résident britannique}\\(Buea) — contrôle général};
\node[box, align=center] (na) at (0,3.0){\textbf{Native Authority}\\chef traditionnel reconnu};
\node[boxg, align=center] (nc) at (-3.2,1.4){\textbf{Native Court}\\justice coutumière};
\node[boxg, align=center] (nt) at (3.2,1.4){\textbf{Native Treasury}\\impôts et budget local};
\node[box, align=center] (vil) at (0,0){\textbf{Village}\\population, impôts, corvées};
\draw[arr] (res) -- (na);
\draw[arr] (na) -- (nc);
\draw[arr] (na) -- (nt);
\draw[arr] (nc) |- (vil);
\draw[arr] (nt) |- (vil);
\draw[arr] (na) -- (vil);
\end{tikzpicture}
\end{center}
\legende{Schéma fonctionnel de l'indirect rule dans le Southern Cameroons. 1 : le Résident contrôle la Native Authority. 2 : la Native Authority dirige la Native Court (justice) et la Native Treasury (finances). 3 : les ordres, l'impôt et la justice coutumière descendent jusqu'au village.}
\end{popfigure}

Ce système présente un avantage pour les Britanniques : il coûte peu cher et s'appuie sur des autorités que la population connaît. Mais il a des limites : certains chefs, transformés en agents coloniaux, perdent leur légitimité ; dans les régions sans chefs centralisés (sociétés segmentaires de la Grassfield ou de la forêt), les Britanniques créent des « \cle{chefs de warrant} » (warrant chiefs) artificiels, ce qui provoque des conflits et de la résistance.

\courssub{Une économie dominée par le Nigeria}

L'économie du British Cameroons est entièrement tournée vers le Nigeria, qui absorbe ses produits et lui fournit ses importations.

\textbf{Les plantations.} Les riches plantations allemandes de bananes et de palmiers à huile des pentes du mont Cameroun, confisquées pendant la guerre, sont d'abord gérées par l'administration (Custodian of Enemy Property), puis revendues — souvent à leurs anciens propriétaires allemands dans les années 1920. La production d'\textbf{huile de palme}, de \textbf{bananes} et de \textbf{cacao} est exportée vers le Nigeria puis vers l'Europe.

\begin{exemplebox}[La Cameroon Development Corporation (CDC)]
En \textbf{1947}, l'administration britannique regroupe les anciennes plantations allemandes au sein de la \cle{Cameroon Development Corporation} (CDC), entreprise publique créée par l'ordonnance n° 17 de 1947, qui cultive le palmier à huile, la banane, l'hévéa et le thé. La CDC devient \textbf{l'un des plus grands employeurs agricoles d'Afrique de l'Ouest} : elle emploie des dizaines de milliers de travailleurs, construit des camps de logement, des dispensaires et des écoles autour de ses plantations (Tiko, Bota, Ekona, Mbonge). Elle structure encore aujourd'hui l'économie du Sud-Ouest camerounais.
\end{exemplebox}

\textbf{La fiscalité et le travail.} Comme dans le Cameroun français, l'impôt (la \emph{poll tax}, impôt par tête) est la principale ressource du territoire. Pour le payer, les paysans ont besoin de monnaie : beaucoup partent travailler quelques mois dans les plantations côtières ou migrent vers le Nigeria, où les salaires sont plus élevés. Cette \cle{migration de main-d'œuvre} vide les villages et renforce la dépendance économique envers le voisin nigérian.

\courssub{L'enseignement confessionnel : le rôle majeur des missions}

L'administration britannique investit très peu dans l'éducation : elle laisse ce domaine aux \cle{sociétés missionnaires}, conformément à sa politique d'économie. Deux grandes missions dominent :

\begin{itemize}
\item la \cle{Basel Mission} (mission protestante de Bâle), présente depuis l'époque allemande, expulsée pendant la guerre, puis remplacée par l'Église presbytérienne (Presbyterian Church), qui ouvre des écoles primaires et des écoles de formation dans la Grassfield et sur la côte ;
\item les missions \cle{catholiques} (notamment les Pères du Sacré-Cœur de Saint-Quentin, arrivés en 1922), qui développent un réseau scolaire autour de Buea et de Victoria.
\end{itemize}

L'enseignement se fait en \textbf{anglais}, langue qui devient celle de l'administration, de l'école et des élites. C'est de ce fait scolaire et linguistique que naît une identité « anglophone » distincte de celle des Camerounais francophones.

\begin{savaistu}[L'origine lointaine de la question anglophone]
C'est dans le Southern Cameroons que naîtra plus tard la \cle{question anglophone}, héritée directement de cette double administration : quarante ans de gouvernement britannique, d'école en anglais et de rattachement au Nigeria ont forgé une identité politique et culturelle différente de celle du Cameroun français. Les débats actuels sur l'anglophonie au Cameroun plongent ainsi leurs racines dans la période du mandat.
\end{savaistu}

\courssub{La marginalisation du territoire et les premières revendications}

Le bilan de l'administration britannique est marqué par le \cle{sous-investissement} : presque pas de routes modernes, peu d'hôpitaux, un réseau scolaire laissé aux missions, un budget dépendant des décisions de Lagos. Les populations ont le sentiment d'être les « oubliées » du mandat : le territoire est une simple annexe du Nigeria.

Ce sentiment d'abandon nourrit les \textbf{premières revendications} :

\begin{itemize}
\item une petite élite anglophone, formée dans les écoles de mission (instituteurs, catéchistes, employés de la CDC, commerçants), prend conscience des retards du territoire ;
\item des \cle{unions tribales} et associations d'originaires (unions d'amélioration) se créent pour défendre les intérêts des villages et financer des écoles ou des bourses ;
\item des \textbf{pétitions} sont adressées à la Société des Nations, qui reçoit chaque année le rapport de l'administration mandataire : les signataires dénoncent l'absence d'investissements et demandent que le territoire soit administré pour lui-même, et non comme une annexe du Nigeria.
\end{itemize}

\begin{methode}[Comparer deux documents historiques]
Pour comparer, par exemple, un rapport britannique sur le Southern Cameroons et une pétition d'élites à la SDN :
\begin{enumerate}
\item \textbf{Identifier la nature} de chaque document : rapport officiel ? pétition ? lettre ? discours ?
\item \textbf{Identifier l'auteur} : administrateur colonial, chef traditionnel, élite scolarisée ? Quel est son point de vue ?
\item \textbf{Replacer le contexte} : date, lieu, situation du territoire à ce moment.
\item \textbf{Dégager les convergences} : sur quels faits les deux documents s'accordent-ils (retard économique, dépendance au Nigeria) ?
\item \textbf{Dégager les différences} : divergences d'appréciation, d'intérêts, de solutions proposées.
\item \textbf{Conclure} en expliquant ce que la comparaison révèle de la situation historique.
\end{enumerate}
\end{methode}

\begin{exoresolu}[Un administrateur et un planteur face à l'indirect rule]
\textbf{Énoncé.} Document 1 (rapport d'un officier britannique, 1930) : « Le système des Native Authorities nous permet d'administrer le territoire à peu de frais, en nous appuyant sur les chefs que la population reconnaît. » Document 2 (mémoire d'une union tribale à la SDN, 1938) : « Notre territoire, traité comme une simple province du Nigeria, ne reçoit ni routes ni écoles ; nous demandons à être administrés pour nous-mêmes. » Comparez ces deux documents et dégagez ce qu'ils révèlent de l'administration britannique au Cameroun.
\tcblower
\textbf{Solution détaillée.}
\begin{enumerate}
\item \emph{Nature et auteurs} : le document 1 est un rapport officiel rédigé par un administrateur colonial ; le document 2 est une pétition d'élites camerounaises adressée à la SDN.
\item \emph{Contexte} : les deux textes relèvent de la période du mandat britannique, marquée par l'indirect rule et le rattachement au Nigeria.
\item \emph{Convergences} : les deux documents reconnaissent implicitement que le territoire est administré à moindre coût et dépend du Nigeria.
\item \emph{Différences} : l'administrateur valorise l'économie de moyens et la stabilité ; les élites dénoncent le sous-investissement et l'absence d'autonomie. Le point de vue colonial s'oppose au point de vue des administrés.
\item \emph{Conclusion} : la comparaison révèle la logique d'économie de l'administration britannique et le sentiment d'abandon des populations, source des premières revendications anglophones.
\end{enumerate}
\end{exoresolu}

\begin{aretenir}[L'essentiel de la section]
Le British Cameroons (Northern et Southern Cameroons) n'a pas d'administration autonome : il est gouverné depuis le Nigeria (Enugu, Kaduna, Lagos). Les Britanniques y appliquent l'\cle{indirect rule} : Native Authorities, Native Courts, Native Treasuries. L'économie, dominée par le Nigeria, repose sur les plantations (huile de palme, bananes ; CDC créée en 1947) et la migration de main-d'œuvre. L'enseignement en anglais est assuré par les missions (Basel Mission/Presbytériens, catholiques). Le sous-investissement provoque un sentiment d'abandon, à l'origine des premières revendications des élites anglophones et des pétitions à la SDN. Cette période (1922-1946) pose les bases de la future question anglophone.
\end{aretenir}

\begin{exercice}[Pour s'entraîner]
\begin{enumerate}
\item Expliquez en quelques lignes pourquoi on dit que le British Cameroons est administré « à distance ».
\item Définissez la Native Authority et citez les deux institutions qui l'accompagnent dans le système de l'indirect rule.
\item Montrez, à l'aide de deux exemples, que l'économie du Southern Cameroons dépend du Nigeria.
\item Rédigez un paragraphe argumenté (10 à 15 lignes) répondant à la question : « L'administration britannique au Cameroun : une politique d'économie aux conséquences durables. »
\end{enumerate}
\end{exercice}

\begin{corrige}[Exercice « Pour s'entraîner »]
\textbf{1.} Le territoire n'a ni gouverneur ni budget propres : le Southern Cameroons dépend de la région de l'Est du Nigeria (Enugu) et le Northern Cameroons de la région du Nord (Kaduna) ; les décisions majeures viennent de Lagos. \textbf{2.} La Native Authority est une autorité traditionnelle reconnue par les Britanniques, dotée de pouvoirs administratifs, fiscaux et judiciaires ; elle est accompagnée des Native Courts (tribunaux coutumiers) et des Native Treasuries (trésoreries indigènes). \textbf{3.} Les produits des plantations (huile de palme, bananes) sont exportés via le Nigeria ; la main-d'œuvre migre vers le Nigeria où les salaires sont plus élevés. \textbf{4.} Le paragraphe doit articuler : choix britannique de l'économie (indirect rule, rattachement au Nigeria, éducation laissée aux missions), conséquences (sous-investissement, dépendance économique, identité anglophone) et conclusion sur la durabilité de cet héritage (question anglophone, difficile intégration post-réunification).
\end{corrige}

\coursec{Bilan et transition : du mandat à la tutelle (1939-1946)}

Les sections précédentes ont montré comment, dans les deux zones du Cameroun, les populations vivaient au quotidien sous des administrations aux logiques distinctes : centralisation et assimilation à la française d'un côté, \emph{indirect rule} et préservation des chefferies traditionnelles de l'autre. Pourtant, au-delà de ces différences, les deux territoires partageaient une condition commune : celle de territoires sous mandat de la Société des Nations (SDN), administrés par des puissances mandataires. Cette période s'achève avec la Seconde Guerre mondiale, pour s'ouvrir sur un nouveau régime international : la tutelle des Nations unies.

\courssub{Bilan comparatif des deux administrations mandataires}

Avant d'aborder la rupture de 1939-1946, il convient de dresser un bilan structuré des années de mandat (1922-1939). Ce bilan permet de comprendre ce que la tutelle va hériter, transformer ou prolonger.

\begin{popfigure}
\begin{tikzpicture}[scale=0.9, every node/.style={font=\small}]
% En-tête
\fill[popblue!20] (0,0) rectangle (18,0.8);
\node[anchor=west] at (0.3,0.4) {\textbf{Thème}};
\node[anchor=center] at (6.5,0.4) {\textbf{Cameroun sous mandat français}};
\node[anchor=center] at (13.75,0.4) {\textbf{Cameroun sous mandat britannique}};

% Cadre et lignes
\draw[popline] (0,0) -- (18,0);
\draw[popline] (0,0.8) -- (18,0.8);
\draw[popline] (0,0) -- (0,8.8);
\draw[popline] (3.5,0) -- (3.5,8.8);
\draw[popline] (9.5,0) -- (9.5,8.8);
\draw[popline] (18,0) -- (18,8.8);

\foreach \y in {1.6,2.4,3.2,4.0,4.8,5.6,6.4,7.2,8.0,8.8} {
  \draw[popline] (0,\y) -- (18,\y);
}

% Contenu
\node[anchor=west] at (0.3,1.2) {\textbf{Cadre juridique}};
\node[anchor=center] at (6.5,1.2) {Mandat de classe B (SDN, 1922)};
\node[anchor=center] at (13.75,1.2) {Mandat de classe B (SDN, 1922)};

\node[anchor=west] at (0.3,2.0) {\textbf{Administration}};
\node[anchor=center] at (6.5,2.0) {Administration directe, centralisée};
\node[anchor=center] at (13.75,2.0) {Indirect rule, décentralisée};

\node[anchor=west] at (0.3,2.8) {\textbf{Langue officielle}};
\node[anchor=center] at (6.5,2.8) {Français};
\node[anchor=center] at (13.75,2.8) {Anglais};

\node[anchor=west] at (0.3,3.6) {\textbf{Droit applicable}};
\node[anchor=center] at (6.5,3.6) {Droit français + coutumier (subordonné)};
\node[anchor=center] at (13.75,3.6) {Common law + droit coutumier reconnu};

\node[anchor=west] at (0.3,4.4) {\textbf{Éducation}};
\node[anchor=center] at (6.5,4.4) {Écoles officielles et missionnaires, programme français};
\node[anchor=center] at (13.75,4.4) {Écoles confessionnelles, programme britannique};

\node[anchor=west] at (0.3,5.2) {\textbf{Économie}};
\node[anchor=center] at (6.5,5.2) {Plantations, travail forcé, port de Douala};
\node[anchor=center] at (13.75,5.2) {Petites exploitations, intégration au Nigeria};

\node[anchor=west] at (0.3,6.0) {\textbf{Fiscalité}};
\node[anchor=center] at (6.5,6.0) {Impôt de capitation, prestations};
\node[anchor=center] at (13.75,6.0) {Impôt direct, collecte par les chefferies};

\node[anchor=west] at (0.3,6.8) {\textbf{Représentation}};
\node[anchor=center] at (6.5,6.8) {Aucune (sujets, non citoyens)};
\node[anchor=center] at (13.75,6.8) {Native Authorities (chefferies reconnues)};

\node[anchor=west] at (0.3,7.6) {\textbf{Points communs}};
\node[anchor=center] at (6.5,7.6) {Exploitation, ségrégation, main-d'œuvre contrainte};
\node[anchor=center] at (13.75,7.6) {Exploitation, ségrégation, main-d'œuvre contrainte};

\node[anchor=west] at (0.3,8.4) {\textbf{Héritage après 1946}};
\node[anchor=center] at (6.5,8.4) {État centralisé, francophonie, droit civil};
\node[anchor=center] at (13.75,8.4) {Tradition fédérale, common law, anglophonie};
\end{tikzpicture}
\legende{Tableau de synthèse comparatif : mandat français et mandat britannique au Cameroun (1922-1939).}
\end{popfigure}

\begin{aretenir}[Points communs et différences structurelles]
\textbf{Points communs (héritage commun du mandat) :}
\begin{itemize}
\item Exploitation économique au profit de la métropole (plantations, main-d'œuvre, matières premières).
\item Fiscalité lourde (impôt de capitation, prestations en nature ou en travail).
\item Absence de citoyenneté politique : les Camerounais sont des « sujets », pas des citoyens.
\item Ségrégation structurelle : résidentielle, judiciaire, salariale.
\item Obligation théorique de rendre des comptes à la SDN (rapports annuels), contrôle en pratique très faible.
\end{itemize}

\textbf{Différences majeures (héritage dual) :}
\begin{itemize}
\item \textbf{Centralisation contre décentralisation :} administration directe française (circonscriptions, chefs de canton nommés) contre \emph{indirect rule} britannique (Native Authorities, chefferies traditionnelles reconnues).
\item \textbf{Langue et droit :} français et droit écrit de type civil contre anglais et common law.
\item \textbf{Éducation :} écoles officielles et missionnaires francophones contre écoles confessionnelles anglophones.
\item \textbf{Intégration territoriale :} le Cameroun français est administré comme une entité distincte (capitales Douala puis Yaoundé) ; le Cameroun britannique est rattaché administrativement au Nigeria voisin.
\end{itemize}
Ces dualismes (juridique, scolaire, linguistique, administratif) structurent encore le Cameroun contemporain.
\end{aretenir}

\courssub{Le Cameroun dans la Seconde Guerre mondiale : le ralliement à la France libre}

L'entrée en guerre de la France contre l'Allemagne (3 septembre 1939) place le Cameroun sous mandat français en situation de territoire concerné par le conflit. Après la défaite française et l'armistice du 22 juin 1940, l'administration coloniale locale reste d'abord fidèle au régime de Vichy. Mais la donne change radicalement en août 1940.

\begin{exemplebox}[Le ralliement du Cameroun à la France libre (août 1940)]
\begin{itemize}
\item \textbf{26 août 1940 :} le Tchad, sous l'impulsion du gouverneur Félix Éboué, se rallie au général de Gaulle : c'est le premier territoire d'Afrique équatoriale française (AEF) à le faire.
\item \textbf{27 août 1940 :} le colonel Philippe Leclerc de Hauteclocque, envoyé par de Gaulle, arrive à Douala avec un petit détachement. Sous la pression des partisans de la France libre, l'autorité coloniale cède : le Cameroun se rallie à la France libre.
\item \textbf{28-29 août 1940 :} le ralliement s'étend à l'ensemble du territoire ; le Congo (Brazzaville) rallie à son tour. Le Cameroun devient l'un des \textbf{premiers territoires africains} engagés aux côtés de la France libre.
\item \textbf{Conséquence immédiate :} le Cameroun devient une base arrière stratégique pour les Forces françaises libres (FFL) : recrutement de soldats et de porteurs, production de caoutchouc, de bois et de vivres pour l'effort de guerre.
\end{itemize}
\end{exemplebox}

\begin{exoresolu}[Analyse d'un document : l'appel à la mobilisation de 1940]
\textbf{Document :} Extrait d'un message diffusé à Douala lors du ralliement, août 1940 : « Camerounais, la France a besoin de vous. Le général de Gaulle, chef des Français libres, vous appelle à la lutte pour la libération de la patrie. »

\textbf{Question :} En quoi ce document illustre-t-il la rupture politique de 1940 et ses ambiguïtés ?

\textbf{Solution détaillée :}
\begin{enumerate}
\item \textbf{Contexte :} le document émane des partisans du général de Gaulle venus rallier le Cameroun. Il s'adresse aux « Camerounais », terme qui rompt avec le vocabulaire administratif colonial (« indigènes », « sujets »), marquant une rhétorique nouvelle d'inclusion.
\item \textbf{Rupture :} le ralliement brise la continuité avec le régime de Vichy. Le territoire bascule du côté des Alliés, ce qui a des conséquences militaires (base de départ des campagnes d'Afrique du Nord) et symboliques (l'un des premiers territoires africains ralliés).
\item \textbf{Ambiguïté :} l'appel à la « libération de la patrie » désigne la France, non le Cameroun. Les soldats camerounais combattent pour la France, mais reviennent sans droits politiques accrus. La « patrie » reste la métropole.
\item \textbf{Conclusion :} le document montre l'instrumentalisation du patriotisme français pour mobiliser les populations colonisées, tout en semant les germes d'une conscience politique nouvelle chez les anciens combattants.
\end{enumerate}
\end{exoresolu}

\begin{savaistu}[Le 27 août 1940, date charnière]
Le \textbf{27 août 1940}, le Cameroun se rallie au général de Gaulle, devenant l'un des \textbf{tout premiers territoires d'Afrique subsaharienne} à rejoindre la France libre, juste après le Tchad (26 août). Plusieurs milliers de Camerounais seront ensuite recrutés comme soldats ou comme porteurs dans les Forces françaises libres, participant aux campagnes d'Afrique du Nord, d'Italie et de France (débarquement de Provence, août 1944).
\end{savaistu}

\courssub{Les réquisitions de guerre et leurs effets sur les populations}

L'effort de guerre repose sur une intensification brutale de l'exploitation coloniale. L'administration de la France libre met en place une économie de guerre dirigée.

\begin{propriete}[Mécanismes de l'économie de guerre au Cameroun français (1940-1945)]
\begin{itemize}
\item \textbf{Travail forcé accentué :} réquisition de main-d'œuvre pour les plantations (caoutchouc, huile de palme, cacao), les chantiers routiers et les ports.
\item \textbf{Réquisitions de produits vivriers :} riz, manioc, bananes, maïs prélevés pour nourrir les troupes et les villes (Douala, Yaoundé), à des prix imposés par l'administration.
\item \textbf{Contrôle des prix et monopoles :} l'administration fixe des prix d'achat très bas aux producteurs et revend à prix fort ; les bénéfices financent l'effort de guerre.
\item \textbf{Recrutement militaire :} engagement présenté comme volontaire mais fortement incité (pressions exercées par l'intermédiaire des chefs de canton).
\end{itemize}
\end{propriete}

\begin{attention}[Effets sociaux dramatiques]
\begin{itemize}
\item \textbf{Disettes et famines locales :} les réquisitions vivrières privent les villages de leurs réserves ; des épisodes de pénurie grave sont signalés en 1942-1943.
\item \textbf{Fuite des populations :} de nombreux villageois quittent les zones de recrutement forcé pour se réfugier en brousse ou de l'autre côté de la frontière, au Cameroun britannique.
\item \textbf{Mortalité accrue :} conditions sanitaires dégradées, paludisme, malnutrition ; les porteurs (non combattants) paient un tribut très lourd (maladies, accidents, épuisement).
\item \textbf{Prise de conscience politique :} les anciens combattants reviennent avec l'expérience de l'égalité au front (mêmes risques que les soldats métropolitains) et la conviction que la domination coloniale n'est pas légitime.
\end{itemize}
\end{attention}

\courssub{De la SDN à l'ONU : un nouveau cadre international}

La disparition de la Société des Nations (SDN), incapable d'empêcher la guerre, et la création de l'Organisation des Nations unies (ONU) en 1945 (Charte de San Francisco, 26 juin 1945) modifient radicalement le statut juridique des territoires sous mandat.

\begin{definition}[La tutelle : définition juridique]
La \textbf{tutelle} est un régime international institué par la \textbf{Charte des Nations unies (Chapitre XII, articles 75 à 85)} en 1945, succédant au système des mandats de la SDN. Elle s'applique aux territoires qui ne sont pas encore autonomes. Son \textbf{objectif déclaré} est de conduire progressivement ces territoires vers l'\textbf{autonomie} ou l'\textbf{indépendance}, sous le contrôle de l'Assemblée générale et du Conseil de tutelle de l'ONU. La puissance administrante (autorité administrante) rend compte annuellement à l'ONU et accepte des missions de visite ainsi que les pétitions des habitants.
\end{definition}

\begin{attention}[Erreur fréquente : tutelle $\neq$ indépendance]
\begin{itemize}
\item La \textbf{tutelle n'est pas l'indépendance}. C'est une \textbf{étape intermédiaire} sous contrôle international.
\item La puissance administrante (France ou Royaume-Uni) conserve l'administration effective (police, justice, finances, armée).
\item L'ONU exerce une \textbf{surveillance} (rapports, pétitions, missions) mais ne gère pas le territoire.
\item L'indépendance ne sera acquise qu'ultérieurement (1\textsuperscript{er} janvier 1960 pour le Cameroun français ; 1961 pour les Cameroons britanniques).
\end{itemize}
\end{attention}

\begin{popfigure}
\begin{tikzpicture}[scale=0.95, every node/.style={font=\small}]
% Frise chronologique 1939-1949
\draw[popblue, line width=2pt] (0,0) -- (16,0);
\foreach \x/\annee in {0/1939, 1.6/1940, 3.2/1941, 4.8/1942, 6.4/1943, 8/1944, 9.6/1945, 11.2/1946, 12.8/1947, 14.4/1948, 16/1949} {
  \draw[popblue, line width=1.5pt] (\x,-0.15) -- (\x,0.15);
  \node[anchor=north, popblue] at (\x,-0.35) {\annee};
}

% Événements au-dessus
\node[anchor=south, popdark, align=center] at (0,0.5) {\textbf{3 sept. 1939}\\Entrée en guerre};
\node[anchor=south, poporange, align=center] at (1.6,1.3) {\textbf{22 juin 1940}\\Armistice (Vichy)};
\node[anchor=south, popgreen, align=center] at (3.2,0.5) {\textbf{27 août 1940}\\Ralliement du Cameroun\\à la France libre};
\node[anchor=south, poppurple, align=center] at (4.8,1.3) {\textbf{1942-1943}\\Réquisitions maximales,\\disettes, recrutement forcé};
\node[anchor=south, poppink, align=center] at (8,0.5) {\textbf{8 mai 1945}\\Capitulation de l'Allemagne};
\node[anchor=south, poppink, align=center] at (9.6,1.3) {\textbf{26 juin 1945}\\Charte de San Francisco\\(création de l'ONU)};
\node[anchor=south, popgold, align=center] at (11.2,0.5) {\textbf{13 déc. 1946}\\Accords de tutelle ONU\\(France et Royaume-Uni)};
\node[anchor=south, popblue, align=center] at (12.8,1.3) {\textbf{1947}\\Premiers rapports de tutelle,\\pétitions à l'ONU};

% Événements en-dessous
\node[anchor=north, popmuted, align=center] at (1.6,-0.6) {\textbf{Cameroun britannique :}\\administré depuis le Nigeria};
\node[anchor=north, popmuted, align=center] at (6.4,-0.6) {\textbf{1944 : Conférence de Brazzaville}\\(réformes promises par la France libre)};
\node[anchor=north, popmuted, align=center] at (9.6,-1.6) {\textbf{1946 : abolition du travail forcé}\\(loi du 11 avril 1946)};
\node[anchor=north, popmuted, align=center] at (14.4,-0.6) {\textbf{1948-1949 : création}\\de partis politiques\\(UPC au Cameroun français)};

% Flèches de repère
\draw[popline, ->, thick] (3.2,0.15) -- (3.2,0.45);
\draw[popline, ->, thick] (11.2,0.15) -- (11.2,0.45);
\end{tikzpicture}
\legende{Frise chronologique 1939-1949 : du ralliement à la France libre à l'accord de tutelle de l'ONU.}
\end{popfigure}

\courssub{L'accord de tutelle du 13 décembre 1946}

Le 13 décembre 1946, l'Assemblée générale de l'ONU approuve les \textbf{accords de tutelle} pour le Cameroun, conclus avec la France et le Royaume-Uni. Ces accords entrent en vigueur en 1947.

\begin{propriete}[Contenu des accords de tutelle (13 décembre 1946)]
\begin{itemize}
\item \textbf{Statut :} le Cameroun français et les Cameroons britanniques deviennent des \textbf{territoires sous tutelle} de l'ONU (Chapitre XII de la Charte).
\item \textbf{Autorités administrantes :} la France pour le Cameroun oriental (ex-Cameroun français), le Royaume-Uni pour les Northern et Southern Cameroons (toujours administrés depuis le Nigeria).
\item \textbf{Obligations des puissances administrantes :}
  \begin{itemize}
  \item Promouvoir l'avancement politique, économique, social et éducatif des populations.
  \item Développer les institutions politiques locales (assemblées territoriales, suffrage).
  \item Garantir les libertés fondamentales (presse, association, réunion).
  \item Transmettre un \textbf{rapport annuel} à l'ONU.
  \item Accepter les \textbf{pétitions} des habitants et les \textbf{missions de visite} du Conseil de tutelle.
  \end{itemize}
\item \textbf{Objectif final :} l'autonomie ou l'indépendance, sans calendrier imposé.
\end{itemize}
\end{propriete}

\courssub{Continuités et changements entre mandat et tutelle}

Le passage du mandat à la tutelle n'est pas une rupture totale : c'est une \textbf{transition} qui conserve des structures lourdes tout en ouvrant des espaces nouveaux.

\begin{popfigure}
\begin{tikzpicture}[scale=0.9, every node/.style={font=\small}]
% Tableau Mandat vs Tutelle
\fill[poppurple!15] (0,0) rectangle (17,0.7);
\node[anchor=west] at (0.3,0.35) {\textbf{Dimension}};
\node[anchor=center] at (7.75,0.35) {\textbf{Mandat (SDN, 1922-1946)}};
\node[anchor=center] at (13.75,0.35) {\textbf{Tutelle (ONU, 1946-1960/61)}};

\draw[popline] (0,0) -- (17,0);
\draw[popline] (0,0.7) -- (17,0.7);
\draw[popline] (0,0) -- (0,7.7);
\draw[popline] (5,0) -- (5,7.7);
\draw[popline] (10.5,0) -- (10.5,7.7);
\draw[popline] (17,0) -- (17,7.7);

\foreach \y in {1.4,2.1,2.8,3.5,4.2,4.9,5.6,6.3,7.0,7.7} {
  \draw[popline] (0,\y) -- (17,\y);
}

\node[anchor=west] at (0.3,1.05) {\textbf{Organe de contrôle}};
\node[anchor=center] at (7.75,1.05) {Commission des mandats (SDN) : faible, sans visite};
\node[anchor=center] at (13.75,1.05) {Conseil de tutelle (ONU) : missions, pétitions, rapports publics};

\node[anchor=west] at (0.3,1.75) {\textbf{Objectif officiel}};
\node[anchor=center] at (7.75,1.75) {« Bien-être et développement » (formule vague)};
\node[anchor=center] at (13.75,1.75) {« Autonomie ou indépendance » (explicite, art. 76)};

\node[anchor=west] at (0.3,2.45) {\textbf{Reddition de comptes}};
\node[anchor=center] at (7.75,2.45) {Rapport annuel à la SDN, contrôle limité};
\node[anchor=center] at (13.75,2.45) {Rapport public à l'ONU + examen contradictoire};

\node[anchor=west] at (0.3,3.15) {\textbf{Participation politique}};
\node[anchor=center] at (7.75,3.15) {Nulle (sujets, non citoyens)};
\node[anchor=center] at (13.75,3.15) {Assemblées territoriales (1946), suffrage élargi (1956)};

\node[anchor=west] at (0.3,3.85) {\textbf{Libertés publiques}};
\node[anchor=center] at (7.75,3.85) {Répression (indigénat, partis interdits)};
\node[anchor=center] at (13.75,3.85) {Libertés syndicales (1946), presse, partis autorisés};

\node[anchor=west] at (0.3,4.55) {\textbf{Économie}};
\node[anchor=center] at (7.75,4.55) {Exploitation coloniale, monopoles, travail forcé};
\node[anchor=center] at (13.75,4.55) {Plans de développement (FIDES, 1946), travail forcé aboli};

\node[anchor=west] at (0.3,5.25) {\textbf{Droit du travail}};
\node[anchor=center] at (7.75,5.25) {Régime répressif de la main-d'œuvre};
\node[anchor=center] at (13.75,5.25) {Code du travail outre-mer (1952), syndicats reconnus};

\node[anchor=west] at (0.3,5.95) {\textbf{Éducation}};
\node[anchor=center] at (7.75,5.95) {Faible scolarisation, écoles missionnaires};
\node[anchor=center] at (13.75,5.95) {Expansion scolaire, création de lycées, bourses};

\node[anchor=west] at (0.3,6.65) {\textbf{Nationalisme}};
\node[anchor=center] at (7.75,6.65) {Réprimé, expression politique muselée};
\node[anchor=center] at (13.75,6.65) {Espace légal pour les partis (UPC, 1948 ; partis anglophones)};

\node[anchor=west] at (0.3,7.35) {\textbf{Frontières}};
\node[anchor=center] at (7.75,7.35) {Frontière issue du partage de 1919};
\node[anchor=center] at (13.75,7.35) {Frontières confirmées, devenues étatiques (1960-1961)};
\end{tikzpicture}
\legende{Tableau de synthèse : ruptures et continuités entre le régime de mandat (SDN) et le régime de tutelle (ONU) au Cameroun.}
\end{popfigure}

\courssub{Héritages durables de la période de mandat (1919-1946)}

La période de mandat, bien que close juridiquement en 1946, laisse des empreintes structurelles qui fondent l'État camerounais contemporain.

\begin{savaistu}[Le bilinguisme et le dualisme juridico-scolaire : héritages directs du mandat]
\begin{itemize}
\item \textbf{Bilinguisme officiel (français-anglais) :} héritage direct du partage franco-britannique de 1919. La Constitution de 1996 consacre l'égalité des deux langues officielles.
\item \textbf{Dualisme juridique :} coexistence du \textbf{droit civil d'inspiration française} (zone francophone) et de la \textbf{common law d'inspiration britannique} (zone anglophone) : le système judiciaire camerounais est bijuridique.
\item \textbf{Dualisme scolaire :} deux sous-systèmes éducatifs (francophone et anglophone, avec le GCE Ordinary et Advanced Level dans la zone anglophone). Les réformes successives tentent une harmonisation, mais les deux logiques persistent.
\item \textbf{Frontières actuelles :} la ligne de partage de 1919 est devenue frontière entre les deux entités de 1961, puis limite entre les régions du Nord-Ouest et du Sud-Ouest (anglophones) et les huit régions francophones.
\item \textbf{Administration territoriale :} le découpage actuel reprend largement les circonscriptions héritées du mandat français et les divisions du mandat britannique.
\item \textbf{Chefferies traditionnelles :} reconnues et institutionnalisées notamment par l'\emph{indirect rule} britannique, elles sont intégrées après l'indépendance comme auxiliaires de l'administration.
\end{itemize}
\end{savaistu}

\begin{methode}[Rédiger une conclusion de composition d'Histoire (Probatoire technique)]
\textbf{Objectif :} synthétiser la démonstration et ouvrir sur la période suivante (la tutelle, puis l'indépendance).
\begin{enumerate}
\item \textbf{Rappeler la problématique} (reformulée) : « En quoi la période de mandat (1919-1946) a-t-elle structuré le Cameroun contemporain, entre exploitation coloniale et émergence d'une conscience politique ? »
\item \textbf{Synthétiser les acquis majeurs (trois piliers) :}
  \begin{itemize}
  \item \textbf{Structuration territoriale et institutionnelle :} partition franco-britannique, dualisme juridique, scolaire et linguistique, frontières actuelles.
  \item \textbf{Exploitation et résistances :} économie de plantation, travail forcé, fiscalité, mais aussi premières résistances (mouvements religieux, grèves, pétitions).
  \item \textbf{Rupture de 1940-1946 :} ralliement à la France libre, effort de guerre, abolition du travail forcé (1946), entrée dans le régime de tutelle de l'ONU.
  \end{itemize}
\item \textbf{Ouvrir sur la suite (transition) :} « La tutelle (1946-1960) va transformer ce cadre : l'ONU impose un contrôle, les libertés publiques s'élargissent, les partis politiques naissent (UPC, partis anglophones), conduisant aux indépendances de 1960 (Cameroun français) et 1961 (réunification). Le dualisme hérité du mandat devient le défi majeur de l'État unifié. »
\item \textbf{Formule finale :} « Le mandat a dessiné la carte ; la tutelle a écrit la loi ; l'indépendance a signé l'acte de naissance, mais le dualisme colonial reste inscrit dans les fondations du Cameroun. »
\end{enumerate}
\end{methode}

\begin{exoresolu}[Sujet de composition type Probatoire]
\textbf{Sujet :} « La période de mandat (1919-1946) a-t-elle été une simple parenthèse coloniale ou le fondement structurel du Cameroun contemporain ? »

\textbf{Plan détaillé et éléments de réponse :}

\textbf{Introduction}
\begin{itemize}
\item Accroche : référence au partage franco-britannique de 1919, conséquence directe de la Première Guerre mondiale.
\item Définition : mandat de la SDN, partage du territoire en deux zones administrées séparément.
\item Problématique : voir ci-dessus.
\item Annonce du plan : 1) fondements structurels (territoire, institutions, dualismes) ; 2) exploitation et résistances (économie, société) ; 3) rupture de 1940-1946 (guerre, ralliement, tutelle).
\end{itemize}

\textbf{I. Le mandat comme matrice territoriale et institutionnelle (1919-1939)}
\begin{itemize}
\item A. Partition et frontières : deux entités distinctes, avec capitales et administrations séparées.
\item B. Dualisme administratif : centralisation française contre \emph{indirect rule} britannique.
\item C. Dualisme juridique et scolaire : droit civil contre common law ; école francophone contre école anglophone confessionnelle.
\item D. Héritage : ces structures survivent à l'indépendance (1960-1961) et à l'unification (1972).
\end{itemize}

\textbf{II. Une économie d'exploitation et des sociétés en résistance}
\begin{itemize}
\item A. Économie de plantation et d'extraction : travail forcé, port de Douala, chemin de fer, monopoles commerciaux.
\item B. Fiscalité et main-d'œuvre : impôt de capitation, prestations, migrations contraintes.
\item C. Résistances multiformes : mouvements religieux, grèves urbaines, pétitions adressées à la SDN, fuites vers la zone britannique.
\item D. Émergence d'une élite lettrée (écoles missionnaires, rares boursiers) porteuse de revendications.
\end{itemize}

\textbf{III. 1940-1946 : la guerre comme accélérateur de l'histoire}
\begin{itemize}
\item A. Ralliement à la France libre (27 août 1940, Leclerc) : basculement géopolitique, engagement militaire des Camerounais.
\item B. Économie de guerre : réquisitions maximales, disettes, prise de conscience des inégalités.
\item C. Premières réformes : Conférence de Brazzaville (1944), abolition du travail forcé (loi du 11 avril 1946).
\item D. Création de l'ONU et accords de tutelle (13 décembre 1946) : nouveau cadre juridique, objectif d'autonomie, contrôle international, ouverture d'un espace politique (partis, syndicats, assemblées).
\end{itemize}

\textbf{Conclusion}
\begin{itemize}
\item Bilan : le mandat n'est pas une parenthèse : il a \textbf{façonné le territoire, les institutions, les dualismes et les frontières}.
\item Mais la tutelle (1946-1960) transforme l'héritage : l'ONU impose une évolution politique, les nationalismes s'organisent, l'indépendance se prépare.
\item Ouverture : la réunification de 1961 tente de dépasser le dualisme, mais la question anglophone montre la permanence de l'héritage mandataire.
\end{itemize}
\end{exoresolu}

\newpage

\coursec{Activité d'intégration}

\coursec{La période de mandat : Activité d'intégration — Le Cameroun sous mandat (1919-1946)}

\courssub{Objectifs de l'activité}
\noindent Cette activité d'intégration vise à mobiliser l'ensemble des savoirs et savoir-faire acquis au cours de la leçon sur la période de mandat au Cameroun (1919-1946). Elle permet à l'élève de :
\begin{itemize}
    \item \textbf{Maîtriser la chronologie et la géographie} du partage du Kamerun entre la France et la Grande-Bretagne à l'issue de la Première Guerre mondiale.
    \item \textbf{Analyser et comparer} les modes d'administration (administration directe française vs \emph{indirect rule} britannique) et leurs impacts concrets sur les populations (impôts, travail forcé, éducation, justice, économie).
    \item \textbf{Pratiquer la cartographie historique} : construction d'une carte légendée, orientée, échelonnée, respectant les conventions sémiologiques.
    \item \textbf{Exploiter un dossier documentaire hétérogène} (traité, cartes, données chiffrées, photographies, témoignage oral) pour en extraire des informations pertinentes.
    \item \textbf{Rédiger un texte argumenté} (pétition à la SDN) structuré, étayé par des faits précis, adoptant un point de vue historique situé.
\end{itemize}

\courssub{Situation}
\noindent \textbf{Contexte :} Nous sommes en \textbf{1935}. La Commission permanente des mandats (CPM) de la Société des Nations (SDN) siège à Genève pour examiner les rapports annuels fournis par la France et le Royaume-Uni sur l'administration de leurs territoires sous mandat au Cameroun. En tant que délégués de la \textbf{Chefferie supérieure de Bafoussam} (région des Grassfields, zone de contact entre les deux administrations), vous êtes mandatés par le Fo'o (chef supérieur) pour porter la voix de votre communauté auprès de la Commission. Votre peuple a été coupé en deux par la frontière de 1919 : une partie vit dans le \textbf{Cameroun français} (circonscription du Noun, puis région de l'Ouest), l'autre dans les \textbf{British Cameroons} (division du Bamenda, province du Nord).

\noindent \textbf{Dossier documentaire mis à votre disposition (extraits) :}

\medskip
\noindent \textbf{Document 1 : Extrait de l'article 119 du Traité de Versailles (28 juin 1919).}
\emph{« L'Allemagne renonce en faveur des Principales Puissances alliées et associées à tous ses droits et titres sur ses territoires d'outre-mer... Le Kamerun est confié sous mandat à la France et à la Grande-Bretagne. »}

\medskip
\noindent \textbf{Document 2 : Données chiffrées du partage (Arrêté Milner-Simon, 10 juillet 1919, confirmé par la SDN en 1922).}

\begin{center}
\begin{tabularx}{\linewidth}{|l|X|X|}\hline
\rowcolor{popblueL}
\textbf{Indicateur} & \textbf{Cameroun français (4/5 Est)} & \textbf{Cameroun britannique (1/5 Ouest)} \\ \hline
Superficie & \num{432 000} km\textsuperscript{2} ($ \approx 91\% $) & \num{88 000} km\textsuperscript{2} ($ \approx 9\% $) \\ \hline
Population (estim. 1922) & \num{1 500 000} hab. & \num{300 000} hab. \\ \hline
Chef-lieu / Capitale & Duala (puis Yaoundé 1921) & Buéa \\ \hline
Monnaie & Franc CFA (dès 1945, Franc français avant) & Livre sterling / Shilling \\ \hline
Langue officielle & Français & Anglais \\ \hline
Voie d'accès principal & Chemin de fer Transcamerounais (Douala-Ngaoundéré) & Route et chemin de fer vers le Nigéria (Port Harcourt) \\ \hline
\end{tabularx}
\end{center}

\medskip
\noindent \textbf{Document 3 : Témoignage oral recueilli en 1934 auprès de Fo'o N., chef de groupement Bamiléké partagé par la frontière.}
\emph{« Avant l'arrivée des Blancs, nos marchés étaient un seul flux. Aujourd'hui, mes sujets du Nord paient l'impôt en shillings aux Anglais qui les jugent selon nos coutumes, mais leur interdisent de vendre leur café au meilleur prix à Douala. Mes sujets du Sud paient en francs aux Français, construisent la voie ferrée de force (prestations), vont à l'école laïque, mais ne peuvent plus rendre la justice selon la tradition. La frontière a coupé les familles, les chefferies et les esprits des ancêtres. »}

\medskip
\noindent \textbf{Document 4 : Photographies (légendées).}
\begin{itemize}
    \item \textit{Photo A :} Chantier du Transcamerounais (vers 1930) : travailleurs africains (prestataires) posant des rails sous la surveillance de contremaîtres français. Légende : « Travail forcé (prestations) pour l'intérêt général : 12 jours/an/homme valide. »
    \item \textit{Photo B :} Tribunal indigène à Bamenda (vers 1935) : un chef traditionnel (Warrant Chief) juge un litige foncier assisté d'un officier de district britannique. Légende : « Indirect Rule : justice coutumière supervisée par l'Administration. »
    \item \textit{Photo C :} École de la mission catholique à Dschang (zone française) vs École gouvernementale à Bamenda (zone britannique). Différence de programmes, de langue, de calendrier.
\end{itemize}

\medskip
\noindent \textbf{Document 5 : Carte muette du Cameroun (fournie sur feuille A3) :} Contours du Kamerun 1911, tracé de la frontière Milner-Simon (1919), fleuves principaux (Sanaga, Bénoué, Logone), villes repères (Douala, Yaoundé, Buéa, Bamenda, Garoua, Ngaoundéré).

\courssub{Consignes}
\noindent Travaillant en groupes de 4 à 5 élèves, vous réaliserez les trois productions suivantes pour la séance de restitution orale (15 min par groupe) :

\begin{enumerate}
    \item \textbf{Cartographie (30 min) :} Sur la carte muette fournie, construisez la \textbf{carte de la partition du Cameroun sous mandat (1919-1946)}. Elle doit comporter impérativement :
    \begin{itemize}
        \item Un \textbf{titre} précis et daté.
        \item Une \textbf{légende organisée} (zones française/britannique, capitale, chefs-lieux de circonscription/province, chemin de fer Transcamerounais, route vers le Nigéria, frontière internationale, frontière de mandat).
        \item L'\textbf{orientation} (flèche Nord) et une \textbf{échelle graphique} (1 cm = 100 km).
        \item Le coloriage distinct des deux zones (ex: rose pour la France, rose pâle/rayé pour UK) et le tracé du chemin de fer.
    \end{itemize}
    \item \textbf{Tableau comparatif synthétique (20 min) :} À partir des documents 2, 3, 4 et de vos connaissances du cours, dressez un tableau comparant la \textbf{vie quotidienne} dans les deux zones sur les quatre piliers suivants : \textbf{Fiscalité \& Travail}, \textbf{Éducation}, \textbf{Justice \& Administration}, \textbf{Économie \& Circulation}. Utilisez un vocabulaire précis (\emph{prestations}, \emph{indirect rule}, \emph{assimilation}, \emph{chefferies de warrant}, \emph{code de l'indigénat}).
    \item \textbf{Rédaction de la pétition à la SDN (40 min) :} Rédigez, \textbf{à la première personne du pluriel (« Nous, chefs et notables de... »)}, une pétition d'environ \textbf{15 à 20 lignes} adressée à la Commission permanente des mandats. Le texte doit :
    \begin{itemize}
        \item \textbf{Identifier} l'expéditeur et le destinataire (en-tête formel).
        \item \textbf{Exposer le fait générateur} : la partition de 1919 et la rupture de l'unité territoriale/ethnique.
        \item \textbf{Argumenter} en opposant ou en nuancant les deux systèmes (ex: « D'un côté..., de l'autre... ») en citant \textbf{au moins 3 faits précis} issus du dossier (chiffres, noms de lois, exemples concrets).
        \item \textbf{Formuler une demande claire} (révision de la frontière, harmonisation des droits, suppression du travail forcé, libre circulation...).
        \item Respecter le \textbf{ton diplomatique mais ferme} propre aux pétitions de l'entre-deux-guerres.
    \end{itemize}
\end{enumerate}

\courssub{Ressources à mobiliser}
\noindent Pour réussir cette activité, vous devez réactiver les connaissances et méthodes suivantes :

\begin{aretenir}[Savoirs clés de la leçon N06]
\textbf{1. Chronologie repère :} 1914-1916 (Campagne du Kamerun) $\rightarrow$ 1919 (Traité de Versailles, Arrêté Milner-Simon) $\rightarrow$ 1922 (Confirmation mandats SDN) $\rightarrow$ 1939-1945 (Seconde Guerre mondiale) $\rightarrow$ 1946 (Tutelle ONU).
\textbf{2. Les deux modèles administratifs :}
\begin{itemize}
    \item \textbf{Zone française :} \cle{Administration directe}. Division en circonscriptions $\rightarrow$ subdivisions $\rightarrow$ postes. Chefs nommés/révoqués par l'administrateur. \cle{Code de l'indigénat} (justice sommaire, travail forcé/prestations 12 j/an, impôt de capitation). Politique d'\cle{assimilation} culturelle (école laïque française, langue française).
    \item \textbf{Zone britannique :} \cle{Indirect Rule} (Lugard). Administration via les \cle{chefferies traditionnelles} (Warrant Chiefs) sous contrôle du \emph{District Officer}. Justice coutumière (\cle{Native Courts}) codifiée. Pas de code de l'indigénat formel, mais travail forcé déguisé (routes). Administration depuis le \cle{Nigéria} (intégration économique/administrative). École majoritairement missionnaire (anglais).
\end{itemize}
\textbf{3. Conséquences de la partition :} Morcellement ethnique (Bamiléké, Bamoun, Tikar, Peuls, Haoussa), rupture des circuits commerciaux traditionnels (vers Douala vs vers Port Harcourt), double système monétaire/juridique/linguistique.
\end{aretenir}

\begin{methode}[Construire une carte historique commentée]
\begin{enumerate}
    \item \textbf{Analyser la consigne} : thème, date, échelle, éléments obligatoires.
    \item \textbf{Structurer la légende} : hiérarchiser l'information (Surfaces : zones ; Linéaires : frontières, voies de communication ; Ponctuels : capitales, villes).
    \item \textbf{Choisir la sémiologie} : couleurs contrastées pour les zones (hachures si impression N/B), symboles normalisés (étoile = capitale, cercle = ville, trait plein = frontière int., trait mixte = mandat, trait pointillé = chemin de fer).
    \item \textbf{Réaliser le tracé} : légèreté au crayon, encrage final, coloriage uniforme.
    \item \textbf{Vérifier} : Titre ? Légende complète ? Nord ? Échelle ? Source ? Propreté.
\end{enumerate}
\end{methode}

\begin{methode}[Rédiger une pétition / texte argumenté historique]
\begin{enumerate}
    \item \textbf{Identifier l'énonciateur et le destinataire} (contexte institutionnel).
    \item \textbf{Poser le diagnostic} (constat factuel, chiffré, daté).
    \item \textbf{Structurer l'argumentation} : Thèse (partition = préjudice) $\rightarrow$ Arguments 1, 2, 3 (comparaison binaire zone FR / zone UK) $\rightarrow$ Nuance (reconnaissance d'apports : santé, routes, écriture) $\rightarrow$ Revendication.
    \item \textbf{Utiliser le vocabulaire spécifique} du programme (mandat, SDN, indigénat, indirect rule, prestations, chefferie de warrant).
    \item \textbf{Soigner la forme} : salutations protocolaires, paragraphes aérés, absence de fautes majeures.
\end{enumerate}
\end{methode}

\courssub{Démarche et corrigé commenté}
\begin{exoresolu}[Activité d'intégration : Partition du Cameroun - Corrigé type]
\textbf{Énoncé :} Réaliser la carte, le tableau comparatif et la pétition à la SDN (voir consignes ci-dessus).

\tcblower

\textbf{PARTIE 1 : LA CARTE DE LA PARTITION (Corrigé type)}

\noindent \textbf{Titre suggéré :} \emph{« Le Cameroun sous mandat français et britannique (1919-1946) : partage territorial et infrastructures principales »}

\begin{popfigure}
\begin{tikzpicture}[scale=0.75, every node/.style={font=\footnotesize}]
    % Échelle : 1 cm = 100 km environ. Carte ~10cm large = 1000km (largeur réelle ~700-800km, acceptable pour schéma)
    % Coordonnées approximatives (x=Est, y=Nord). Origine ~Douala.
    % Contour simplifié du Kamerun 1914 (sans Neukamerun)
    \draw[popdark, very thick] (0,0) coordinate (Douala) -- (-0.5,1.5) coordinate (Victoria) -- (-1,3) -- (0,5) -- (1.5,7) -- (4,10) coordinate (LacTchad) -- (6,8) -- (6,3) -- (4,0.5) -- (0,0);
    
    % Frontière de mandat Milner-Simon (approx 1/5 à l'Ouest)
    \draw[popblue, thick, dash pattern=on 5pt off 3pt] (-0.2,0.5) -- (0.5,3) -- (1,6) -- (1.5,9);
    
    % Zones colorées
    \fill[poppink!30] (0,0) -- (4,0.5) -- (6,3) -- (6,8) -- (4,10) -- (1.5,9) -- (1,6) -- (0.5,3) -- (-0.2,0.5) -- (-1,3) -- (-0.5,1.5) -- (0,0); % Zone FR (Est) - grosso modo à droite de la ligne
    \fill[popblue!20] (0,0) -- (-0.5,1.5) -- (-1,3) -- (0,5) -- (0.5,3) -- (-0.2,0.5) -- (0,0); % Zone UK (Ouest) - bande étroite à gauche
    
    % Correction: Le coloriage TikZ "à la main" est complexe. On utilise des "clips" ou on dessine les deux polygones séparément.
    % Pour simplifier le code et garantir la compilation : on dessine le contour, puis on remplit la zone UK (gauche), puis la zone FR (droite) par différence ou polygones distincts.
    % Ici, je redéfinis proprement les deux polygones.
    
    % Zone UK (Ouest) - Polygone : Côte -> Nord-Ouest -> Frontière Mandat -> Sud
    \fill[popblue!20] (-1,3) -- (-0.5,1.5) -- (0,0) -- (-0.2,0.5) -- (0.5,3) -- (1,6) -- (1.5,9) -- (-1,5) -- cycle; % Approx
    % Zone FR (Est) - Le reste
    \fill[poppink!30] (0,0) -- (4,0.5) -- (6,3) -- (6,8) -- (4,10) -- (1.5,9) -- (1,6) -- (0.5,3) -- (-0.2,0.5) -- (0,0);
    
    % Re-tracer les frontières par-dessus
    \draw[popdark, very thick] (0,0) -- (-0.5,1.5) -- (-1,3) -- (0,5) -- (1.5,7) -- (4,10) -- (6,8) -- (6,3) -- (4,0.5) -- (0,0);
    \draw[popblue, thick, dash pattern=on 5pt off 3pt] (-0.2,0.5) -- (0.5,3) -- (1,6) -- (1.5,9);
    
    % Villes (cercles)
    \node[circle, fill=popdark, inner sep=1.5pt, label=below right:\textbf{Douala}] at (0.5,1) {};
    \node[circle, fill=popdark, inner sep=1.5pt, label=right:\textbf{Yaoundé}] at (2.5,3) {};
    \node[circle, fill=popdark, inner sep=1.5pt, label=right:\textbf{Ngaoundéré}] at (3,7) {};
    \node[circle, fill=popdark, inner sep=1.5pt, label=right:\textbf{Garoua}] at (4.5,8) {};
    \node[circle, fill=popdark, inner sep=1.5pt, label=right:\textbf{Maroua}] at (5,9) {};
    
    \node[circle, fill=popblue, inner sep=1.5pt, label=left:\textbf{Victoria (Limbe)}] at (-0.5,1.5) {};
    \node[circle, fill=popblue, inner sep=1.5pt, label=left:\textbf{Buéa}] at (0,2.5) {};
    \node[circle, fill=popblue, inner sep=1.5pt, label=left:\textbf{Kumba}] at (0.2,2.8) {};
    \node[circle, fill=popblue, inner sep=1.5pt, label=left:\textbf{Bamenda}] at (0.8,5) {};
    
    % Capitales (étoiles)
    \node[star, star points=5, star point ratio=2.25, fill=poppink, draw=popdark, inner sep=2pt, label=right:\textbf{Yaoundé (Cap. FR)}] at (2.5,3) {};
    \node[star, star points=5, star point ratio=2.25, fill=popblue, draw=popdark, inner sep=2pt, label=left:\textbf{Buéa (Cap. UK)}] at (0,2.5) {};
    
    % Chemins de fer / Routes
    \draw[popdark, very thick, double=popdark!50] (0.5,1) -- (2.5,3) -- (3,7); % Transcamerounais
    \draw[poporange, thick, dash pattern=on 3pt off 3pt] (0.8,5) -- (-1.5,5) node[left, pos=1.2, poporange] {\textbf{Port Harcourt (Nigéria)}}; % Vers Nigéria
    
    % Fleuves
    \draw[popblue!70, thick] (0.5,1) .. controls (2,2) .. (2.5,3); % Sanaga approx
    \draw[popblue!70, thick] (3,7) .. controls (4,7.5) .. (4.5,8); % Bénoué approx
    \draw[popblue!70, thick] (4,10) -- (5,9.5); % Logone/Chari vers Lac Tchad
    
    % Lac Tchad
    \node[align=center, text=popblue!70] at (4,10.5) {\textbf{Lac Tchad}};
    
    % Flèche Nord
    \draw[popdark, thick, ->] (7, 10.5) -- (7, 11.5) node[above] {\textbf{N}};
    
    % Échelle graphique 1cm = 100km
    \draw[popdark, thick] (7, 0.5) -- (8, 0.5) node[midway, below] {100 km};
    \draw[popdark, thick] (7, 0.4) -- (7, 0.6);
    \draw[popdark, thick] (8, 0.4) -- (8, 0.6);
    \node[below] at (7.5, 0.3) {\textbf{Échelle : 1 cm = 100 km}};
    
    % Légende interne (optionnel, mais la légende formelle est dans \legende)
\end{tikzpicture}
\legende{Carte schématique de la partition du Cameroun (1919-1946). \textbf{Rose} : Cameroun français (administration directe, capitale Yaoundé). \textbf{Bleu clair} : Cameroun britannique (Indirect Rule, capitale Buéa, administré depuis Lagos/Nigéria). \textbf{Trait plein noir} : Frontière internationale (1911/1919). \textbf{Trait mixte bleu} : Frontière de mandat Milner-Simon (1919). \textbf{Double trait noir} : Chemin de fer Transcamerounais (Douala-Yaoundé-Ngaoundéré). \textbf{Trait pointillé orange} : Liaison routière/ferroviaire vers Port Harcourt (Nigéria). \textbf{Étoiles} : Capitales. \textbf{Cercles noirs} : Villes zone française. \textbf{Cercles bleus} : Villes zone britannique. \textbf{Flèche N} : Orientation. \textbf{Barre} : Échelle graphique.}
\end{popfigure}

\noindent \textbf{Éléments attendus sur la copie élève (check-list) :}
\begin{itemize}
    \item Titre précis, daté (1919-1946), mentionnant les deux mandats.
    \item Flèche Nord (haut de carte).
    \item Échelle graphique : 1 cm = 100 km (soit $\approx$ 7-8 cm pour la largeur Est-Ouest du pays).
    \item Coloriage net, sans débordement sur la frontière de mandat ; distinction visuelle claire zone FR / zone UK.
    \item Légende encadrée, organisée par thèmes (Limites, Zones, Villes/Capitales, Transports, Hydrographie), symboles cohérents avec la carte.
    \item Distinction \textbf{Frontière internationale} (trait plein) vs \textbf{Frontière de mandat} (trait mixte).
    \item Tracé du Transcamerounais (Douala $\rightarrow$ Yaoundé $\rightarrow$ Ngaoundéré) et flèche vers Port Harcourt (Nigéria).
    \item Indication « Adminéré depuis Lagos (Nigéria) » pour la zone britannique (flèche ou mention).
\end{itemize}

\medskip
\noindent \textbf{PARTIE 2 : TABLEAU COMPARATIF DE LA VIE QUOTIDIENNE (Corrigé type)}

\begin{center}
\begin{tabularx}{\linewidth}{|l|X|X|}\hline
\rowcolor{popblueL}
\textbf{Thème} & \textbf{Zone Française (Administration Directe)} & \textbf{Zone Britannique (Indirect Rule)} \\ \hline
\textbf{Fiscalité \& Travail} & \textbf{Impôt de capitation} (francs), obligatoire, perçu par l'administration. \textbf{Prestations} : 12 jours/an/homme valide (travail forcé légal : rails, routes, portages). Sanctions pénales (prison) en cas de refus. Code de l'indigénat. & \textbf{Impôt direct} (shillings), perçu par les chefs (\emph{Native Treasuries}). \textbf{Travail communautaire} : obligation de travailler sur les routes/chef-lieu (\emph{Native Authority Ordinance}). Moins systématisé qu'en zone FR, mais contrainte réelle. Pas de code de l'indigénat formel. \\ \hline
\textbf{Éducation} & École \textbf{laïque publique} (français, programme métropole) + Écoles privées missionnaires (subventionnées). Objectif : \cle{assimilation} (évolués). Certificat d'études primaires (CEP). & École quasi exclusivement \textbf{missionnaire} (protestantes/catholiques), subventionnées par \emph{Native Authorities}. Langue : \textbf{anglais}. Programme adapté au contexte local. Pas d'école publique laïque avant 1940. \\ \hline
\textbf{Justice \& Administration} & \textbf{Code de l'indigénat} : justice d'exception, administrateur = juge/procureur. Peines : amendes, prison, relégation. Chefs = auxiliaires nommés/révoqués. Découpage : \textbf{Circonscriptions / Subdivisions}. & \textbf{Native Courts} : justice coutumière codifiée, jugée par chefs (\cle{Warrant Chiefs}) sous contrôle \emph{District Officer}. Appel possible devant cour britannique. Découpage : \textbf{Provinces / Divisions}. Chef = pivot du système. \\ \hline
\textbf{Économie \& Circulation} & Économie tournée vers \textbf{Douala} (port) et la métropole (France). Monnaie : Franc. \textbf{Transcamerounais} = colonne vertébrale. Cultures d'exportation imposées (cacao, café, coton Nord). & Économie intégrée au \textbf{Nigéria} (Port Harcourt). Monnaie : Livre/Shilling. Routes vers le Nigéria. Cultures : cacao, palme, thé (hauteurs). \textbf{Rupture des circuits traditionnels} (ex: Bamiléké vers Douala coupés). \\ \hline
\end{tabularx}
\end{center}

\medskip
\noindent \textbf{PARTIE 3 : PÉTITION À LA COMMISSION PERMANENTE DES MANDATS DE LA SDN (Corrigé type / Modèle de rédaction)}

\medskip
\noindent \textit{En-tête :}
\begin{flushleft}
\textbf{À l'attention de Monsieur le Président de la Commission Permanente des Mandats,}\\
\textbf{Société des Nations, Genève.}\\[0.5cm]
\textbf{Objet : Pétition relative aux conséquences de la partition du Kamerun sur les populations des Grassfields (Bamiléké/Bamoun).}\\[0.5cm]
\textbf{Excellences, Mesdames, Messieurs les Commissaires,}
\end{flushleft}

\noindent \textit{Corps du texte :}
\emph{Nous, Fo'o et notables des groupements Bamiléké et Bamoun, établis de part et d'autre de la frontière tracée en 1919 par l'Arrêté Milner-Simon, venons respectueusement porter à votre connaissance les souffrances engendrées par cette division artificielle de notre peuple unique.}

\emph{Depuis la confirmation des mandats en 1922, nos familles sont soumises à deux lois contraires. Au Sud, dans le Cameroun français, nos sujets paient l'impôt de capitation en francs, sont astreints aux \textbf{prestations} (12 jours de travail forcé par an) pour construire le \textbf{Transcamerounais}, et sont jugés sans avocat par l'administrateur au nom du \textbf{Code de l'indigénat}. Nos enfants vont à l'école française laïque, mais nos chefferies sont affaiblies par la nomination d'auxiliaires zélés.}

\emph{Au Nord, dans les British Cameroons administrés depuis le \textbf{Nigéria}, nos frères paient l'impôt en shillings aux \textbf{Warrant Chiefs}, jugés par des \textbf{Native Courts} qui dénaturent notre coutume, et voient leur café obligé de partir vers \textbf{Port Harcourt} alors que Douala est trois fois plus proche. L'\emph{Indirect Rule} fige des chefferies créées de toutes pièces par l'occupant.}

\emph{\textbf{Cette frontière coupe les chefferies, les marchés, les mariages et le culte des ancêtres.} Elle crée deux citoyennetés inégales pour un même peuple.}

\emph{\textbf{En vertu de l'Article 22 du Pacte de la SDN qui vous charge du « bien-être et du développement » des populations placées sous tutelle, nous demandons :}}
\begin{enumerate}
    \item La \textbf{libre circulation des personnes et des biens} à travers la ligne de démarcation pour les populations riveraines.
    \item L'\textbf{harmonisation des régimes fiscaux et du travail} (suppression effective des prestations forcées en zone française).
    \item Le \textbf{respect de l'unité coutumière} des chefferies transfrontalières dans l'exercice de la justice.
\end{enumerate}
\emph{Nous espérons que la Société des Nations saura faire primer le droit des peuples sur les intérêts stratégiques des Puissances mandataires.}\\[0.5cm]
\textit{Fait à Bafoussam, le 15 mars 1935.}\\
\textit{Signatures : Fo'o N. (Bafoussam), Fo'o M. (Bangangté), Notable K. (Bamenda-Nord), ...}

\medskip
\noindent \textbf{Commentaire pédagogique du corrigé :}
\begin{itemize}
    \item \textbf{Carte :} La réussite repose sur la distinction nette \emph{Frontière internationale} vs \emph{Frontière de mandat} (souvent confondues). Le tracé du Transcamerounais (Douala-Ngaoundéré) est le marqueur structurel de l'aménagement français. L'indication « Adminéré depuis Lagos » est cruciale pour la zone UK.
    \item \textbf{Tableau :} L'opposition \emph{Prestations (legally forced labor)} vs \emph{Travail communautaire (Native Authority)} doit être nuancée : contrainte dans les deux cas, mais cadre juridique différent. L'intégration économique au Nigéria (zone UK) vs métropole (zone FR) est le point clé pour l'économie.
    \item \textbf{Pétition :} Le modèle utilise le « Je/Nous » historique. Les arguments s'appuient sur le dossier (Doc 2 : monnaies, capitales ; Doc 3 : témoignage ; Doc 4 : prestations/chemin de fer). Le vocabulaire technique (\emph{indigénat, warrant chief, native courts, prestations, Milner-Simon}) est mobilisé correctement. La référence à l'Art. 22 du Pacte ancre l'action dans le droit international de l'époque. La forme (en-tête, objet, formules de politesse) est respectée.
\end{itemize}
\end{exoresolu}

\courssub{Bilan de l'activité}
\noindent Cette activité d'intégration a permis de vérifier l'appropriation des concepts centraux du chapitre :
\begin{itemize}
    \item \textbf{La partition n'est pas une ligne abstraite :} elle a découpé des espaces vécus (ethnies, bassins économiques, chefferies). La carte le visualise ; le tableau et la pétition le prouvent par le vécu quotidien (monnaie, impôt, justice, école).
    \item \textbf{Deux logiques coloniales distinctes :} L'\cle{administration directe française} (centralisée, assimilationniste, coercitive via l'indigénat/prestations) vs l'\cle{Indirect Rule britannique} (décentralisée, associationniste, instrumentalisant la tradition via les Warrant Chiefs/Native Courts). L'élève a dû manipuler ce couple conceptuel pour remplir le tableau et argumenter la pétition.
    \item \textbf{L'unité camerounaise comme construction historique :} La pétition met en lumière le paradoxe du mandat : la SDN charge les puissances du « bien-être » des populations, mais la partition même (décidée par les vainqueurs de 1919) structure une division durable (héritage linguistique, juridique, économique actuel).
    \item \textbf{Compétences transversales validées :}
    \begin{itemize}
        \item \textbf{Cartographie :} Respect des conventions (titre, légende, échelle, nord), choix sémiologiques pertinents.
        \item \textbf{Analyse documentaire :} Croisement de sources quantitatives (chiffres 1919), qualitatives (témoignage, photos), juridiques (traité).
        \item \textbf{Écriture argumentée :} Adoption d'une posture historique (énonciateur situé), structuration thèse/arguments/nuance/revendication, vocabulaire précis.
    \end{itemize}
\end{itemize}

\begin{attention}[Pièges fréquents relevés lors de la correction]
\begin{itemize}
    \item Confondre \textbf{mandat} (SDN, 1922-1946) et \textbf{tutelle} (ONU, 1946-1960) : ne pas projeter les partis politiques (UPC, KNDP) ou le référendum de 1961 dans cette activité située en 1935.
    \item Oublier que le Cameroun britannique est \textbf{administré depuis le Nigéria} (capitale Lagos, monnaie livre, flux vers Port Harcourt). C'est la clé de l'intégration économique britannique.
    \item Réduire l'\emph{Indirect Rule} à « les chefs gardent le pouvoir » : il s'agit d'une \textbf{reconstruction coloniale de la chefferie} (Warrant Chiefs) soumise au District Officer.
    \item Négliger le \textbf{Code de l'indigénat} en zone française : c'est l'outil juridique principal de la contrainte (travail, justice, déplacement).
    \item Carte : absence d'échelle graphique, légende désorganisée (mélange villes/rivières/frontières), frontière de mandat non distinguée de la frontière internationale.
\end{itemize}
\end{attention}

\begin{savaistu}[Le Cameroun, « laboratoire » du mandat de la SDN]
Le Cameroun est l'\textbf{unique territoire africain} partagé entre deux puissances mandataires (France et UK) sous l'égide de la SDN. Cela en fait un cas d'école unique pour l'histoire du droit international : la Commission permanente des mandats recevait \textbf{deux rapports annuels distincts} pour un même peuple. Les pétitions réelles (comme celles des chefs Bamiléké/Bamoun dans les années 30 adressées à la SDN, ou celle de Rudolf Duala Manga Bell en 1914 adressée aux autorités allemandes) sont conservées aux Archives de la SDN (Genève) et aux Archives Nationales du Cameroun (Yaoundé). Elles prouvent que les populations n'étaient pas passives : elles utilisaient le droit international pour contester la partition, le travail forcé ou l'arbitraire administratif, préfigurant le nationalisme des années 50.
\end{savaistu}

\newpage

\coursec{Méthodes et exercices résolus}

\coursec{La période de mandat}

\courssub{Méthode 1 : Appliquer les notions de l'unité à une situation concrète}

\begin{methode}[Mobiliser les notions du cours dans une situation de la vie courante]
Pour appliquer les notions de l'unité (partition, mandat, administration directe, \cle{indirect rule}) à une situation concrète :
\begin{enumerate}
\item \textbf{Lire attentivement la situation} proposée et repérer les indices : dates (1916, 1922, 1946...), lieux (Douala, Buea, Yaoundé...), acteurs (administrateur français, chef traditionnel, planteur...).
\item \textbf{Identifier la notion du cours} correspondant à chaque indice : par exemple, un chef coutumier qui collecte l'impôt pour le compte des Britanniques renvoie à l'\cle{indirect rule}.
\item \textbf{Rappeler la définition exacte} de la notion (date, traité ou texte fondateur : traité de Versailles 1919, mandats de la SDN 1922).
\item \textbf{Faire le lien explicite} entre la situation et la notion : « cette situation illustre... car... ».
\item \textbf{Conclure} en montrant ce que la situation révèle de la période du mandat (conséquence de la Première Guerre mondiale, différence entre administration française et britannique).
\end{enumerate}
\textbf{Réflexe :} toujours dater et localiser avant de nommer la notion. Une réponse sans repère chronologique perd des points.
\end{methode}

\begin{exoresolu}[Une situation concrète : le chef et l'impôt]
\textbf{Énoncé.} En 1930, dans la région de Bamenda (Cameroun britannique), un chef traditionnel collecte l'impôt auprès des villageois et le reverse à l'officier de district britannique, qui le laisse par ailleurs juger les petits litiges selon la coutume. À l'aide de tes connaissances, explique quel mode d'administration cette situation illustre et pourquoi les Britanniques l'ont choisi.
\tcblower
\textbf{Solution détaillée.}
\begin{enumerate}
\item \textbf{Repérage des indices.} La date (1930) situe la scène dans la période du \cle{mandat} (1922-1946). Le lieu, Bamenda, se trouve dans la partie occidentale du Cameroun (Southern Cameroons), attribuée à la Grande-Bretagne après la défaite allemande de 1916 et confirmée par les mandats de la SDN en 1922. Les acteurs sont un chef traditionnel et un officier britannique (District Officer).
\item \textbf{Notion concernée.} Le fait que le chef coutumier administre la population pour le compte de la puissance mandataire renvoie à l'\cle{indirect rule} (administration indirecte).
\item \textbf{Définition.} L'indirect rule, théorisé par Lord Lugard au Nigeria, consiste à gouverner les populations colonisées par l'intermédiaire de leurs propres chefs et institutions traditionnelles (Native Authorities, Native Courts, Native Treasuries), encadrés par un petit nombre d'administrateurs britanniques.
\item \textbf{Lien avec la situation.} Le chef de Bamenda collecte l'impôt (Native Treasury) et rend la justice coutumière (Native Court) : il exerce les fonctions administratives à la place des Britanniques, mais sous leur contrôle (il reverse l'impôt à l'officier de district). C'est exactement le fonctionnement de l'indirect rule.
\item \textbf{Justification du choix britannique.} Ce système coûte peu cher (peu de fonctionnaires européens), s'appuie sur l'autorité déjà reconnue des chefs, et limite les résistances en donnant l'apparence du respect des coutumes.
\end{enumerate}
\textbf{Conclusion.} Cette situation illustre l'indirect rule britannique au Cameroun : la Grande-Bretagne, puissance mandataire depuis 1922, administre sa partie du Cameroun à moindre coût en s'appuyant sur les chefs traditionnels, contrairement à la France qui pratique l'administration directe.
\end{exoresolu}

\courssub{Méthode 2 : Rédiger et justifier une réponse argumentée}

\begin{methode}[Construire une réponse rédigée et justifiée en Histoire]
Pour rédiger une réponse justifiée (question de synthèse ou de commentaire) :
\begin{enumerate}
\item \textbf{Analyser le sujet} : souligner les mots-clés, délimiter le temps (1914-1946) et l'espace (Cameroun français / britannique).
\item \textbf{Formuler une réponse claire dès la première phrase} (prise de position ou réponse directe à la question).
\item \textbf{Développer avec des faits précis} : chaque argument s'appuie sur une date, un traité, un acteur (traité de Versailles 1919, SDN, mandats de 1922, conférence de Milner-Simon 1919...).
\item \textbf{Justifier chaque affirmation} par un exemple concret (impôt, travail forcé, chemin de fer, plantations...).
\item \textbf{Conclure} en répondant à nouveau au sujet et, si possible, en ouvrant (du mandat à la tutelle en 1946).
\end{enumerate}
\textbf{Formule à retenir :} une affirmation + une preuve datée + une explication = un paragraphe valable.
\end{methode}

\begin{exoresolu}[Question de synthèse : les causes de la partition du Cameroun]
\textbf{Énoncé.} « La partition du Cameroun est une conséquence directe de la Première Guerre mondiale. » Justifie cette affirmation en un paragraphe rédigé et argumenté (10 à 15 lignes).
\tcblower
\textbf{Solution détaillée.}
\begin{enumerate}
\item \textbf{Analyse du sujet.} Mots-clés : partition, conséquence, Première Guerre mondiale. Période : 1914-1922. Il faut montrer la chaîne de causalité : guerre $\rightarrow$ conquête militaire $\rightarrow$ partage juridique.
\item \textbf{Rédaction.}

Avant 1914, le Cameroun est une colonie allemande, le \emph{Kamerun}, depuis 1884. La Première Guerre mondiale bouleverse cette situation : dès août 1914, les troupes françaises et britanniques attaquent la colonie depuis le Tchad, le Nigeria et le Gabon. Après la prise de Douala (septembre 1914) et de Yaoundé (janvier 1916), les Allemands capitulent à Mora en février 1916 : le Cameroun est militairement conquis. Dès mars 1916, la convention franco-britannique de Picot-Simon organise un partage provisoire entre les deux vainqueurs. Cette occupation est ensuite légalisée par le traité de Versailles (28 juin 1919), qui dépouille l'Allemagne de ses colonies, puis par la Société des Nations qui, le 20 juillet 1922, attribue officiellement le Cameroun sous forme de \cle{mandats} : environ les quatre cinquièmes du territoire à la France, le reste (deux bandes occidentales le long du Nigeria) à la Grande-Bretagne. Ainsi, sans la guerre de 1914-1918 et la défaite allemande, le Cameroun serait resté une colonie allemande unie : sa partition en deux zones d'administration est bien une conséquence directe du conflit mondial.
\item \textbf{Vérification.} La réponse contient : une thèse claire, des dates (1884, 1914, 1916, 1919, 1922), des acteurs (France, Grande-Bretagne, Allemagne, SDN), et une conclusion qui répond au sujet.
\end{enumerate}
\textbf{Conclusion.} La défaite allemande de 1916, confirmée par le traité de Versailles et les mandats de la SDN de 1922, transforme le Cameroun uni sous domination allemande en deux territoires administrés séparément par la France et la Grande-Bretagne.
\end{exoresolu}

\courssub{Méthode 3 : Commenter un document historique}

\begin{methode}[Commentaire de document (texte ou carte) en 5 étapes]
\begin{enumerate}
\item \textbf{Présenter le document} : nature (texte, carte, photographie), auteur ou source, date, contexte (période du mandat, 1922-1946).
\item \textbf{Dégager le thème} : de quoi parle le document ? (partition, administration, vie économique...).
\item \textbf{Extraire les informations} : relever 2 ou 3 idées principales en citant des éléments précis du document.
\item \textbf{Confronter avec ses connaissances} : compléter ou expliquer le document grâce au cours (ex. : le document mentionne le mandat $\rightarrow$ rappeler la SDN et 1922).
\item \textbf{Conclure} : intérêt et limites du document pour l'historien (point de vue, date, auteur).
\end{enumerate}
\textbf{Réflexe :} ne jamais paraphraser : chaque citation du document doit être suivie d'une explication issue du cours.
\end{methode}

\begin{exoresolu}[Commentaire d'un extrait de texte sur le mandat]
\textbf{Énoncé.} Document : « Le Cameroun, ancienne colonie allemande, est placé sous mandat de la Société des Nations. La France et la Grande-Bretagne administrent leurs zones respectives au nom de la communauté internationale et doivent lui rendre compte. » (Extrait d'un manuel, résumant la décision de la SDN de 1922.) Présente ce document, dégage son idée principale et explique ce qu'est un mandat de la SDN.
\tcblower
\textbf{Solution détaillée.}
\begin{enumerate}
\item \textbf{Présentation.} Il s'agit d'un texte historique (extrait de manuel) qui résume la décision de la Société des Nations de 1922 plaçant le Cameroun sous \cle{mandat}. Le contexte est l'après-Première Guerre mondiale : l'Allemagne, vaincue, perd ses colonies au profit des vainqueurs.
\item \textbf{Idée principale.} Le Cameroun n'est pas annexé comme une colonie ordinaire : il est confié par la SDN à la France et à la Grande-Bretagne, qui doivent en rendre compte.
\item \textbf{Extraction des informations.} Le texte indique trois choses : (a) le Cameroun était allemand ; (b) il est placé sous mandat de la SDN ; (c) les puissances mandataires ont une obligation de compte-rendu.
\item \textbf{Apport des connaissances.} Un \cle{mandat} est une mission d'administration confiée par la SDN (créée en 1919) à une puissance victorieuse sur un territoire pris aux vaincus. En théorie, le mandataire doit préparer le territoire à s'administrer lui-même et envoyer des rapports annuels à la Commission des mandats à Genève. En pratique, la France administre sa partie comme une colonie (administration directe, impôt de capitation, travail forcé/prestations), et la Grande-Bretagne rattache la sienne au Nigeria (indirect rule). En 1946, les mandats deviennent des territoires sous \cle{tutelle} de l'ONU.
\item \textbf{Conclusion.} Ce document est intéressant car il montre la différence juridique entre colonie et mandat ; mais il est limité car il ne dit pas que, sur le terrain, l'administration ressemble fort à la colonisation classique.
\end{enumerate}
\end{exoresolu}

\courssub{Méthode 4 : Comparer deux situations (Cameroun français / Cameroun britannique)}

\begin{methode}[Construire une comparaison rigoureuse]
\begin{enumerate}
\item \textbf{Choisir des critères de comparaison} identiques pour les deux cas : statut juridique, mode d'administration, place des chefs, économie, vie sociale.
\item \textbf{Traiter chaque critère pour les deux territoires} : jamais un bloc sur la France puis un bloc sur l'Angleterre sans points communs.
\item \textbf{Utiliser un tableau} si l'exercice s'y prête : lignes = critères, colonnes = territoires.
\item \textbf{Dégager ressemblances et différences} : ressemblance majeure = domination étrangère et exploitation ; différence majeure = administration directe contre indirect rule.
\item \textbf{Conclure} par un bilan nuancé.
\end{enumerate}
\end{methode}

\begin{exoresolu}[Comparer les deux modes d'administration]
\textbf{Énoncé.} À l'aide d'un tableau puis d'un court paragraphe, compare l'administration du Cameroun français et celle du Cameroun britannique entre 1922 et 1939 (critères : rattachement, mode d'administration, rôle des chefs traditionnels).
\tcblower
\textbf{Solution détaillée.}
\begin{center}
\begin{tabularx}{\linewidth}{|l|X|X|}
\hline
\rowcolor{popblueL} \textbf{Critère} & \textbf{Cameroun français} & \textbf{Cameroun britannique} \\
\hline
Rattachement & Administré directement depuis Yaoundé (Commissaire de la République, puis Gouverneur) & Rattaché aux provinces voisines du Nigeria (Northern et Southern Cameroons) \\
\hline
Mode d'administration & Administration \cle{directe} : fonctionnaires français (cercles, subdivisions) commandent & \cle{Indirect rule} : les Britanniques gouvernent à travers les chefs (Native Authorities) \\
\hline
Rôle des chefs & Chefs nommés ou contrôlés, simples auxiliaires de l'administration (cantonniers, collecte impôt) & Chefs traditionnels associés : impôt (Native Treasury), justice coutumière (Native Courts), ordre local \\
\hline
\end{tabularx}
\end{center}
\textbf{Paragraphe de synthèse.} Entre 1922 et 1939, les deux puissances mandataires administrent leur zone sous le contrôle théorique de la SDN, mais de manières très différentes. La France applique l'administration directe : ses fonctionnaires exercent le pouvoir à tous les niveaux et les chefs ne sont que des auxiliaires. La Grande-Bretagne, fidèle à l'indirect rule expérimenté au Nigeria, gouverne à travers les chefs coutumiers, ce qui coûte moins cher. Dans les deux cas cependant, la population camerounaise reste dominée, imposée et exploitée : la différence porte sur la méthode, non sur la nature de la domination.
\textbf{Conclusion.} La comparaison par critères montre une opposition de méthodes (directe / indirecte) au sein d'une même réalité : la domination coloniale sous couvert de mandat.
\end{exoresolu}

\courssub{Méthode 5 : Construire et exploiter une frise chronologique}

\begin{methode}[Frise chronologique : construire et lire]
\begin{enumerate}
\item \textbf{Choisir l'échelle} : pour 1914-1946, 1 cm pour 4 ou 5 ans convient.
\item \textbf{Placer les dates repères} dans l'ordre : 1914 (début de la guerre au Cameroun), 1916 (capitulation de Mora, partage provisoire), 1919 (traité de Versailles), 1922 (mandats de la SDN), 1939 (Seconde Guerre mondiale), 1946 (tutelle de l'ONU).
\item \textbf{Colorier les grandes périodes} : guerre et conquête (1914-1916), mandat (1922-1939), transition (1939-1946).
\item \textbf{Exploiter la frise} : repérer les ruptures (1916, 1922, 1946) et les continuités (la domination étrangère).
\end{enumerate}
\textbf{Réflexe :} au Bac, une frise demandée doit comporter un titre, une échelle et des dates accompagnées d'une légende courte.
\end{methode}

\begin{exoresolu}[Construire la frise de la période du mandat]
\textbf{Énoncé.} Construis une frise chronologique de 1914 à 1946 comportant au moins six dates clés de l'histoire du Cameroun, puis explique en trois lignes pourquoi 1922 est une date charnière.
\tcblower
\textbf{Solution détaillée.}

\begin{popfigure}
\begin{tikzpicture}[x=0.32cm]
\draw[->,thick,popdark] ({1913-1913},0) -- ({1947-1913},0);
\foreach \x in {1915,1920,1925,1930,1935,1940,1945}
  \draw[popmuted] ({\x-1913},0.08) -- ({\x-1913},-0.08);
% périodes colorées
\fill[poporange!40] ({1914-1913},0.15) rectangle ({1916-1913},0.55);
\fill[popblue!35] ({1922-1913},0.15) rectangle ({1939-1913},0.55);
\fill[popgreen!35] ({1939-1913},0.15) rectangle ({1946-1913},0.55);
\node[font=\scriptsize] at ({1915-1913},0.85) {Guerre};
\node[font=\scriptsize] at ({1930.5-1913},0.85) {Mandat};
\node[font=\scriptsize] at ({1942.5-1913},0.85) {Transition};
% dates clés
\foreach \x/\t in {1914/{Guerre au\\Cameroun},1916/{Mora :\\partage},1919/{Versailles},1922/{Mandats\\SDN},1939/{2$^{\text{e}}$ guerre\\mondiale},1946/{Tutelle\\ONU}}
  {\draw[popdark] ({\x-1913},0) -- ({\x-1913},-0.7); \node[below,font=\scriptsize,align=center] at ({\x-1913},-0.7) {\textbf{\x}\\\t};}
\end{tikzpicture}
\legende{Frise chronologique du Cameroun, 1914-1946 : de la conquête franco-britannique à la tutelle de l'ONU.}
\end{popfigure}

\textbf{Explication de la date charnière.} 1922 est une date charnière car la SDN attribue officiellement les mandats : le partage militaire provisoire de 1916 devient un partage juridique international. Le Cameroun cesse d'être un territoire conquis pour devenir deux territoires sous mandat, français et britannique, statut qui durera jusqu'à la tutelle de 1946.
\textbf{Conclusion.} La frise fait apparaître trois temps : la conquête (1914-1916), le mandat (1922-1939) et la transition vers la tutelle (1939-1946), dont 1922 est le pivot juridique.
\end{exoresolu}

\courssub{Méthode 6 : Lire et exploiter une carte ou un croquis de partition}

\begin{methode}[Analyser un croquis de partage territorial]
\begin{enumerate}
\item \textbf{Identifier le titre, l'orientation et la légende} du croquis.
\item \textbf{Localiser les deux zones} : la grande zone orientale (mandat français, capitale Yaoundé, port de Douala) et les deux bandes occidentales le long du Nigeria (mandat britannique, Buea, Bamenda).
\item \textbf{Mesurer le rapport de taille} : la France reçoit environ les $\dfrac{4}{5}$ du territoire, la Grande-Bretagne environ $\dfrac{1}{5}$.
\item \textbf{Interpréter} : le partage suit la ligne de séparation des armées en 1916 (accord Picot-Simon), confirmé en 1919 (Milner-Simon) et 1922 (SDN).
\item \textbf{Conclure} sur les conséquences : populations séparées, économies orientées différemment (vers le Nigeria à l'ouest, vers l'AEF à l'est).
\end{enumerate}
\end{methode}

\begin{exoresolu}[Exploiter un croquis de la partition]
\textbf{Énoncé.} Sur un croquis du Cameroun en 1922, la zone française couvre environ 432 000 km\textsuperscript{2} et la zone britannique environ 108 000 km\textsuperscript{2}, pour un total d'environ 540 000 km\textsuperscript{2}. (a) Calcule la part de chaque puissance en pourcentage. (b) Explique l'origine de ce partage.
\tcblower
\textbf{Solution détaillée.}

(a) \textbf{Calculs des pourcentages.}
\begin{align*}
\text{Part française} &= \dfrac{432\,000}{540\,000} \times 100 = 80\ \% \\
\text{Part britannique} &= \dfrac{108\,000}{540\,000} \times 100 = 20\ \%
\end{align*}
Vérification : $80 + 20 = 100\ \%$. La France reçoit donc les quatre cinquièmes du territoire ($\dfrac{4}{5} = 0,8 = 80\ \%$), la Grande-Bretagne un cinquième.

(b) \textbf{Origine du partage.} Ce partage résulte de la conquête militaire du Cameroun allemand pendant la Première Guerre mondiale : après la capitulation allemande de Mora (février 1916), les vainqueurs fixent une ligne de partage provisoire (accord Picot-Simon, mars 1916), précisée par l'accord Milner-Simon (1919) et confirmée juridiquement par les mandats de la SDN en 1922. La Grande-Bretagne, qui dispose déjà du Nigeria voisin, se contente de deux bandes frontalières occidentales ; la France, qui veut relier le Cameroun à son empire d'Afrique équatoriale, obtient l'essentiel du territoire.

\textbf{Conclusion.} Le calcul confirme le déséquilibre de la partition (80 \% / 20 \%), héritage direct du rapport de forces militaire et diplomatique de 1916-1922.
\end{exoresolu}

\courssub{Méthode 7 : Analyser la vie économique et sociale sous mandat}

\begin{methode}[Étudier les réalités vécues par les populations]
Pour traiter la vie économique et sociale (points 4 et 5 du plan), procède par thèmes communs aux deux zones :
\begin{enumerate}
\item \textbf{Fiscalité et travail :} impôt de capitation (tête), prestations (travail forcé légalisé côté français), migration de main-d'œuvre côté britannique.
\item \textbf{Économie de plantation et d'exportation :} cultures de rente (cacao, café, banane, hévéa, palme), compagnies concessionnaires (ex. CDC côté britannique, sociétés françaises côté est), chemin de fer (Douala-Yaoundé, Nord).
\item \textbf{Éducation et santé :} réseaux missionnaires (catholiques/protestants) et laïques (français), disparités fortes (taux de scolarisation plus élevé au Sud britannique).
\item \textbf{Urbain et migrations :} croissance de Douala, Yaoundé, Buea, Victoria (Limbe) ; migrations internes et vers le Nigeria (plantations, mines).
\end{enumerate}
\end{methode}

\begin{exemplebox}[Vie économique et sociale : deux expériences contrastées]
\textbf{Cameroun français (1922-1939).}
\begin{itemize}
\item \textbf{Travail forcé (prestations) :} obligation de 10 à 12 jours/an/homme pour travaux d'intérêt public (routes, chemin de fer). Réquisitions massives pour le chemin de fer Douala-Yaoundé (1927-1933) et le « Nord » (1930s) : mortalité élevée, révoltes (ex. Sanaga-Maritime).
\item \textbf{Économie de plantation :} développement des cultures d'exportation (cacao, café, banane, hévéa) par des sociétés (ex. Société des Plantations de la Mbam, HEVECAM ancêtres). Douala devient le grand port d'exportation.
\item \textbf{Fiscalité :} impôt de capitation (1923), perçu par les chefs cantonniers. Monétarisation forcée de l'économie villageoise.
\item \textbf{Scolarisation :} faible taux, monopole quasi total des missions catholiques (Spiritains, Pallottins) et protestantes (Mission de Bâle, Société des Missions de Paris), subventionnées par l'État français à partir de 1924.
\end{itemize}

\textbf{Cameroun britannique (Southern \& Northern Cameroons, 1922-1939).}
\begin{itemize}
\item \textbf{Intégration nigériane :} économie tournée vers le Nigeria (port de Lagos, chemin de fer nigérian). Pas de grand port ni de chemin de fer propre.
\item \textbf{Plantations :} Cameroon Development Corporation (CDC) créée en 1947 (mais racines dans les années 1920-30) : palmiers, hévéa, banane, thé à Tiko, Muyuka, Victoria. Main-d'œuvre migrante (Nigériens, Bamilekes).
\item \textbf{Indirect rule et fiscalité :} Native Treasuries gèrent l'impôt et financent dispensaires, écoles, routes locales. Moins de travail forcé direct, mais migration économique massive vers les plantations du Sud et les mines du Nord (Jos Plateau).
\item \textbf{Scolarisation :} missions protestantes (Bâle, Presbytérienne) et catholiques très actives. Taux de scolarisation plus élevé qu'à l'Est, élite anglophone formée très tôt (ex. E. M. L. Endeley, J. N. Foncha).
\end{itemize}
\textbf{Point commun :} dans les deux zones, l'économie est extravertie (tournée vers l'exportation), la main-d'œuvre est contrainte (prestations ou migration forcée par la nécessité monétaire), et l'État social est quasi inexistant (laissé aux missions).
\end{exemplebox}

\courssub{Méthode 8 : Bilan et transition (1939-1946)}

\begin{methode}[Du mandat à la tutelle : ruptures et continuités]
\begin{enumerate}
\item \textbf{1939-1940 :} Guerre, ralliement du Cameroun français à la France Libre (août 1940, De Gaulle à Douala), engagement des tirailleurs camerounais.
\item \textbf{1944 :} Conférence de Brazzaville (réforme coloniale, fin de l'indigénat \emph{théorique}, citoyenneté pour quelques-uns).
\item \textbf{1945-1946 :} Création de l'ONU (juin 1945). Accords de tutelle signés (décembre 1946) : les mandats deviennent \cle{territoires sous tutelle}. Objectif affiché : autonomie progressive. Maintien du partage territorial (France / Grande-Bretagne).
\item \textbf{Continuité :} structures administratives, économie de plantation, code de l'indigénat (jusque 1946 côté français) persistent.
\end{enumerate}
\end{methode}

\begin{exoresolu}[La transition 1939-1946 : rupture ou continuité ?]
\textbf{Énoncé.} « Le passage du mandat à la tutelle en 1946 change le vocabulaire mais pas la réalité coloniale. » Discute cette affirmation en t'appuyant sur deux arguments pour et deux contre.
\tcblower
\textbf{Solution détaillée.}
\textbf{Arguments POUR (continuité) :}
\begin{enumerate}
\item \textbf{Statut international :} La tutelle (ONU) remplace le mandat (SDN) mais le principe reste : administration déléguée à une puissance étrangère (France/UK) avec reddition de comptes. Le partage territorial de 1922 est maintenu.
\item \textbf{Structures sur place :} Administration directe / indirect rule inchangées. Code de l'indigénat aboli formellement en 1946 (Loi Lamine Guèye) mais pratiques de travail forcé et répression persistent. Économie d'exportation inchangée.
\end{enumerate}
\textbf{Arguments CONTRE (rupture) :}
\begin{enumerate}
\item \textbf{Droit nouveau :} La Charte de l'ONU (Art. 76) fixe explicitement l'objectif d'\emph{autonomie} ou d'\emph{indépendance} (pas seulement « bien-être et développement » comme la SDN). Création du Conseil de tutelle, visites de pétitionnaires camerounais à New York possibles.
\item \textbf{Vie politique :} 1944-1946 : naissance des premiers partis politiques (USC, UPC côté français ; KNC, KPP côté britannique), droit syndical (décret 1944), premières élections (collèges électoraux) en 1946. Les Camerounais deviennent des acteurs politiques.
\end{enumerate}
\textbf{Synthèse.} 1946 est une \textbf{rupture juridique et politique majeure} (droit à l'indépendance, entrée en politique) mais une \textbf{continuité administrative et économique} sur le terrain immédiat. C'est le point de départ de la décolonisation, non son aboutissement.
\end{exoresolu}

\begin{attention}[Les 5 erreurs qui coûtent des points au Bac]
\begin{enumerate}
\item \textbf{Confondre mandat et colonie.} Le Cameroun n'est pas annexé : il est placé sous \cle{mandat} de la SDN (1922), puis sous \cle{tutelle} de l'ONU (1946). Écrire « colonie française » pour la période 1922-1946 est une erreur de statut juridique.
\item \textbf{Inverser les zones.} La France administre la \textbf{grande partie orientale} (environ 80 \%, Yaoundé, Douala) ; la Grande-Bretagne les \textbf{deux bandes occidentales} le long du Nigeria (Buea, Bamenda, Garoua). Ne jamais attribuer Douala aux Britanniques.
\item \textbf{Confondre administration directe et indirect rule.} Directe = fonctionnaires français aux commandes (cercles, subdivisions) ; indirect rule = gouvernement par les chefs traditionnels (Native Authorities, système britannique, hérité de Lugard au Nigeria). Vérifie toujours quel territoire applique quelle méthode.
\item \textbf{Oublier de dater les événements.} Une réponse qui dit « après la guerre, le Cameroun est partagé » sans citer 1916 (Mora), 1919 (Versailles) et 1922 (mandats SDN) perd l'essentiel des points de justification.
\item \textbf{Paraphraser le document au lieu de le commenter.} Recopier le texte sans apport personnel ne rapporte rien : chaque élément cité doit être expliqué avec tes connaissances (SDN, obligation de rendre compte, réalité de l'exploitation : impôt, prestations, travail forcé).
\end{enumerate}
\end{attention}

\newpage

\coursec{Exercices}

\courssub{Je vérifie mes connaissances}

\begin{exercice}[QCM : les notions clés de la période de mandat]
Pour chaque question, choisis la bonne réponse et justifie ton choix par une phrase.
\begin{enumerate}
\item La Société des Nations (SDN), qui confie les mandats sur le Cameroun, est créée en :
\begin{enumerate}
\item 1914 \quad b) 1919 \quad c) 1939 \quad d) 1946
\end{enumerate}
\item Le mandat sur le Cameroun français est confirmé par la SDN en :
\begin{enumerate}
\item 1916 \quad b) 1919 \quad c) 1922 \quad d) 1945
\end{enumerate}
\item L'\emph{indigénat} désigne :
\begin{enumerate}
\item un régime juridique d'exception appliqué aux populations colonisées ;
\item une taxe payée par les chefs traditionnels ;
\item une école réservée aux fils de notables ;
\item un territoire autonome.
\end{enumerate}
\item La \emph{Native Authority} est une institution du système :
\begin{enumerate}
\item français d'administration directe ;
\item britannique d'\emph{indirect rule} ;
\item allemand du \emph{Schutzgebiet} ;
\item de la tutelle de l'ONU.
\end{enumerate}
\end{enumerate}
\end{exercice}

\begin{exercice}[Vrai ou faux, avec justification]
Indique si chaque affirmation est vraie ou fausse, puis justifie ta réponse en une ou deux phrases.
\begin{enumerate}
\item La partition du Cameroun entre la France et la Grande-Bretagne est une conséquence de la défaite allemande en 1918.
\item La France et la Grande-Bretagne deviennent propriétaires des territoires camerounais qu'elles administrent sous mandat.
\item La capitation est un impôt par tête d'habitant exigé dans le Cameroun sous administration française.
\item Le travail forcé est totalement absent du Cameroun britannique entre 1922 et 1939.
\item La période de mandat s'achève en 1946 avec la mise en place du régime de tutelle de l'ONU.
\end{enumerate}
\end{exercice}

\begin{exercice}[Définitions]
Définis en deux ou trois phrases chacun des termes suivants, en donnant pour chacun un exemple camerounais précis :
\begin{enumerate}
\item le \cle{mandat} (de type B) ;
\item la \cle{partition} du Cameroun ;
\item l'\cle{administration directe} ;
\item l'\cle{indirect rule} ;
\item la \cle{tutelle}.
\end{enumerate}
\end{exercice}

\begin{exercice}[Repères chronologiques et calculs simples]
\begin{enumerate}
\item Replace dans l'ordre chronologique les événements suivants : confirmation des mandats par la SDN ; début de la Première Guerre mondiale ; conquête franco-britannique du Cameroun ; conférence de San Remo ; passage du mandat à la tutelle.
\item La France administre environ les quatre cinquièmes du territoire de l'ancien Kamerun, évalué à \num{500000} km$^2$. Calcule la superficie approximative administrée par la France, puis celle administrée par la Grande-Bretagne.
\item Entre 1916 (conquête) et 1946 (tutelle), combien d'années le Cameroun a-t-il vécu sous administration franco-britannique ? Exprime cette durée en pourcentage du XX\textsuperscript{e} siècle (1901-2000).
\end{enumerate}
\end{exercice}

\courssub{Je m'entraîne}

\begin{exercice}[Question à réponse construite : deux modes d'administration]
En une vingtaine de lignes, explique les différences entre l'administration directe française et l'\emph{indirect rule} britannique au Cameroun. Ta réponse doit mobiliser au moins deux exemples précis pour chaque puissance mandataire (institutions, rôle des chefs traditionnels, langue, justice).
\end{exercice}

\begin{exercice}[Construction cartographique]
Sur une carte muette du Cameroun fournie par l'enseignant :
\begin{enumerate}
\item trace la frontière séparant la zone française de la zone britannique (deux bandes le long de la frontière nigériane) ;
\item place et nomme les capitales Yaoundé et Buéa, ainsi que la ville de Douala ;
\item figure par un trait le chemin de fer Transcamerounais (Douala--Yaoundé, puis prolongement vers Ngaoundéré) ;
\item indique par un symbole le port de Douala et le port de Victoria (Limbé) ;
\item rédige une légende complète et donne un titre à ta carte.
\end{enumerate}
\end{exercice}

\begin{exercice}[Analyse de tableau statistique]
Le tableau ci-dessous donne des ordres de grandeur de la production d'exportation du Cameroun français.
\begin{center}
\begin{tabularx}{\linewidth}{|l|X|X|}\hline
\rowcolor{popblueL} Produit & Années 1920 (tonnes) & Années 1930 (tonnes) \\ \hline
Cacao & \num{15000} & \num{40000} \\ \hline
Café & \num{1000} & \num{8000} \\ \hline
Bananier & faible & \num{20000} \\ \hline
\end{tabularx}
\end{center}
\begin{enumerate}
\item Calcule le taux de variation de la production de cacao entre les deux périodes.
\item À quelles politiques économiques de la puissance mandataire (mise en valeur, travail forcé, grands travaux) attribues-tu cette progression ?
\item Quelles conséquences sociales cette économie de traite entraîne-t-elle pour les populations camerounaises ? Cite deux exemples.
\end{enumerate}
\end{exercice}

\begin{exercice}[Frise chronologique]
Construis une frise chronologique de 1914 à 1946 comportant au moins huit dates repères de la période de mandat au Cameroun (guerre, conquête, partage, confirmation des mandats, grandes réalisations, crise de 1929, Seconde Guerre mondiale, tutelle). Pour chaque date, rédige une légende d'une ligne.
\end{exercice}

\courssub{Je me prépare au Bac}

\begin{exercice}[Dissertation type Probatoire]
\textbf{Sujet :} « La partition du Cameroun de 1919 : une simple conséquence de la guerre ou un bouleversement durable pour les populations ? »

\begin{enumerate}
\item Définis les termes du sujet et délimite le cadre spatio-temporel.
\item Rédige une introduction comportant une accroche, une présentation du sujet, une problématique et l'annonce du plan.
\item Développe en deux grandes parties : (a) la partition comme résultat direct de la Première Guerre mondiale et des décisions des vainqueurs ; (b) la partition comme bouleversement durable (double administration, morcellement des peuples frontaliers, héritages linguistiques et institutionnels).
\item Rédige une conclusion qui répond à la problématique et ouvre sur la période de tutelle.
\end{enumerate}
\end{exercice}

\begin{exercice}[Étude critique de documents]
\textbf{Document 1 (extrait adapté d'un rapport annuel de la France à la SDN, années 1920) :} « La puissance mandataire assure dans le territoire la paix, l'ordre et la bonne administration ; elle veille au progrès matériel et moral des populations, interdit la traite des esclaves et le travail forcé sauf pour les travaux publics essentiels, et rend compte chaque année à la Commission des mandats. »

\textbf{Document 2 (témoignage adapté d'un administré du Centre-Cameroun, années 1920-1930) :} « Les hommes valides étaient requis pour porter les charges et construire la route et le chemin de fer. Beaucoup ne revenaient pas des chantiers. L'impôt de capitation obligeait chacun à trouver de l'argent, même quand la récolte était mauvaise. »

\begin{enumerate}
\item Présente chaque document (nature, auteur, date, contexte).
\item Dégage du document 1 les obligations de la puissance mandataire envers la SDN.
\item Montre, à partir du document 2, l'écart entre ces obligations et les réalités vécues par les administrés.
\item Conclus en discutant la portée et les limites du contrôle exercé par la SDN sur le Cameroun français.
\end{enumerate}
\end{exercice}

\begin{exercice}[Sujet type Baccalauréat technique]
\textbf{Sujet :} « La vie des Camerounais sous administration française et anglaise (1922-1946) : exploitation ou modernisation ? »

\begin{enumerate}
\item Rappelle les conditions dans lesquelles la France et la Grande-Bretagne obtiennent leurs mandats sur le Cameroun.
\item Montre, avec au moins deux arguments documentés, les aspects d'exploitation des populations (impôts, travail forcé et prestations, réquisitions, indigénat, plantations).
\item Montre, avec au moins deux arguments documentés, les réalisations présentées comme une modernisation (Transcamerounais, port de Douala, écoles, dispensaires, cultures d'exportation).
\item Compare les situations vécues dans les deux zones d'administration (administration directe contre \emph{indirect rule}).
\item Rédige une conclusion argumentée prenant position sur l'alternative « exploitation ou modernisation ».
\end{enumerate}
\textbf{Barème indicatif :} introduction 2 pts ; développement 12 pts ; conclusion 2 pts ; cohérence et expression 4 pts.
\end{exercice}

\newpage

\coursec{Corrigés des exercices}

\begin{corrige}[Exercice 1 — QCM : les notions clés de la période de mandat]
\begin{enumerate}
\item \textbf{Réponse b) 1919.} La Société des Nations est créée par le traité de Versailles (1919), au lendemain de la Première Guerre mondiale, pour garantir la paix et organiser notamment le système des mandats.
\item \textbf{Réponse c) 1922.} Après la conférence de San Remo (1920), la SDN confirme officiellement les mandats français et britannique sur le Cameroun le 20 juillet 1922.
\item \textbf{Réponse a).} L'indigénat est un régime juridique d'exception : il soumet les populations colonisées à des sanctions administratives (amendes, emprisonnement) sans procès régulier, contrairement aux citoyens français.
\item \textbf{Réponse b).} La \emph{Native Authority} est l'institution centrale de l'\emph{indirect rule} britannique : le pouvoir est exercé par les chefs traditionnels sous le contrôle du résident britannique.
\end{enumerate}
\end{corrige}

\begin{corrige}[Exercice 2 — Vrai ou faux, avec justification]
\begin{enumerate}
\item \textbf{Vrai.} La défaite de l'Allemagne en 1918 entraîne la perte de ses colonies ; le Kamerun, conquis en 1916 par les Alliés, est partagé entre la France et la Grande-Bretagne, partage confirmé par les mandats de la SDN.
\item \textbf{Faux.} Le mandat n'est pas une annexion : les puissances mandataires administrent les territoires \emph{au nom de la SDN} et doivent lui rendre compte chaque année ; elles n'en sont pas propriétaires.
\item \textbf{Vrai.} La capitation est un impôt par tête d'habitant, exigé en argent dans le Cameroun français ; il pousse les paysans à vendre leur travail ou leurs récoltes pour se procurer des liquidités.
\item \textbf{Faux.} Le travail forcé existe aussi dans la zone britannique (portage, prestations, construction de routes), même s'il y est moins systématique que dans la zone française.
\item \textbf{Vrai.} En 1946, le régime de tutelle de l'Organisation des Nations unies (ONU) succède au système des mandats de la SDN, dissoute après la Seconde Guerre mondiale.
\end{enumerate}
\end{corrige}

\begin{corrige}[Exercice 3 — Définitions]
\begin{enumerate}
\item \textbf{Le mandat (de type B)} : système institué par la SDN en 1919 par lequel une puissance victorieuse administre une ancienne colonie allemande au nom de la communauté internationale, avec obligation de rendre compte chaque année à la Commission des mandats. \emph{Exemple :} la France administre le Cameroun oriental sous mandat de type B à partir de 1922.
\item \textbf{La partition du Cameroun} : partage de l'ancien Kamerun allemand entre la France et la Grande-Bretagne après la conquête de 1916, officialisé en 1919-1922. \emph{Exemple :} la France reçoit environ les quatre cinquièmes du territoire (capitale Yaoundé), la Grande-Bretagne deux bandes occidentales le long de la frontière nigériane (capitale Buéa).
\item \textbf{L'administration directe} : mode de gouvernement colonial français dans lequel les administrateurs français commandent directement les populations, les chefs traditionnels étant réduits au rôle d'auxiliaires nommés et révocables. \emph{Exemple :} au Cameroun français, les chefs de canton sont désignés par l'administration et appliquent ses ordres (impôt, recrutement de main-d'œuvre).
\item \textbf{L'indirect rule} : mode d'administration britannique consistant à gouverner par l'intermédiaire des chefs traditionnels reconnus, regroupés dans les \emph{Native Authorities}, sous le contrôle d'un résident britannique. \emph{Exemple :} dans le Cameroun britannique, les chefs collectent l'impôt et jugent les litiges coutumiers dans les \emph{Native Courts}.
\item \textbf{La tutelle} : régime institué par l'ONU en 1946 pour succéder aux mandats ; il vise officiellement à conduire les territoires vers l'autonomie et l'indépendance. \emph{Exemple :} en 1946, le Cameroun français et le Cameroun britannique passent sous tutelle des Nations unies.
\end{enumerate}
\end{corrige}

\begin{corrige}[Exercice 4 — Repères chronologiques et calculs simples]
\begin{enumerate}
\item Ordre chronologique :
\begin{enumerate}
\item début de la Première Guerre mondiale (1914) ;
\item conquête franco-britannique du Cameroun (1914-1916) ;
\item conférence de San Remo (1920) ;
\item confirmation des mandats par la SDN (1922) ;
\item passage du mandat à la tutelle (1946).
\end{enumerate}
\item La superficie totale est de $\num{500000}$ km$^2$ et la France administre les quatre cinquièmes :
\[ S_{\text{France}} = \frac{4}{5} \times \num{500000} = \num{400000} \text{ km}^2. \]
La Grande-Bretagne administre le reste :
\[ S_{\text{GB}} = \num{500000} - \num{400000} = \num{100000} \text{ km}^2. \]
\emph{Vérification :} $\num{400000} + \num{100000} = \num{500000}$ km$^2$ : le calcul est cohérent.
\item Durée sous administration franco-britannique :
\[ 1946 - 1916 = 30 \text{ ans}. \]
Le XX\textsuperscript{e} siècle (1901-2000) compte 100 ans, donc :
\[ \frac{30}{100} \times 100 = 30\,\%. \]
Le Cameroun a donc vécu 30 ans sous administration franco-britannique, soit \textbf{30 \%} du XX\textsuperscript{e} siècle.
\end{enumerate}
\end{corrige}

\begin{corrige}[Exercice 5 — Question à réponse construite : deux modes d'administration]
\textbf{Éléments de réponse attendus (environ vingt lignes).}

Après la partition du Cameroun, la France et la Grande-Bretagne appliquent deux modes d'administration très différents.

Dans la zone française, c'est l'\textbf{administration directe} : les administrateurs français gouvernent sans intermédiaire. Les chefs traditionnels ne sont que des auxiliaires nommés et révocables par l'administration (premier exemple : les chefs de canton, chargés de collecter la capitation et de fournir la main-d'œuvre). La justice est rendue selon le droit français pour les citoyens, tandis que les administrés relèvent de l'\textbf{indigénat}, régime d'exception permettant de punir sans procès (deuxième exemple). La langue d'administration et d'enseignement est le français, et l'école vise à diffuser la culture française (politique d'assimilation).

Dans la zone britannique, c'est l'\textbf{\emph{indirect rule}} : la Grande-Bretagne gouverne par l'intermédiaire des chefs traditionnels reconnus, regroupés dans les \emph{Native Authorities} (premier exemple), qui collectent l'impôt et administrent sous le contrôle d'un résident britannique. Les litiges coutumiers sont jugés par les \emph{Native Courts} selon la coutume locale (deuxième exemple). La langue d'administration est l'anglais, et l'enseignement est largement laissé aux missions. Le Cameroun britannique est d'ailleurs administré comme une annexe du Nigeria voisin.

En somme, la France commande directement et cherche à assimiler, tandis que la Grande-Bretagne s'appuie sur les structures traditionnelles pour gouverner à moindre coût. Dans les deux cas, cependant, les populations restent soumises à l'impôt et aux réquisitions.

\textbf{Points de vigilance :} exiger deux exemples par puissance (institutions, rôle des chefs, langue, justice) ; éviter l'opposition caricaturale « bons Britanniques / mauvais Français » : les deux systèmes reposent sur la domination coloniale.
\end{corrige}

\begin{corrige}[Exercice 6 — Construction cartographique]
\textbf{Corrigé indicatif (à reporter sur la carte muette).}

\begin{enumerate}
\item \textbf{Frontière :} tracer une ligne nord-sud le long de la frontière avec le Nigeria, matérialisant les deux bandes britanniques à l'ouest (du lac Tchad au golfe de Guinée) ; la grande partie orientale est la zone française.
\item \textbf{Villes :} Yaoundé (capitale de la zone française, au centre-sud) ; Buéa (capitale de la zone britannique, au pied du mont Cameroun) ; Douala (grande ville portuaire, sur l'estuaire du Wouri).
\item \textbf{Chemin de fer :} tracer un trait Douala--Yaoundé (ligne du Centre), puis un prolongement vers le nord jusqu'à Ngaoundéré (ligne du Nord, achevée après la période de mandat mais prolongeant la logique des grands travaux).
\item \textbf{Ports :} symbole d'ancre à Douala (port principal) et à Victoria, actuelle Limbé (port britannique au sud-ouest).
\item \textbf{Légende :} couleur 1 = zone sous administration française ; couleur 2 = zone sous administration britannique ; carré = capitale ; point = ville ; trait noir = chemin de fer ; ancre = port. \textbf{Titre proposé :} « Le Cameroun sous mandats français et britannique (1922-1946) ».
\end{enumerate}

\textbf{Points de vigilance :} ne pas oublier le titre, la légende et l'orientation (flèche du nord) ; la frontière doit suivre la limite avec le Nigeria, non couper le pays au milieu.
\end{corrige}

\begin{corrige}[Exercice 7 — Analyse de tableau statistique]
\begin{enumerate}
\item \textbf{Taux de variation du cacao} entre les années 1920 (\num{15000} t) et les années 1930 (\num{40000} t) :
\[ t = \frac{\num{40000} - \num{15000}}{\num{15000}} \times 100 = \frac{\num{25000}}{\num{15000}} \times 100 \approx 166{,}7\,\%. \]
La production de cacao a donc augmenté d'environ \textbf{167 \%} (elle a été multipliée par environ $2{,}7$).
\item \textbf{Politiques explicatives :} la politique de \emph{mise en valeur} du territoire (développement des cultures d'exportation : cacao, café, bananes) ; le recours au \emph{travail forcé} et aux prestations pour les plantations et les chantiers ; les \emph{grands travaux} (chemin de fer Transcamerounais, routes, port de Douala) qui permettent l'évacuation des produits vers la métropole.
\item \textbf{Conséquences sociales (deux exemples) :}
\begin{itemize}
\item les déplacements forcés de main-d'œuvre vers les chantiers et les plantations, qui dépeuplent les villages et provoquent de nombreux décès (chantier du Transcamerounais) ;
\item la pression fiscale (capitation) qui oblige les paysans à abandonner les cultures vivrières au profit des cultures de rente, entraînant disettes et paupérisation.
\end{itemize}
\end{enumerate}
\end{corrige}

\begin{corrige}[Exercice 8 — Frise chronologique]
\textbf{Exemple de frise attendue (huit dates minimum) :}
\begin{itemize}
\item \textbf{1914 :} début de la Première Guerre mondiale ; les troupes alliées pénètrent au Kamerun.
\item \textbf{1916 :} capitulation allemande ; conquête franco-britannique du Cameroun et partage militaire provisoire.
\item \textbf{1919 :} traité de Versailles ; création de la SDN ; la partition est entérinée.
\item \textbf{1920 :} conférence de San Remo, attribution des mandats à la France et à la Grande-Bretagne.
\item \textbf{1922 :} confirmation officielle des mandats par la SDN (20 juillet).
\item \textbf{1921-1927 :} construction du chemin de fer Douala--Yaoundé (Transcamerounais, ligne du Centre) au prix d'une lourde mortalité.
\item \textbf{1929 :} la crise économique mondiale frappe les cultures d'exportation camerounaises (effondrement des prix du cacao).
\item \textbf{1939-1945 :} Seconde Guerre mondiale ; le Cameroun français se rallie à la France libre (1940).
\item \textbf{1946 :} le régime de tutelle de l'ONU succède au mandat.
\end{itemize}
\textbf{Points de vigilance :} respecter l'ordre chronologique, espacer les dates régulièrement sur la frise, rédiger une légende courte (une ligne) par date.
\end{corrige}

\begin{corrige}[Exercice 9 — Dissertation type Probatoire]
\textbf{1. Définition des termes et cadre.} La \emph{partition} désigne le partage de l'ancien Kamerun allemand entre la France et la Grande-Bretagne ; un \emph{bouleversement durable} renvoie à des transformations profondes et durables pour les populations. Cadre : le Cameroun, de 1919 (entérinement du partage) à 1946 (fin du mandat).

\textbf{2. Introduction (modèle).} \emph{Accroche :} en 1916, les armées franco-britanniques chassent l'Allemagne du Kamerun. \emph{Présentation :} les vainqueurs se partagent le territoire, partage confirmé par les mandats de la SDN en 1922. \emph{Problématique :} la partition de 1919 est-elle une simple conséquence de la guerre, ou marque-t-elle un bouleversement durable pour les populations camerounaises ? \emph{Annonce du plan :} nous verrons d'abord que la partition résulte directement de la guerre et des décisions des vainqueurs, puis qu'elle constitue un bouleversement durable.

\textbf{3. Développement.}
\emph{Première partie — une conséquence directe de la guerre :} la défaite allemande de 1918 prive l'Allemagne de ses colonies ; la conquête de 1914-1916 a déjà créé un partage militaire de fait ; les traités de 1919 et la SDN (mandats confirmés en 1922) ne font qu'officialiser le rapport de forces issu de la guerre.

\emph{Deuxième partie — un bouleversement durable :} double administration (directe française, \emph{indirect rule} britannique) aux institutions, langues et monnaies différentes ; morcellement de peuples frontaliers (peuples du Grassfield, populations du nord séparées du Nigeria ou du Tchad) ; héritages linguistiques durables (Cameroun francophone et anglophone) ; nouvelles contraintes économiques (impôts, travail forcé, cultures d'exportation).

\textbf{4. Conclusion (modèle).} La partition est bien née de la guerre, mais elle dépasse largement ce cadre : elle recompose durablement la vie politique, linguistique et sociale du Cameroun. \emph{Ouverture :} en 1946, le passage à la tutelle de l'ONU prolonge cette double tutelle jusqu'aux indépendances.

\textbf{Barème indicatif :} introduction 3 pts ; développement (2 parties équilibrées, exemples précis) 12 pts ; conclusion 2 pts ; cohérence et expression 3 pts.

\textbf{Points de vigilance :} définir les termes dès l'introduction ; ne pas réduire la copie à un récit de la guerre ; mobiliser des exemples précis (San Remo 1920, mandats 1922, indigénat, \emph{Native Authority}).
\end{corrige}

\begin{corrige}[Exercice 10 — Étude critique de documents]
\begin{enumerate}
\item \textbf{Présentation des documents.} Document 1 : texte administratif, extrait d'un rapport annuel adressé par la France à la SDN dans les années 1920, dans le cadre de l'obligation de reddition de comptes du système des mandats. Document 2 : témoignage oral d'un administré du Centre-Cameroun, datant des années 1920-1930, période des grands chantiers (routes, Transcamerounais) et de la mise en place de la fiscalité coloniale.
\item \textbf{Obligations de la puissance mandataire (document 1) :} assurer la paix, l'ordre et la bonne administration ; veiller au progrès matériel et moral des populations ; interdire la traite des esclaves ; interdire le travail forcé sauf pour les travaux publics essentiels ; rendre compte chaque année à la Commission des mandats.
\item \textbf{Écart avec les réalités (document 2) :} le travail forcé, officiellement limité, devient massif : les hommes valides sont requis pour le portage, la route et le chemin de fer ; la mortalité sur les chantiers est élevée (« beaucoup ne revenaient pas ») ; la capitation impose une charge fiscale écrasante, même en cas de mauvaise récolte. Le « progrès matériel et moral » proclamé contraste donc avec l'exploitation vécue.
\item \textbf{Portée et limites du contrôle de la SDN :} la SDN impose une obligation de rapports annuels et une Commission des mandats, ce qui constitue une première forme de contrôle international. Mais ce contrôle reste limité : la SDN n'a pas de moyen de coercition, elle s'appuie sur les rapports des puissances elles-mêmes, qui enjolivent la réalité ; les plaintes des populations sont rarement entendues. Le mandat apparaît ainsi comme une tutelle de principe plus que de fait.
\end{enumerate}
\textbf{Points de vigilance :} toujours présenter les documents avant de les exploiter ; citer le texte pour appuyer chaque argument ; ne pas confondre SDN (1919) et ONU (1945).
\end{corrige}

\begin{corrige}[Exercice 11 — Sujet type Baccalauréat technique]
\textbf{1. Obtention des mandats.} Le Kamerun allemand est conquis par les troupes franco-britanniques en 1914-1916 pendant la Première Guerre mondiale. La défaite allemande (1918) et le traité de Versailles (1919) entérinent le partage ; la conférence de San Remo (1920) attribue les mandats, confirmés par la SDN en 1922 : environ les quatre cinquièmes du territoire à la France, deux bandes occidentales à la Grande-Bretagne.

\textbf{2. Aspects d'exploitation (au moins deux arguments).}
\begin{itemize}
\item la \emph{capitation}, impôt par tête payé en argent, qui contraint les paysans à vendre leur travail ;
\item le \emph{travail forcé} et les prestations : portage, construction des routes et du Transcamerounais, avec une forte mortalité sur les chantiers ;
\item l'\emph{indigénat} en zone française, régime de sanctions sans procès ; les réquisitions et les plantations européennes accaparant les meilleures terres.
\end{itemize}

\textbf{3. Réalisations présentées comme modernisation (au moins deux arguments).}
\begin{itemize}
\item le chemin de fer Transcamerounais (Douala--Yaoundé, 1921-1927) et l'équipement du port de Douala ;
\item l'ouverture d'écoles (et d'écoles de mission en zone britannique) et de dispensaires ; le développement des cultures d'exportation (cacao, café, bananes) qui intègrent le Cameroun aux échanges mondiaux.
\end{itemize}

\textbf{4. Comparaison des deux zones.} Zone française : administration directe, chefs réduits à des auxiliaires nommés, indigénat, francisation, forte pression fiscale et de main-d'œuvre. Zone britannique : \emph{indirect rule}, gouvernement par les \emph{Native Authorities} et les \emph{Native Courts}, enseignement largement missionnaire, territoire administré depuis le Nigeria ; l'exploitation y est réelle (impôts, portage) mais moins intensive.

\textbf{5. Conclusion argumentée (modèle).} La période de mandat combine indissociablement exploitation et modernisation : les infrastructures et les écoles sont réelles, mais elles servent d'abord l'économie de traite et reposent sur le travail forcé et la fiscalité. Il s'agit donc d'une modernisation au service de l'exploitation, dont les effets durables préparent, sous tutelle, les revendications nationalistes.

\textbf{Barème indicatif :} introduction 2 pts ; développement 12 pts ; conclusion 2 pts ; cohérence et expression 4 pts.

\textbf{Points de vigilance :} répondre à l'alternative en prenant position argumentée ; documenter chaque argument (dates, institutions, exemples) ; comparer explicitement les deux zones au lieu de juxtaposer deux descriptions.
\end{corrige}

\newpage

\coursec{Fiche bilan}

\begin{popfigure}
\begin{tikzpicture}[mindmap, grow cyclic, every node/.style={concept, text=white, font=\small\bfseries},
  root concept/.append style={concept color=popdark, font=\normalsize\bfseries, minimum size=3.2cm},
  level 1/.append style={level distance=3.6cm, sibling angle=60},
  level 2/.append style={level distance=2.2cm, font=\scriptsize\bfseries, minimum size=1.6cm}]
\node[root concept]{La période de mandat (1922-1946)}
  child[concept color=popblue]{ node {Partition du Cameroun}
    child { node[concept color=popblueL, text=popdark] {1914-1916 : conquête} }
    child { node[concept color=popblueL, text=popdark] {1919 : traité de Versailles} }
    child { node[concept color=popblueL, text=popdark] {1922 : mandats SDN} }
  }
  child[concept color=poporange]{ node {Partage territorial}
    child { node[concept color=poporangeL, text=popdark] {4/5 à la France} }
    child { node[concept color=poporangeL, text=popdark] {1/5 à l'Angleterre} }
  }
  child[concept color=popgreen]{ node {Modes d'administration}
    child { node[concept color=popgreenL, text=popdark] {Directe (France)} }
    child { node[concept color=popgreenL, text=popdark] {Indirect rule (GB)} }
  }
  child[concept color=poppurple]{ node {Vie sous mandat français}
    child { node[concept color=poppurpleL, text=popdark] {Travail forcé} }
    child { node[concept color=poppurpleL, text=popdark] {Chefferies, écoles} }
  }
  child[concept color=poppink]{ node {Cameroun britannique}
    child { node[concept color=poppinkL, text=popdark] {Rattaché au Nigeria} }
    child { node[concept color=poppinkL, text=popdark] {Plantations CDC} }
  }
  child[concept color=popgold]{ node {Du mandat à la tutelle}
    child { node[concept color=popgoldL, text=popdark] {1939 : SDN paralysée} }
    child { node[concept color=popgoldL, text=popdark] {1946 : tutelle ONU} }
  };
\end{tikzpicture}
\legende{Carte mentale de la leçon : les six idées essentielles de la période de mandat au Cameroun.}
\end{popfigure}

\begin{bilan}
\textbf{Définitions clés à connaître par c\oe ur}
\begin{itemize}
  \item \cle{Mandat} : système créé par la Société des Nations (SDN) en 1919 confiant l'administration d'un territoire à une puissance victorieuse, sous le contrôle de la SDN.
  \item \cle{Mandat de type B} : catégorie de mandat (dont relève le Cameroun) où la puissance mandataire administre le territoire mais ne peut ni l'annexer ni y construire de bases militaires.
  \item \cle{Administration directe} : mode de gouvernement français où les administrateurs français commandent directement, les chefs locaux n'étant que des auxiliaires nommés et révocables.
  \item \cle{Indirect rule} : mode de gouvernement britannique qui administre par l'intermédiaire des chefs traditionnels reconnus (\emph{Native Authorities}).
  \item \cle{Tutelle} : régime international institué par l'ONU en 1946, succédant au mandat, avec pour objectif déclaré de conduire les peuples vers l'autonomie.
\end{itemize}

\textbf{Repères chronologiques essentiels}
\begin{center}
\begin{tabularx}{\linewidth}{|l|X|}
\hline
\rowcolor{popblueL} \textbf{Date} & \textbf{Événement} \\
\hline
1914-1916 & Campagne militaire franco-britannique ; capitulation allemande à Mora (février 1916). \\
\hline
1919 & Traité de Versailles : l'Allemagne renonce à ses colonies ; principe du mandat adopté. \\
\hline
20 juillet 1922 & Confirmation des mandats par la SDN : Cameroun français (environ 4/5) et Cameroun britannique (environ 1/5). \\
\hline
1921-1927 & Construction du chemin de fer Douala-Yaoundé (travail forcé, lourdes pertes humaines). \\
\hline
1939-1945 & Seconde Guerre mondiale : la SDN est paralysée ; le Cameroun français se rallie à la France libre (1940). \\
\hline
13 décembre 1946 & Accords de tutelle de l'ONU : fin du régime de mandat. \\
\hline
\end{tabularx}
\end{center}

\textbf{Comparaison à connaître : les deux administrations}
\begin{center}
\begin{tabularx}{\linewidth}{|X|X|X|}
\hline
\rowcolor{popblueL} \textbf{Critère} & \textbf{Cameroun français} & \textbf{Cameroun britannique} \\
\hline
Superficie & Environ 4/5 du territoire & Environ 1/5 (deux bandes : Northern et Southern Cameroons) \\
\hline
Capitale & Yaoundé & Buea (administré depuis le Nigeria) \\
\hline
Mode d'administration & Administration directe & Indirect rule \\
\hline
Économie & Mise en valeur, grands travaux (rail, routes, ports), travail forcé & Plantations (CDC : bananes, palmiers), économie liée au Nigeria \\
\hline
\end{tabularx}
\end{center}

\textbf{Méthode express : réussir un commentaire de document sur le mandat}
\begin{enumerate}
  \item Identifier le document (nature, auteur, date, contexte : SDN, puissance mandataire).
  \item Dégager l'idée principale en une phrase.
  \item Analyser en deux ou trois parties (faits, acteurs, conséquences).
  \item Conclure en reliant le document au bilan de la période de mandat.
\end{enumerate}
\end{bilan}

\begin{aretenir}[Checklist avant le Bac]
\begin{itemize}
  \item $\square$ Je sais expliquer pourquoi la Première Guerre mondiale a entraîné la partition du Cameroun.
  \item $\square$ Je connais la date de capitulation allemande au Cameroun (février 1916, Mora).
  \item $\square$ Je sais définir le système de mandat de la SDN et le mandat de type B.
  \item $\square$ Je retiens la date du 20 juillet 1922 (confirmation des mandats).
  \item $\square$ Je connais la répartition territoriale : environ 4/5 à la France, 1/5 à la Grande-Bretagne.
  \item $\square$ Je sais distinguer administration directe (France) et indirect rule (Grande-Bretagne).
  \item $\square$ Je peux citer les réalisations économiques du mandat français (chemin de fer Douala-Yaoundé, routes, port de Douala).
  \item $\square$ Je sais que le travail forcé et les impôts ont pesé lourdement sur les populations.
  \item $\square$ Je sais que le Cameroun britannique était administré comme une dépendance du Nigeria.
  \item $\square$ Je connais le rôle des plantations (CDC) dans le Cameroun britannique.
  \item $\square$ Je sais expliquer le passage du mandat à la tutelle en 1946 (création de l'ONU).
  \item $\square$ Je sais rédiger une introduction, un développement en parties et une conclusion pour un commentaire de document.
\end{itemize}
\end{aretenir}

\findecours

\end{document}
